\documentclass[11pt]{amsart}

\usepackage{amsmath,amssymb,amsthm}
\usepackage[a4paper,margin=25mm,includeheadfoot]{geometry}
\usepackage[T1]{fontenc}
\usepackage{lmodern}
\usepackage[hidelinks]{hyperref}

\newenvironment{romanlist}{\begin{enumerate}\renewcommand{\labelenumi}{(\roman{enumi})}%
  \renewcommand{\theenumi}{\roman{enumi}}}{\end{enumerate}}
\newenvironment{alphalist}{\begin{enumerate}\renewcommand{\labelenumi}{(\alph{enumi})}%
  \renewcommand{\theenumi}{\alph{enumi}}}{\end{enumerate}}

\newtheorem{theorem}{Theorem}[section]
\newtheorem{proposition}[theorem]{Proposition}
\newtheorem{lemma}[theorem]{Lemma}
\newtheorem{corollary}[theorem]{Corollary}
\newtheorem{fact}{Fact}
\renewcommand{\thefact}{AK\arabic{fact}}
\theoremstyle{definition}
\newtheorem{definition}[theorem]{Definition}
\theoremstyle{remark}
\newtheorem{remark}[theorem]{Remark}

\newcommand{\C}{\mathbb{C}}
\newcommand{\R}{\mathbb{R}}
\newcommand{\Rp}{\mathbb{R}_{>0}}
\newcommand{\ord}{\operatorname{ord}}
\newcommand{\code}[1]{\begingroup\def\_{\textunderscore\allowbreak}\texttt{#1}\endgroup}
\setlength{\emergencystretch}{3em}

\title[Five-body central configurations]{Finiteness of planar central configurations of five
bodies for all positive masses: On Smale's sixth problem}

\author{Zhe Chen}
\address{School of Wind Energy, Technical University of Denmark, 2800 Kongens Lyngby, Lyngby-Taarbæk Municipality, Copenhagen, Denmark}
\email{hitchenzhe@gmail.com}

\author{Zhicheng Tong}
\address{School of Mathematics, Jilin University, Changchun 130012, P. R. China}
\email{tongzc25@jlu.edu.cn}
\pdfinfo{/Author (Zhe Chen and Zhicheng Tong)}


\begin{document}

\begin{abstract}
{
We prove that for every choice of five positive masses the planar five-body problem has only
finitely many central configurations up to rotation, translation and scaling.  This answers
Smale's sixth problem for five bodies.  Albouy and Kaloshin proved the statement for all masses
outside a closed algebraic subset of codimension two of the mass space.  We keep their
classification of singular sequences and combine local boundary analysis with a global
argument on an algebraic curve.  We exclude Diagrams~1, 2, 4, 6 and~8 for complex germs,
Diagram~7 on curves through real configurations, and Diagram~3 on curves with a positive arc.  For Diagrams~5 and~9, direct order estimates show that the product
$\Delta$ of all squared mutual distances tends to zero.  On a hypothetical curve containing a
positive arc, a first pole of a four-product therefore gives a zero of $\Delta$.  On the
completed curve, $\Delta$ must then have a pole elsewhere.  The identity
$\Delta^4=\prod_{i<j}p_{ij}$ forces a four-product to have a pole there, contradicting the same
vanishing conclusion.  The computer-assisted inputs concern Diagrams~2, 3 and~6 and use exact
algebra, interval arithmetic and ball arithmetic.
}
\end{abstract}

\maketitle

\tableofcontents


\section{Introduction}\label{sec:intro}

A configuration of $n$ point masses $m_1,\dots,m_n>0$ at positions $q_1,\dots,q_n\in\R^2$ is a
\emph{central configuration} if there is $\lambda\in\R$ with
\[
  \lambda (q_k-c) = \sum_{l\neq k} m_l\,\frac{q_l-q_k}{|q_l-q_k|^3},\qquad k=1,\dots,n,
\]
where $c$ is the centre of mass.  Central configurations give the relative equilibria of the
planar $n$-body problem and the homographic solutions; they are invariant under rotations,
translations and scalings, and one counts them modulo these symmetries.  Euler and Lagrange
found the five classes for $n=3$; Moulton~\cite{Moulton} showed that there are exactly $n!/2$
collinear classes for every $n$.  Wintner~\cite{Wintner} and Smale~\cite{Smale} asked whether
the number of classes is finite for every choice of positive masses; this is Smale's sixth
problem.  Hampton and Moeckel~\cite{HM} proved finiteness for $n=4$.  Albouy and
Kaloshin~\cite{AK} gave a new proof for $n=4$ and proved the following for $n=5$.

\begin{theorem}[{\cite[Thm.~2]{AK}}]
There is a closed algebraic subset $A$ of codimension two such that for every
$(m_1,\dots,m_5)\in\Rp^5\setminus A$ there are finitely many positive normalized central
configurations of the planar five-body problem.
\end{theorem}
The following theorem removes the exceptional set, subject to the computational
statements listed in Appendix~\ref{app:verification}.  Positive normalized configurations
are defined in Section~\ref{sec:ak}.

\begin{theorem}\label{thm:main}
Assume the computational statements listed in Appendix~\ref{app:verification}.
For every $(m_1,\dots,m_5)\in\Rp^5$ there are finitely many positive normalized central
configurations of the planar five-body problem.  Equivalently, the planar five-body problem
has finitely many central configurations up to rotation, translation and scaling.
\end{theorem}
\begin{remark}
Recently, \cite{FiveT} independently established a conclusion identical to ours. It should be emphasized that our approach and that of \cite{FiveT} do not overlap, as the underlying proof strategies differ substantially. Therefore, both works should be viewed as independent contributions to Smale's sixth problem.
\end{remark}
\subsection*{Related work}
The generic finiteness of \cite{AK} for five bodies has been reproved by tropical
geometry~\cite{JL} and by an algorithmic version of the method of \cite{AK}~\cite{CC1,CC2},
which also reduces the planar six-body problem to finitely many diagrams.  In the spatial
five-body problem finiteness is known outside an explicit set of exceptional
masses~\cite{HJ}, and four-body central configurations are finite for all nonzero
masses~\cite{YZ}.  Computer-assisted methods classify the central configurations of five equal
masses~\cite{LS,MZ19} and, recently, give exact counts (between $294$ and $450$ classes) for
finitely many mass vectors on the loci (23), (29) and (32) of \cite[\S8]{AK}, the last of which
lies in the exceptional set~\cite{MZ26}; these computations treat individual mass vectors.
For all positive masses, finiteness is known for special classes such as the equilateral
pentagonal central configurations~\cite{DH}.  We are not aware of a previous proof of
Theorem~\ref{thm:main} on the whole exceptional set.

\subsection*{Strategy}
Albouy and Kaloshin embed the central configuration equations in a polynomial system and
study \emph{singular sequences}: unbounded sequences of complex normalized solutions.  To such a
sequence they attach a two-coloured \emph{diagram} recording which terms of the equations are
of maximal order, and they prove that every singular sequence approaches one of sixteen
diagrams~\cite[Fig.~11]{AK}.  If a mass vector admits infinitely many normalized central
configurations, the solutions contain an algebraic curve; along it some four-product of squared
distances is unbounded.  The boundedness argument of \cite[Prop.~8]{AK} excludes
Diagrams~10--16 at such a limit.  We retain Diagram~9 and treat it separately below.
The mass relations and the two-limit argument of \cite[\S\S8--9]{AK} provide the starting point
for our analysis.  Albouy and Kaloshin then push a second quantity to a second singular
limit on the same curve to obtain an independent relation, which yields the codimension-two set
$A$.

{
Our proof replaces the construction of $A$ by seven diagram exclusions and a vanishing result
for the product of all
squared distances in Diagrams~5 and~9.  For Diagrams~1--4 and~6--8 we introduce weighted blow-up
charts for the normalized equations.  The estimates of \cite[\S3]{AK} show that every germ
has a finite limit in the appropriate chart.  The boundary equations then give the necessary
mass conditions and the obstructions to continuation into the interior.  For Diagram~5 we
work directly with the two outer bodies and the shrinking triangle; its mass condition
comes from the isolated triangle.  We give the derivations of the mass conditions used
below, including those stated with short or no proofs in \cite[\S\S8.3--8.8]{AK}.
Diagram~7 requires no mass condition.  For Diagram~9, the distance estimates of
\cite[\S8.9]{AK} directly give the required vanishing of the total product.

These local arguments exclude Diagrams~1, 2, 4, 6 and~8 for complex germs.  They exclude
Diagram~7 with an unbounded four-product on a curve through a real point, and Diagram~3
on a curve through a positive point.  For every germ of Diagram~5 or~9, we prove that
$\Delta=\prod_{k<l}r_{kl}^2$ tends to zero.  This is enough for the global argument.
If a curve through a positive arc existed, a pole of a four-product could only have
Diagram~5 or~9 and would therefore give a zero of $\Delta$.  Thus $\Delta$ would be nonconstant
and would have a pole elsewhere on the completed curve.  The identity
$\Delta^4=\prod_{i<j}p_{ij}$ gives an unbounded four-product at this second pole.
The same exclusions leave only Diagrams~5 and~9 again, hence $\Delta\to0$, a contradiction.

The distinction between real and positive configurations is essential.  On a curve through a real
configuration, $U=I$ is a nonzero constant (Lemma~\ref{lem:U-constant}; compare \cite[Lemma~2]{AK}).
The inertia identity for Diagram~7 forces $I\to0$ whenever a four-product is unbounded.
For Diagram~3, the existence of a real point is insufficient: the modified Roberts continuum~\cite[p.~543]{AK}
is a curve of normalized central configurations at positive masses that contains an arc of
real configurations with signed inverse distances and carries a Diagram~3 germ.  The remaining
germs satisfy the identity $\delta_{15}+2\delta_{13}=0$ on the whole curve, which contradicts
positivity at any positive point.  The estimates for Diagrams~5 and~9 hold without a reality assumption
or a nonzero inertia.
}

\subsection*{Computer assistance}
{The remaining computer-assisted inputs are confined to Diagrams~2, 3 and~6:
exact polynomial and jet identities, full parameter interval and ball covers, and local
obstruction coefficients.  The required statements are identified in
Appendix~\ref{app:verification}.  Diagram~1 uses explicit triangle classification and rank
arguments; Diagram~4 uses a direct positivity argument in place of the sign census;
Diagram~5 uses force-order estimates, a Morse--Bott argument and the total product of
squared distances; and Diagram~8 uses exact distance closure and two scalar Taylor expansions.
These arguments, and those for Diagrams~7 and~9, are proved in the text using the stated results
of \cite{AK}.  Several boundary classifications and rank calculations in the other diagrams
are also proved directly.  Appendix~\ref{app:verification} distinguishes these proofs from
the computational statements on which the theorem still depends.

}\subsection*{Results used from \texorpdfstring{\cite{AK}}{Albouy and Kaloshin}}
We use without re-proving it the classification of singular sequences (Figure~11 of
\cite{AK} and its completeness), Proposition~8 of \cite{AK} restricted to Diagrams~10--16, the rules and estimates of
\cite[\S3]{AK}, and Lemma~1 and Proposition~1 of \cite{AK}.  These are listed precisely in
Section~\ref{sec:ak}.  For Diagram~9 we also use the local estimates of \cite[\S8.9]{AK},
recalled in Lemma~\ref{lem:d9-product}.  We re-derive the mass conditions of all diagrams that need one
(Diagrams~1--6 and~8), several of which \cite[\S\S8.3--8.8]{AK} state with short or no proofs,
and we prove directly that $U$ is constant along curves of solutions.

\subsection*{Organization}
{
Section~\ref{sec:ak} recalls the framework of \cite{AK}.  Section~\ref{sec:reduction} reduces
Theorem~\ref{thm:main} to a curve with an unbounded four-product and records the total-product
identity.  Section~\ref{sec:exclusion} states the boundary alternatives, proves the
Diagram~9 estimate and proves Theorem~\ref{thm:main} by taking two poles on that curve.  Sections~\ref{sec:d1}--\ref{sec:d8}
establish the diagram statements.  Appendix~\ref{app:verification} lists the computational statements used in the proof.
}


\section{The framework of Albouy and Kaloshin}\label{sec:ak}

\subsection{Normalized central configurations}
Fix masses $m=(m_1,\dots,m_5)\in\Rp^5$.  Following \cite[\S2]{AK} we use the complex coordinates
$z_k=x_k+iy_k$, $w_k=x_k-iy_k$ and inverse distances $\delta_{kl}=\delta_{lk}$, and write
$z_{kl}=z_l-z_k$, $w_{kl}=w_l-w_k$.  The normalized system (system (4) of \cite{AK}, written as
system (10) there) is
\begin{equation}\label{eq:system}
\begin{gathered}
  z_k = \sum_{l\ne k} m_l\,\delta_{kl}^3\,(z_k-z_l),\qquad
  w_k = \sum_{l\ne k} m_l\,\delta_{kl}^3\,(w_k-w_l)\qquad (k=1,\dots,5),\\
  \delta_{kl}^2\, z_{kl}\,w_{kl} = 1\quad (1\le k<l\le 5),\qquad z_{12} = w_{12}.
\end{gathered}
\end{equation}
Its zero set $A_m\subset\C^{20}$ is the set of \emph{normalized central configurations}.  A
solution is \emph{real} if $w_k=\bar z_k$ for all $k$ and \emph{positive} if moreover
$\delta_{kl}>0$ for all $k<l$ \cite[Def.~2]{AK}.  Positive normalized central configurations are
in two-to-one correspondence with central configurations modulo rotation, translation and
scaling, so Theorem~\ref{thm:main} asserts that $A_m$ has finitely many positive points.
We call the distance constraints in \eqref{eq:system} the \emph{ties}.
The map $\sigma(z,w,\delta)=(\bar w,\bar z,\bar\delta)$ preserves \eqref{eq:system} and fixes
{the real solutions.  System \eqref{eq:system} is symmetric under $z\leftrightarrow w$.  A permutation of the bodies,
with the masses permuted, preserves the equations apart from the choice of gauge.  After
relabelling, write the coordinates again as $(z,w,\delta)$.  A rotation
$(z,w)\mapsto(\lambda z,\lambda^{-1}w)$ with $\lambda^2=w_{12}/z_{12}$ restores the gauge
$z_{12}=w_{12}$.  It preserves $U$, $I$, every $r_{kl}^2$, and the diagram.  Along a germ,
if $\ord\lambda=a$, it adds $a$ to all $z$-orders and subtracts $a$ from all $w$-orders;
thus it preserves order differences within each colour, rather than absolute orders.
We use the term symmetry in this sense, allowing a finite algebraic cover and a finite
ramification of the parameter.

More precisely, on a relabelled irreducible curve $\Gamma$ consider
$\widetilde\Gamma=\{(q,\lambda)\in\Gamma\times\C^*:
\lambda^2z_{12}(q)=w_{12}(q)\}$.  The ties make $z_{12}$ and $w_{12}$ invertible regular
functions, so this is a finite unramified cover, and every irreducible component maps onto
$\Gamma$.  Any chosen germ lifts after a finite ramification.  Its component maps by the
rotation to an irreducible algebraic curve after taking the Zariski closure: some squared
distance is nonconstant by Lemma~\ref{lem:distances}, and remains nonconstant under the
cover and the rotation.  If $\Gamma$ contains a positive arc, each component contains a
local lift of that arc; there $|\lambda|=1$, so its rotated image is again a nonconstant
positive arc.  The same argument preserves the existence of a real point.  All squared
distances, four-products, $U$ and $I$ are unchanged, and a pole of a four-product remains
a pole on a lifted germ.}

\subsection{Potential and inertia}
Let $U=\sum_{k<l}m_km_l\delta_{kl}$ and $I=\sum_k m_kz_kw_k$.

\begin{lemma}\label{lem:U-constant}
Every solution of \eqref{eq:system} satisfies $U=I$, and $U$ is constant on every irreducible
curve $\Gamma\subset A_m$.
\end{lemma}

\begin{proof}
Let $B\subset\C^{20}$ be the set where $\delta_{kl}^2z_{kl}w_{kl}=1$ for all $k<l$, and let
$f=2U+I$.  On $B$ no $z_{kl}w_{kl}$ vanishes, so $B$ is smooth, $(z,w)$ are local coordinates
on $B$, and $d\delta_{kl}=-\frac12\delta_{kl}^3(w_{kl}\,dz_{kl}+z_{kl}\,dw_{kl})$.  Hence, on
$B$,
\[
  \frac{\partial f}{\partial z_k}=m_k\Bigl(w_k-\sum_{l\ne k}m_l\delta_{kl}^3(w_k-w_l)\Bigr),\qquad
  \frac{\partial f}{\partial w_k}=m_k\Bigl(z_k-\sum_{l\ne k}m_l\delta_{kl}^3(z_k-z_l)\Bigr),
\]
so every point of $A_m$ is a critical point of $f|_B$ (the gauge $z_{12}=w_{12}$ plays no
role).  The scaling $(z,w,\delta)\mapsto(sz,sw,s^{-1}\delta)$ preserves $B$ and maps $f$ to
$2s^{-1}U+s^2I$; differentiating at $s=1$ gives $0=-2U+2I$ at every critical point of $f|_B$, so
$U=I$ on $A_m$.  Let $\Gamma\subset A_m$ be an irreducible curve.  At a smooth point of
$\Gamma$ the tangent line lies in the tangent space of $B$, on which $df$ vanishes; so $f$ is
locally constant on the smooth locus of $\Gamma$, which is connected, and hence constant on
$\Gamma$.  As $f=3U$ on $A_m$, $U$ is constant on $\Gamma$.
\end{proof}
The identity $U=I$
is \cite[p.~541]{AK}, and \cite[Lemma~2]{AK} shows more: $U(A_m)$ is finite.

\begin{fact}[{\cite[Lemma~1]{AK}}]\label{fact:lemma1}
If $X\subset\C^N$ is a closed algebraic subset and $f$ a polynomial, then $f(X)$ is finite or
the complement of a finite set (``$f$ is dominating'').
\end{fact}

\subsection{Singular sequences and diagrams}
Let $\mathcal Z=(z_1,\dots,z_5,Z_{12},\dots,Z_{45})$ and $\mathcal W=(w_1,\dots,w_5,W_{12},
\dots,W_{45})$ with $Z_{kl}=z_{kl}\,\delta_{kl}^3$ and $W_{kl}=w_{kl}\,\delta_{kl}^3$; these are,
up to mass factors, the terms of the equations \eqref{eq:system}.  Let $Z=\max_p|\mathcal Z_p|$
and $W=\max_q|\mathcal W_q|$.  A \emph{singular sequence} \cite[Def.~3]{AK} is a sequence of
normalized central configurations with $Z$ or $W$ unbounded, extracted so that the maximal
components are fixed and $Z^{-1}\mathcal Z$, $W^{-1}\mathcal W$ converge.  Its \emph{diagram}
has a $z$-circle at $k$ if $\lim z_k/Z\ne0$, a $z$-stroke between $k$ and $l$ if
$\lim Z_{kl}/Z\ne 0$, and likewise in the second colour \cite[\S3]{AK}.  After multiplying $z$
by $a>0$ and $w$ by $a^{-1}$ one may assume $Z=W=\epsilon^{-2}$ with $\epsilon\to0$
\cite[p.~547]{AK}.  We write $a\prec b$, $a\preccurlyeq b$, $a\approx b$, $a\sim b$ as in
\cite[\S3]{AK}.

\begin{fact}[{\cite[\S3]{AK}}]\label{fact:rules}
Rules 1a--1e and 2a--2e and Estimates 1--2 of \cite[\S3]{AK} hold for every singular sequence.
\end{fact}

\begin{fact}[{\cite[Prop.~1]{AK}}]\label{fact:prop1}
Proposition 1 of \cite{AK} (localization of the origin in a cluster) holds.
\end{fact}

\begin{fact}[{\cite[\S7]{AK}}]\label{fact:fig11}
Every singular sequence of the five-body problem approaches, after renumbering the bodies and
possibly exchanging $z$ and $w$, one of the sixteen diagrams of \cite[Fig.~11]{AK}.
\end{fact}

For pairwise distinct $i,j,k,l,m$ with $\{i,j,k,l,m\}=\{1,\dots,5\}$ the \emph{four-product} is
$p_{ij}=r_{ij}^2r_{kl}^2r_{lm}^2r_{mk}^2$, where $r_{kl}^2=z_{kl}w_{kl}$.

\begin{fact}[{\cite[Prop.~8]{AK}}]\label{fact:prop8}
Along a singular sequence approaching one of Diagrams 10--16 every four-product is bounded.
\end{fact}

We use the results of \cite{AK} stated in Facts~AK1--AK5, together with the Diagram~9
estimates recalled in Lemma~\ref{lem:d9-product}.  Diagram~9 is treated separately there.
The mass conditions of
\cite[\S8]{AK} for Diagrams 1--6 and 8 are re-derived below
(Propositions~\ref{prop:d1-first}, \ref{prop:d2-first}, \ref{prop:d3-first},
\ref{prop:d45-first}, \ref{prop:d6-first} and \ref{prop:d8-first}).


\section{Reduction to germs with an unbounded four-product}\label{sec:reduction}

The first step is to show that the squared distances cannot all remain constant on a
curve of normalized solutions.  We then obtain a positive arc and pass to a branch at
infinity along which a four-product is unbounded.

\begin{lemma}\label{lem:distances}
Let $s\in(\C^*)^{10}$.  Only finitely many points of $A_m$ satisfy $z_{kl}w_{kl}=s_{kl}$ for all
$k<l$.
\end{lemma}

\begin{proof}
For any such point, $\delta_{kl}^2=s_{kl}^{-1}$, so $\delta$ takes finitely many values; fix one.
Multiplying the $z$-equation of body $k$ by $m_k$ and summing gives $\sum_km_kz_k=0$ by
antisymmetry, and likewise $\sum_km_kw_k=0$.

Put $v_k=(z_k-z_1,w_k-w_1)\in\C^2$ and let $B$ be
the symmetric bilinear form with $B((a,b),(a,b))=ab$.  Then $B(v_k-v_l,v_k-v_l)=s_{kl}$, so the
Gram matrix $B(v_k,v_l)=\frac12(s_{1k}+s_{1l}-s_{kl})$ is fixed.  The span $V$ of the $v_k$ is
nondegenerate for $B$: either $V=\C^2$, or $V=\C v_2$ with $B(v_2,v_2)=s_{12}\ne0$.

If
$(v'_k)$ is a second family with the same Gram matrix, its span $V'$ is nondegenerate by the
same argument, and $\sum_kc_kv_k=0$ holds if and only if $\sum_kc_kB(v_k,v_l)=0$ for all $l$,
and likewise for the $v'_k$; so the $v_k$ and the $v'_k$ satisfy the same linear relations, and
$v_k\mapsto v'_k$ is an isometry of nondegenerate subspaces; by Witt's theorem it extends to an
element of the orthogonal group of $B$, which consists of the maps $(a,b)\mapsto(\lambda a,
\lambda^{-1}b)$ and $(a,b)\mapsto(\lambda b,\lambda^{-1}a)$, $\lambda\in\C^*$.

The centre of
mass determines $(z_1,w_1)$ from $(v_k)$.  For a fixed reference family $(v^0_k)$ with first
components $a^0_2$ and second components $b^0_2$ of $v^0_2$ (so $a^0_2b^0_2=s_{12}\ne0$), the gauge
$z_{12}=w_{12}$ leaves $\lambda^2=b^0_2/a^0_2$ or $\lambda^2=a^0_2/b^0_2$: at most four values.
Hence there are finitely many solutions for each $\delta$.
\end{proof}

\begin{lemma}[Arc lemma]\label{lem:arc}
If $A_m$ has infinitely many positive points, there is an irreducible algebraic curve
$\Gamma\subset A_m$ such that
\begin{romanlist}
  \item $\Gamma$ contains a nonconstant real-analytic arc of positive points;
  \item $\sigma(\Gamma)=\Gamma$;
  \item $U=I$ is constant on $\Gamma$, with value $c\in(0,\infty)$.
\end{romanlist}
\end{lemma}

\begin{proof}
In the real coordinates $(x,y,\delta)\in\R^{20}$ the positive points form a semialgebraic set.
It has finitely many connected components, so one of them is infinite, hence of positive
dimension, and the curve selection lemma \cite[Thm.~2.5.5]{BCR} gives a nonconstant Nash arc in
it.  Its Zariski closure $\Gamma\subset A_m$ (in the complexified coordinates; passing to $(z,w)$ is linear) is an
irreducible curve, which proves (i).  Since \eqref{eq:system} has real coefficients, $\sigma(
\Gamma)$ is an irreducible curve containing the arc, so $\sigma(\Gamma)=\Gamma$.  By
Lemma~\ref{lem:U-constant}, $U=I$ is constant on $\Gamma$, and on the arc
$U=\sum m_km_l\delta_{kl}>0$.
\end{proof}

\begin{lemma}\label{lem:germ}
Let $\Gamma\subset A_m$ be an irreducible curve and $(q_n)$ a singular sequence on $\Gamma$.  There
are a Puiseux parametrization $\gamma$ of a branch of $\Gamma$ at infinity, $\gamma(t)\to\infty$ as
$t\to 0$, and $t_n\to0$ with $\gamma(t_n)$ a subsequence of $(q_n)$.  Every coordinate and every
component of $\mathcal Z$, $\mathcal W$ has a Puiseux expansion $c\,t^{\alpha}(1+o(1))$ along
$\gamma$ (or vanishes identically), and the diagram of every singular sequence extracted from
$(\gamma(t_n))$ is determined by these orders and leading coefficients: the circles and strokes
of colour $z$ are the components of $\mathcal Z$ of minimal order, and the relations $\prec$,
$\preccurlyeq$, $\approx$, $\sim$ of \cite[\S3]{AK} become relations between orders and leading
coefficients.  In particular a four-product unbounded along $(q_n)$ has negative order along
$\gamma$.
\end{lemma}

\begin{proof}
The closure of $\Gamma$ in a projective compactification has finitely many points at infinity,
each with finitely many analytic branches; a sequence tending to infinity has a subsequence on
one branch, and the normalization of the branch gives the parameter $t$.  Polynomial functions
of the coordinates are meromorphic in $t$, which gives the expansions.

If $\alpha$ is the least
order among the components of $\mathcal Z$, then $Z=\max_p|\mathcal Z_p|$ has order $\alpha$, and
$\mathcal Z_p/Z\to0$ exactly when $\ord\mathcal Z_p>\alpha$; this does not depend on the
subsequence.

The renormalization of \cite[p.~547]{AK} multiplies all $z$-quantities by one
positive factor $a(t)\approx|t|^{(\beta-\alpha)/2}$ and all $w$-quantities by $a^{-1}$, where
$\beta$ is the least order in $\mathcal W$.  This preserves every relation within one colour and
every product $z_{kl}w_{kl}$.  Replacing $a$ by the Puiseux series $\lambda(t)=(\mathcal W_q(t)/
\mathcal Z_p(t))^{1/2}$ built from maximal components, with $|\lambda|/a$ bounded above and
below, gives an analytic renormalization with the same relations; after reparametrizing
$t=s^N$, $\epsilon$ is comparable to a power of $s$.  A meromorphic function that is unbounded
along $t_n\to0$ has a pole.
\end{proof}

\begin{definition}\label{def:germ}
A \emph{Diagram-$r$ germ} is a Puiseux branch $\gamma$ of a curve in $A_m$ whose diagram, in the
sense of Lemma~\ref{lem:germ}, is Diagram $r$ of \cite[Fig.~11]{AK} after renumbering the bodies
and possibly exchanging $z$ and $w$.  It has an \emph{unbounded four-product} if some $p_{ij}$ has
negative order along $\gamma$.
\end{definition}

\begin{lemma}\label{lem:petersen}
On every irreducible curve $\Gamma\subset A_m$ some four-product is not constant.
\end{lemma}
\begin{proof}
Index the squared distances and the four-products by the ten pairs.  The pairs disjoint from
$ij$ are the sides of the triangle complementary to $ij$, so the exponent vector of $p_{ij}$ in
the variables $r_{ab}^2$ is the column $ij$ of $E=I+P$, where $P$ is the adjacency matrix of the
Petersen graph (two pairs are adjacent if they are disjoint).  The eigenvalues of $E$ are $4$,
$2$ (five times) and $-1$ (four times), so $\det E=128$, and $N=128E^{-1}$ is an integer
matrix.  As no $r_{ab}^2$ vanishes on $A_m$, $(r_{ab}^2)^{128}=\prod_{ij}p_{ij}^{N_{ab,ij}}$ there.
If every $p_{ij}$ were constant on $\Gamma$, every $r_{ab}^2$ would take finitely many values on
$\Gamma$ and hence be constant, and $\Gamma$ would be finite by Lemma~\ref{lem:distances}.
\end{proof}
This gives an algebraic proof of \cite[Prop.~9]{AK}, independent of the diagram analysis in
\cite[Fig.~11]{AK}.

{
\begin{lemma}[The total product]\label{lem:total-product}
Put $\Delta=\prod_{k<l}r_{kl}^2$.  On $A_m$, the function $\Delta$ is regular and nowhere zero,
and
\begin{equation}\label{eq:total-product}
  \prod_{i<j}p_{ij}=\Delta^4.
\end{equation}
\end{lemma}
\begin{proof}
The ties give $r_{kl}^2=z_{kl}w_{kl}\ne0$ on $A_m$.  Each squared distance occurs in its own
four-product and in the three four-products indexed by a pair disjoint from its endpoints.
It therefore occurs exactly four times in the product on the left of \eqref{eq:total-product}.
\end{proof}
}

\begin{proposition}[Reduction]\label{prop:reduction}
Let $m\in\Rp^5$.  If $A_m$ has infinitely many positive points, there are $r\in\{1,\dots,9\}$, a
curve $\Gamma$ as in Lemma~\ref{lem:arc}, and a Diagram-$r$ germ on $\Gamma$ with an unbounded
four-product.
\end{proposition}

\begin{proof}
Let $\Gamma$ be given by Lemma~\ref{lem:arc}.  By Lemma~\ref{lem:petersen} some four-product
$p_{ij}$ is not constant on $\Gamma$, and by Fact~\ref{fact:lemma1} it is dominating there.  A
sequence on $\Gamma$ with $p_{ij}\to\infty$ is unbounded, and a singular sequence extracted from
it approaches one of Diagrams 1--9 by Facts~\ref{fact:fig11} and~\ref{fact:prop8}.
Lemma~\ref{lem:germ} turns it into a germ.
\end{proof}

\begin{remark}\label{rem:levels}
The positivity of the arc must be distinguished from the signs of a branch at infinity.
Germs at infinity never lie on the positive arc itself: on a positive arc $U=I=c$ bounds every
$r_{kl}$ above and below.  Thus the signs of the boundary values alone do not determine whether the curve contains
positive configurations; we need identities valid on all of $\Gamma$.
{The diagram exclusions accordingly distinguish complex germs,
germs on curves through a real point, where $I=U$ is a positive constant, and germs on curves
with a positive arc.  Diagrams~5 and~9 instead give $\Delta\to0$ for all their complex germs;
this local statement is used in the global pole argument below.}
\end{remark}


{
\section{Boundary alternatives and the global argument}\label{sec:exclusion}

\begin{theorem}\label{thm:exclusion}
Let $m\in\Rp^5$, $r\in\{1,\dots,9\}$ and let $\gamma$ be a Diagram-$r$ germ with an unbounded
four-product on an irreducible curve $\Gamma\subset A_m$.
\begin{alphalist}
  \item If $r\in\{1,2,4,6,8\}$, no such germ exists.
  \item If $r=7$, then $\Gamma$ contains no real point.
  \item If $r=3$, then $\Gamma$ contains no arc of positive points.
  \item If $r\in\{5,9\}$, then $\ord_\gamma\Delta>0$, where $\Delta=\prod_{k<l}r_{kl}^2$.
\end{alphalist}
\end{theorem}


\begin{lemma}[Diagram 9]\label{lem:d9-product}
Let $m\in\Rp^5$ and let $\gamma$ be a Diagram~9 germ on an irreducible curve
$\Gamma\subset A_m$.  With the numbering of \cite[\S\S7.5.1.2 and~8.9]{AK}, the
$zw$-triangle is $123$ and the two $w$-edges are $14$ and $25$.  Then
\begin{equation}\label{eq:d9-distances}
\begin{gathered}
 r_{12},r_{13},r_{23}\approx\epsilon,\qquad
 r_{14},r_{25},r_{45}\approx1,\\
 r_{15},r_{24},r_{34},r_{35}\approx\epsilon^{-1/2}.
\end{gathered}
\end{equation}
Consequently $\Delta\approx\epsilon^2\to0$ and $\ord_\gamma\Delta>0$.
\end{lemma}

\begin{proof}
We normalize the germ as in Lemma~\ref{lem:germ} and take $r_{ij}=\delta_{ij}^{-1}$,
so that the signs of the square roots are those of the given solution.
We first exclude the clustering in \cite[\S8.9.1]{AK}.  In that case,
$r_{14}\approx\epsilon^{1/4}$, and equation~(41) of \cite{AK}, together with the
estimates following it, gives
\[
 \frac{m_1m_2}{r_{12}}+\frac{m_1m_3}{r_{13}}+
 \frac{m_2m_3}{r_{23}}=O(\epsilon^2).
\]
Every term of $U$ outside this triangle, apart from $m_1m_4/r_{14}$, is bounded.  Hence
\[
 U=\frac{m_1m_4}{r_{14}}+O(1),\qquad
 \left|\frac{m_1m_4}{r_{14}}\right|\longrightarrow\infty,
\]
contrary to Lemma~\ref{lem:U-constant}.  The same argument applies after the
simultaneous exchanges $1\leftrightarrow2$ and $4\leftrightarrow5$.

We are therefore in the case of \cite[\S8.9.2]{AK}, where both $w$-edges are
maximal.  The estimates for the $zw$-triangle and the maximal $w$-edges give
\[
 z_{12},z_{13},z_{23},w_{12},w_{13},w_{23}\approx\epsilon,
 \qquad z_{14},z_{25}\approx\epsilon^2,
 \qquad w_{14},w_{25}\approx\epsilon^{-2}.
\]
Taking differences gives $z_{45}\approx\epsilon$ and, for
$ij\in\{15,24,34,35\}$,
\[
 z_{ij}\approx\epsilon,\qquad w_{ij}\approx\epsilon^{-2},
 \qquad r_{ij}\approx\epsilon^{-1/2}.
\]
Thus all estimates in \eqref{eq:d9-distances}, except that for $r_{45}$, follow.
To obtain the remaining estimate, subtract the $z$-equations of bodies $4$ and $5$.
The terms from pairs $14$ and $25$ are $O(\epsilon^2)$, and those from pairs
$15,24,34,35$ are $O(\epsilon^{5/2})$.  Therefore
\[
 (z_4-z_5)\bigl(1-(m_4+m_5)r_{45}^{-3}\bigr)=O(\epsilon^2).
\]
Since $z_4-z_5\approx\epsilon$ and $m_4+m_5>0$, we obtain
$r_{45}^3\to m_4+m_5$, and hence $r_{45}\approx1$.

Multiplying the ten squared-distance estimates now gives
\[
 \Delta\approx(\epsilon^2)^3\,1^3\,(\epsilon^{-1})^4
 =\epsilon^2\longrightarrow0.
\]
The meromorphic function $\Delta$ thus has positive order along $\gamma$.
\end{proof}

In the numbering of Lemma~\ref{lem:d9-product}, the same estimates give
\[
 p_{15}=r_{15}^2r_{23}^2r_{34}^2r_{24}^2\approx\epsilon^{-1},\qquad
 p_{24}=r_{24}^2r_{13}^2r_{35}^2r_{15}^2\approx\epsilon^{-1}.
\]
Thus the boundedness assertion in the proof of \cite[Prop.~8]{AK} does not follow
for Diagram~9 from its distance estimates.  We use that proposition only for
Diagrams~10--16 and use Lemma~\ref{lem:d9-product} for Diagram~9.

\begin{proof}[Proof of Theorem~\ref{thm:main}]
Suppose that $A_m$ has infinitely many positive points.  Proposition~\ref{prop:reduction}
gives an irreducible curve $\Gamma$ containing a positive arc and a germ $\gamma_0$ on which
some four-product has a pole.  Theorem~\ref{thm:exclusion}(a)--(c) excludes every diagram
except Diagrams~5 and~9.  By part~(d), $\Delta$ has a zero at this end of $\Gamma$.

Let $\overline\Gamma$ be the smooth projective completion of the normalization of $\Gamma$.
All squared distances, four-products and $\Delta$ pull back to meromorphic functions on
$\overline\Gamma$.  By Lemma~\ref{lem:total-product}, $\Delta$ is nowhere zero on the affine
curve; its zero at $\gamma_0$ therefore implies that it is nonconstant.  A nonconstant
meromorphic function on the compact curve $\overline\Gamma$ has a pole, so choose a pole
of $\Delta$ and let $\gamma_1$ be the corresponding germ.  Taking orders in
\eqref{eq:total-product} shows that at least one four-product has a pole at $\gamma_1$.
As in the proof of Proposition~\ref{prop:reduction}, Facts~\ref{fact:fig11} and
\ref{fact:prop8} and Lemma~\ref{lem:germ} put this germ in one of Diagrams~1--9.
It lies on the same curve $\Gamma$ through the positive arc, so
Theorem~\ref{thm:exclusion}(a)--(c) again leaves only Diagrams~5 and~9.  Part~(d) now gives
$\ord_{\gamma_1}\Delta>0$, contradicting the choice of a pole.
The finite covers and rotations used to number and normalize a germ preserve $\Delta$,
the four-products and the existence of the positive arc, as explained in Section~\ref{sec:ak}.
\end{proof}

The statements of Theorem~\ref{thm:exclusion} for Diagrams~1--8 are proved in
Sections~\ref{sec:d1}--\ref{sec:d8}, and the Diagram~9 statement is Lemma~\ref{lem:d9-product}.
Theorem~\ref{thm:d5} and Lemma~\ref{lem:d9-product} prove part~(d) for every germ of
Diagram~5 or~9, without assuming an unbounded four-product.
Table~\ref{tab:exclusions} records the mass conditions, assumptions and conclusions for each
diagram.

\begin{table}[ht]

\centering
\small
\begin{tabular}{cllp{5.4cm}}
\hline
$r$ & mass condition & configurations & method and conclusion \\
\hline
1 & (23) on the triangle & complex & triangle virial and Morse--Bott estimate are incompatible \\
2 & $K_j$, $L_j$, $P_j$ & complex & isolation of every boundary label at orders 1--3,
with folds of corank one and two \\
3 & (29) $m_1m_3=m_2m_4$ & positive & orders 3--6; identity $\delta_{15}+2\delta_{13}=0$ \\
4 & (23) & complex & incompatibility of the force balances at orders one and three \\
5 & (23) & complex & direct force-order estimates give $\ord\Delta>0$; used in the global pole argument \\
6 & (32) $m_1=m_2$, $m_3=m_4$ & complex & finite family of wing charts \\
7 & none needed & real & inertia identity forces $I\to0$ \\
8 & (40) & complex & order-six and order-nine pairings; elimination on one curve of masses \\
9 & none needed & complex & the distance estimates give $\Delta\approx\epsilon^2\to0$ \\
\hline
\end{tabular}
\caption{The nine boundary arguments.  For Diagrams~5 and~9, the product of squared distances vanishes;
this conclusion is used in the global argument.  Numbers in parentheses refer to the equations of \cite[\S8]{AK}.}
\label{tab:exclusions}
\end{table}

Diagrams~1--4 and~6--8 use charts with coverage propositions and, where needed, mass conditions.
Diagram~5 uses direct estimates for the two outer bodies to prove the vanishing of the total
product, while Diagram~9 follows from its distance estimates.  The global argument combines
these vanishing statements with the other seven diagram results.
}
Throughout, a statement about a diagram is meant for every renumbering of the bodies (with the
masses permuted) and for both colourings; both operations are symmetries of
\eqref{eq:system}.


\section{Diagram 1}\label{sec:d1}

{Diagram~1 is one of the two disconnected diagrams remaining in \cite[\S7.1, p.~565]{AK}.
It consists of the $z$-stroke $12$, with $z$-circles at $1$ and $2$, and the $w$-triangle
with strokes $34,35,45$ and no circle.  There are no other strokes or circles.
We use the numbering of \cite[\S8.1]{AK}: $P=\{1,2\}$ is the pair and $T=\{3,4,5\}$ is the
triangle.  Indices $i$ refer to bodies of $P$, and $k,l$ to bodies of $T$.  Put
$M_P=m_1+m_2$, $M_T=m_3+m_4+m_5$, $Z_P=M_P^{-1}(m_1z_1+m_2z_2)$,
$Z_T=M_T^{-1}\sum_km_kz_k$, and define $W_P,W_T$ in the same way.  The mass-centre identities
$\sum_jm_jz_j=\sum_jm_jw_j=0$ on $A_m$ were proved in Lemma~\ref{lem:distances}; they give
$M_PZ_P+M_TZ_T=0=M_PW_P+M_TW_T$.  Write the coordinate equations of \eqref{eq:system} as
$E^z_j=E^w_j=0$, where}
\[
  E^z_j=z_j-\sum_{l\ne j}m_l\delta_{jl}^3(z_j-z_l),\qquad
  E^w_j=w_j-\sum_{l\ne j}m_l\delta_{jl}^3(w_j-w_l).
\]
Along a germ, $a\prec b$, $a\preccurlyeq b$ and $a\approx b$ in the sense of \cite[\S3]{AK} mean
$\ord a>\ord b$, $\ord a\ge\ord b$ and $\ord a=\ord b$ (Lemma~\ref{lem:germ}).

{We shall exclude every Diagram~1 germ in Theorem~\ref{thm:d1}.
The first difficulty is that a triangle body might approach the $w$-centre of the triangle
more closely than the other two.  The torque on the pair rules this out
(Lemma~\ref{lem:d1-central}), allowing every germ to be represented in one weighted chart
(Proposition~\ref{prop:d1-cover}).  The limiting triangle then has zero multiplier and its
masses satisfy relation~(23) of \cite[\S8.1]{AK} (Proposition~\ref{prop:d1-first}).
Near this limiting triangle, the potential vanishes quadratically in the forces.
The virial identity gives a lower vanishing order along the germ, and the two estimates
yield the required contradiction.}

\subsection{The chart}
{Let $\gamma$ be a Diagram~1 germ at $m\in\Rp^5$.  The symmetries of
\eqref{eq:system} allow us to renumber the bodies and exchange $z$ and $w$ so that its
diagram is as described above.  Since the coordinate equations are homogeneous of degree
one in $z$ and in $w$ separately, and the ties involve only $z_{jl}w_{jl}$, the transformation
$(z,w)\mapsto(cz,c^{-1}w)$ preserves both the equations and the diagram
\cite[p.~547]{AK}.
By Rule~2e \cite[p.~552]{AK}, the three $z$-separations of the triangle have the same order,
as do its three $w$-separations.  Choosing $c$ to be a power of the parameter $t$, we normalize
the $w$-separations of the triangle to have order $0$.  We do not impose the gauge
$z_{12}=w_{12}$.  Define $\lambda,\mu,\rho,\sigma$ by}
\[
  \ord z_{12}=-\lambda,\qquad \ord z_{kl}=\mu,\qquad \ord Z_P=\mu+\rho,\qquad \ord W_P=\sigma,
\]
with $\rho=+\infty$ (resp.\ $\sigma=+\infty$) if $Z_P$ (resp.\ $W_P$) vanishes identically, in
which case the terms carrying $\hat Z$ (resp.\ $\hat W$) below are absent.  Here $\sigma$ is an
exponent, not the conjugation of Section~\ref{sec:ak}.
After a finite ramification of $t$ we write
\begin{equation}\label{eq:d1-chart}
\begin{gathered}
  z_i=t^{\mu+\rho}\hat Z+t^{-\lambda}A_i,\qquad w_i=t^{\sigma}\hat W+t^{\lambda}B_i
  \qquad(i\in P),\\
  z_k=-\tfrac{M_P}{M_T}\,t^{\mu+\rho}\hat Z+t^{\mu}\zeta_k,\qquad
  w_k=-\tfrac{M_P}{M_T}\,t^{\sigma}\hat W+\eta_k\qquad(k\in T),\\
  \delta_{12}=D_{12},\qquad\delta_{ik}=t^{\lambda/2}D_{ik},\qquad\delta_{kl}=t^{-\mu/2}D_{kl},
\end{gathered}
\end{equation}
where $(A_1,A_2)=(m_2,-m_1)\,\pi$, $(B_1,B_2)=(m_2,-m_1)\,\varpi$ and
$\sum_km_k\zeta_k=\sum_km_k\eta_k=0$.  {The chart coordinates are
$\hat Z,\hat W,\pi,\varpi$, the mass-centred triangle offsets $\zeta=(\zeta_k)$,
$\eta=(\eta_k)$, and the ten $D_{jl}$.  To obtain finite limiting equations, we multiply
$E^z_i$, $E^w_i$, $E^z_k$ and $E^w_k$ by $t^{\lambda}$, $t^{-\lambda}$, $t^{\mu/2}$ and
$t^{3\mu/2}$, respectively; the ten ties retain weight $1$.
The exponents will belong to the open cone}
\begin{equation}\label{eq:d1-cone}
  C=\{\lambda>0,\ \mu>0,\ \sigma-\lambda>0,\ \rho+\tfrac32\mu>0,\ \rho+\lambda+\mu>0\}.
\end{equation}
For the triangle block put, with sums over $T$,
\[
  \phi_k=\sum_{l\ne k}m_lD_{kl}^3(\zeta_k-\zeta_l),\qquad
  \psi_k=\sum_{l\ne k}m_lD_{kl}^3(\eta_k-\eta_l),\qquad
  \chi_{kl}=D_{kl}^2(\zeta_l-\zeta_k)(\eta_l-\eta_k)-1,
\]
$U_T=\sum_{k<l}m_km_lD_{kl}$ and $I_T=\sum_km_k\zeta_k\eta_k$; note
$\sum_km_k\phi_k=\sum_km_k\psi_k=0$ identically.  The crossed forces are
$X^z_k=\sum_{i\in P}m_i\delta_{ik}^3(z_k-z_i)$ and $X^w_k=\sum_{i\in P}m_i\delta_{ik}^3(w_k-w_i)$.

{\begin{lemma}\label{lem:d1-chart}
After substitution of \eqref{eq:d1-chart} and multiplication by the prescribed weights,
each of the ten coordinate equations and ten ties is a finite sum $\sum_ec_e\,t^e$.
The exponents $e$ are linear forms in $(\lambda,\mu,\rho,\sigma)$, $e=0$ occurs, and each
$c_e$ is a polynomial in the chart coordinates.  Every exponent $e\ne0$ is a nonnegative
rational combination of the five forms in \eqref{eq:d1-cone}, with at least one positive
coefficient.  The constant coefficients $c_0$ are independent of $\hat Z$ and $\hat W$.
For $k\in T$,
\[
  t^{\mu/2}E^z_k=t^{\mu/2}(z_k-X^z_k)-\phi_k,\qquad
  t^{3\mu/2}E^w_k=t^{3\mu/2}(w_k-X^w_k)-\psi_k,
\]
and the exponents occurring in $t^{\mu/2}(z_k-X^z_k)$, respectively $t^{3\mu/2}(w_k-X^w_k)$, are
\[
\begin{gathered}
  \tfrac32\mu,\ \tfrac12(\lambda+\mu),\ \tfrac32(\lambda+\mu),\ \tfrac32\mu+\rho,\
  \tfrac32(\lambda+\mu)+\rho,\\
  \text{respectively}\\
  \tfrac32\mu,\ \tfrac32(\lambda+\mu),\ \tfrac52\lambda+\tfrac32\mu,\ \sigma+\tfrac32\mu,\
  \tfrac32(\lambda+\mu)+\sigma .
\end{gathered}
\]
\end{lemma}}
{\begin{proof}
Put $h=M_P/M_T$ and, for $i\in P$, $k\in T$, set
\[
 a_{ik}=A_i-t^{\lambda+\mu}\zeta_k
             +(1+h)t^{\lambda+\mu+\rho}\hat Z,\qquad
 b_{ik}=\eta_k-t^\lambda B_i-(1+h)t^\sigma\hat W.
\]
Then $z_i-z_k=t^{-\lambda}a_{ik}$ and $w_i-w_k=-b_{ik}$.
The weighted coordinate equations are
\begin{align*}
 t^\lambda E^z_i
 &=A_i(1-M_PD_{12}^3)+t^{\lambda+\mu+\rho}\hat Z
       -t^{3\lambda/2}\sum_km_kD_{ik}^3a_{ik},\\
 t^{-\lambda}E^w_i
 &=B_i(1-M_PD_{12}^3)+t^{\sigma-\lambda}\hat W
       +t^{\lambda/2}\sum_km_kD_{ik}^3b_{ik},\\
 t^{\mu/2}E^z_k
 &=-\phi_k+t^{3\mu/2}\zeta_k-h t^{3\mu/2+\rho}\hat Z
       +t^{(\lambda+\mu)/2}\sum_im_iD_{ik}^3a_{ik},\\
 t^{3\mu/2}E^w_k
 &=-\psi_k+t^{3\mu/2}\eta_k-h t^{\sigma+3\mu/2}\hat W
       -t^{3(\lambda+\mu)/2}\sum_im_iD_{ik}^3b_{ik}.
\end{align*}
The ties are, respectively,
$M_P^2D_{12}^2\pi\varpi-1$, $-D_{ik}^2a_{ik}b_{ik}-1$ and $\chi_{kl}$.
Every nonconstant power is a nonnegative rational combination of the five forms
defining $C$, with at least one positive coefficient.  Indeed,
$\sigma=\lambda+(\sigma-\lambda)$, and the terms involving $\rho$ are generated by
$\lambda+\mu+\rho$ and $3\mu/2+\rho$.  Expanding $a_{ik}$ and $b_{ik}$ gives the two
lists of exponents in the statement.  The constant terms give the limiting equations
described below.
\end{proof}

}{The constant terms $c_0$ are $A_i(1-M_PD_{12}^3)$ and
$B_i(1-M_PD_{12}^3)$ for the pair equations, $-\phi_k$ and $-\psi_k$ for the triangle
equations, and $M_P^2D_{12}^2\pi\varpi-1$, $-D_{ik}^2A_i\eta_k-1$ and $\chi_{kl}$ for the
ties within the pair, between the pair and triangle, and within the triangle.
We call their common zero set, in the coordinates other than $\hat Z,\hat W$, the
\emph{boundary fibre} $\mathcal B_1$.  On $\mathcal B_1$,
$D_{12}^3=M_P^{-1}$ and $D_{ik}^2A_i\eta_k=-1$, while
$(\zeta,\eta,D_{34},D_{35},D_{45})$ satisfies the \emph{zero-multiplier triangle system}
$\phi=\psi=0$, $\chi=0$.}

\subsection{Every germ lies in the chart}
{We now prove that every Diagram~1 germ has a finite limit in the chart
\eqref{eq:d1-chart}.  We use Rules 1e and 2e and Estimates 1 and 2 of \cite[\S3]{AK}
(Fact~\ref{fact:rules}), together with \cite[Prop.~1]{AK} (Fact~\ref{fact:prop1}).
The following identities express the forces on each component and the torque on the pair.}

\begin{lemma}\label{lem:d1-syzygies}
{The following polynomial identities hold:}
\begin{gather*}
  m_1E^w_1+m_2E^w_2=m_1w_1+m_2w_2-\sum_{i\in P,\,k\in T}m_im_k\delta_{ik}^3(w_i-w_k),\\
  \sum_{k\in T}m_kE^z_k=\sum_{k\in T}m_kz_k-\sum_{i\in P,\,k\in T}m_im_k\delta_{ik}^3(z_k-z_i),\\
  T_P:=\sum_{i\in P,\,k\in T}m_im_k\delta_{ik}^3(z_iw_k-w_iz_k)
  =-\sum_{i\in P}m_i\big(w_iE^z_i-z_iE^w_i\big).
\end{gather*}
In particular the torque $T_P$ of the triangle on the pair vanishes on $A_m$.
\end{lemma}

\begin{proof}
{The internal force terms cancel in pairs within $P$ and within $T$.
The same cancellation in the weighted coordinate equations gives the torque identity.}
\end{proof}

\begin{lemma}\label{lem:d1-orders}
Let $\gamma$ be a Diagram~1 germ, normalized as above.
\begin{romanlist}
  \item $\ord\delta_{12}=0$ and $\ord w_{12}=\lambda>0$; the least order of the components of
  $\mathcal Z$ is $\ord z_1=\ord z_2=\ord Z_{12}=-\lambda$.
  \item $\mu>0$, $\ord\delta_{kl}=\ord Z_{kl}=-\mu/2$, and $\lambda>\mu/2$.
  \item $\ord z_{ik}=-\lambda$, $\ord w_{ik}<\lambda$ and $\ord w_i\ge\lambda$; hence
  $\ord W_T=\ord W_P\ge\lambda$.
  \item For $k\in T$ let $o_k$ be the coefficient of $t^0$ in $w_k-W_T$.  At most one $o_k$
  vanishes.  If $o_k\ne0$, then $\ord w_{ik}=0$ and $\ord\delta_{ik}=\lambda/2$.  If $o_k=0$, then
  $f=\ord(w_k-W_T)$ satisfies $0<f<\lambda$, and $\ord w_{ik}=f$, $\ord\delta_{ik}=(\lambda-f)/2$.
\end{romanlist}
\end{lemma}

\begin{proof}
(i) The isolated component $\{1,2\}$ of the $z$-diagram has circled vertices, so it has a
maximal $z$-stroke by Rule~1e \cite[p.~546]{AK}, which can only be $12$.  Estimate~2
\cite[p.~548]{AK} gives $r_{12}\approx1$, so $\ord\delta_{12}=0$ and $\ord w_{12}=-\ord
z_{12}=\lambda$.  Estimate~1 \cite[p.~548]{AK} with $z$ and $w$ exchanged gives
$w_{12}\approx\epsilon^2$, and Estimate~2 with $z$ and $w$ exchanged gives $w_{34}\succcurlyeq\epsilon$
for the $w$-stroke $34$.  Hence $w_{12}\prec w_{34}$; a comparison within one colour is not
affected by the normalizations, so $\lambda>\ord w_{34}=0$.  Finally $1,2$ are $z$-circled and
$12$ is a $z$-stroke, and $\ord Z_{12}=\ord z_{12}+3\ord\delta_{12}=-\lambda$.

(ii) Bodies $1,2$ carry no $w$-stroke, so $T$ is an isolated component of the $w$-diagram
without $w$-circles.  By Rule~1e it has no maximal $w$-stroke.  Estimate~2 with $z$ and $w$
exchanged gives $r_{kl}\preccurlyeq1$, with $r_{kl}\approx1$ only for a maximal $w$-edge, so
$r_{kl}\prec1$; thus $\mu=\ord(z_{kl}w_{kl})=2\ord r_{kl}>0$, $\ord\delta_{kl}=
-\mu/2$ and $\ord Z_{kl}=\mu-\frac32\mu=-\frac12\mu$.  Diagram~1 has no $z$-stroke inside $T$, so
$\ord Z_{kl}>-\lambda$, that is, $\lambda>\mu/2$.

(iii) The bodies of $T$ are not $z$-circled, so $\ord z_k>-\lambda=\ord z_i$ and $\ord
z_{ik}=-\lambda$.  Since $Z_{ik}^2=(z_{ik}w_{ik}^3)^{-1}$, $\ord Z_{ik}=\frac12(\lambda-3\ord
w_{ik})$, and the absence of a $z$-stroke $ik$ gives $\ord Z_{ik}>-\lambda$, i.e.\ $\ord
w_{ik}<\lambda$.  By Estimate~1 ($z$ and $w$ exchanged) $w_{2k}\succ\epsilon^2\approx w_{12}$,
because $12$ is the only maximal $z$-edge, and $z_{2k}\approx z_{12}$.  So
\cite[Prop.~1, p.~559]{AK} with $z$ and $w$ exchanged, applied to bodies $1,2$ as in the proof of
\cite[Prop.~4, p.~564]{AK}, gives $w_1,w_2\preccurlyeq w_{12}$, i.e.\ $\ord w_i\ge\lambda$.  This
also follows from Lemma~\ref{lem:d1-syzygies}: every term of $M_PW_P=\sum m_im_k\delta_{ik}^3
(w_i-w_k)$ has order $\frac12(3\lambda-\ord w_{ik})>\lambda$, and $w_i-W_P$ has order
$\ord w_{12}=\lambda$.

(iv) Since $\sum_km_ko_k=0$, and the $o_k$ do not all vanish because $\ord w_{kl}=0$, at most one
$o_k$ vanishes.  Now $w_k-w_i=(w_k-W_T)+(W_T-w_i)$ with $\ord(W_T-w_i)\ge\lambda>0$ by (iii).  If
$o_k\ne0$ this gives $\ord w_{ik}=0$.  If $o_k=0$ then $f>0$; $f\ge\lambda$ would give $\ord
w_{ik}\ge\lambda$, against (iii), so $f<\lambda$ and $\ord w_{ik}=f$.  In both cases
$\ord\delta_{ik}=-\frac12(\ord z_{ik}+\ord w_{ik})$.
\end{proof}

{\begin{lemma}\label{lem:d1-torque}
Let $\lambda,\mu,f>0$ with $f<\lambda$ and let $\rho,\sigma$ be real.  Modify
\eqref{eq:d1-chart} by $(\eta_3,\eta_4,\eta_5)=(t^f\omega,\ m_5\vartheta,\
-m_3m_5^{-1}t^f\omega-m_4\vartheta)$ and $\delta_{i3}=t^{(\lambda-f)/2}d_i$ $(i=1,2)$, all other
entries unchanged.  Then
\[
  T_P=t^{(\lambda-f)/2}\Big(m_3\sum_{i=1,2}m_id_i\,\big(d_i^2A_i\omega\big)+\sum_{g\in G}c_g\,t^g\Big),
\]
where the $c_g$ are polynomials in the chart coordinates and $G$ consists of the thirteen
exponents
\[
\begin{gathered}
  \tfrac12f,\ \tfrac32f,\ \sigma-f,\ \sigma+\tfrac12f,\ 2\lambda+\mu-f,\
  2\lambda+\mu+\tfrac12f,\ \lambda+\mu+\rho,\ \lambda+\mu+\sigma-f,\ 2\lambda+\mu+\rho-f,\\
  \lambda+\mu+\rho+\tfrac12f,\ \lambda+\mu+\sigma+\tfrac12f,\ 2\lambda+\mu+\rho+\tfrac12f,\
  \lambda+\mu+\rho+\tfrac32f,
\end{gathered}
\]
each a combination of $\lambda,\ \mu,\ \sigma-\lambda,\ \rho+\lambda+\mu,\ f,\ \lambda-f$ with
nonnegative rational coefficients, at least one of them positive.  Moreover, if
$d_i^2A_i\omega=-1$ for $i=1,2$, then $(m_1d_1)^2-(m_2d_2)^2=-(m_1^3+m_2^3)/(m_1A_1\omega)$.
\end{lemma}}
{\begin{proof}
Put $h=M_P/M_T$.  The determinant in each summand of the torque is
\begin{align*}
 z_iw_k-w_iz_k
 ={}&t^{-\lambda}A_i\eta_k-h t^{\sigma-\lambda}A_i\hat W
       +t^{\mu+\rho}\hat Z\eta_k-t^{\mu+\sigma}\hat W\zeta_k\\
   &+h t^{\lambda+\mu+\rho}B_i\hat Z
       -t^{\lambda+\mu}B_i\zeta_k.
\end{align*}
For $k=3$, multiply this expression by
$t^{3(\lambda-f)/2}d_i^3$ and use $\eta_3=t^f\omega$.
After removing the common factor $t^{(\lambda-f)/2}$, the first term
is $d_i^3A_i\omega$, and the other five exponents are
\[
 \sigma-f,\quad\lambda+\mu+\rho,\quad
 \lambda+\mu+\sigma-f,\quad2\lambda+\mu+\rho-f,\quad2\lambda+\mu-f.
\]
For $k=4,5$, multiply instead by $t^{3\lambda/2}D_{ik}^3$;
the constant and $t^f$ parts of $\eta_k$ give exactly the remaining
eight exponents in $G$.  This proves the displayed expansion.
Each exponent has the asserted nonnegative representation: the only
subtractions needing attention are
$\sigma-f=(\sigma-\lambda)+(\lambda-f)$,
$2\lambda+\mu-f=\lambda+\mu+(\lambda-f)$ and
$2\lambda+\mu+\rho-f=(\rho+\lambda+\mu)+(\lambda-f)$;
the remaining cases follow by adding nonnegative multiples of
$\lambda,\mu,f$ to these or to $\rho+\lambda+\mu$.
Finally $A_2=-m_1A_1/m_2$, so the ties give directly
\[
 (m_1d_1)^2-(m_2d_2)^2
 =-\frac{m_1^2}{A_1\omega}+\frac{m_2^2}{A_2\omega}
 =-\frac{m_1^3+m_2^3}{m_1A_1\omega}.
\]
\end{proof}

}\begin{lemma}\label{lem:d1-central}
Let $m\in\Rp^5$.  Along every Diagram~1 germ all three leading triangle $w$-offsets $o_k$ are
nonzero.
\end{lemma}

\begin{proof}
Suppose not.  By Lemma~\ref{lem:d1-orders}(iv), after renumbering $T$, $o_3=0\ne o_4,o_5$,
$0<f<\lambda$, $\ord\delta_{i3}=\frac12(\lambda-f)$ and $\ord\delta_{ik}=\frac12\lambda$ for $k=4,5$.
In the first two identities of Lemma~\ref{lem:d1-syzygies} the crossed terms have orders
$\frac12(3\lambda-\ord w_{ik})$ and $\frac12(\lambda-3\ord w_{ik})$, so
$\sigma\ge\frac12(3\lambda-f)$ and $\mu+\rho\ge\frac12(\lambda-3f)$, that is,
$\sigma-\lambda\ge\frac12(\lambda-f)>0$ and $\rho+\lambda+\mu\ge\frac32(\lambda-f)>0$.
Hence $\gamma$ is described by the chart of Lemma~\ref{lem:d1-torque}, with chart coordinates of
order $\ge0$, of order $0$ for $d_i,A_i,\omega$ (so these are nonzero at $t=0$), and every
exponent in $G$ is positive.
As $\ord(Z_T-Z_P)=\mu+\rho>-\lambda$ and $\sigma>\lambda>f$, we have $z_3-z_i=-t^{-\lambda}A_i(1+o(1))$
and $w_3-w_i=t^f\omega(1+o(1))$, so the ties $\delta_{i3}^2(z_3-z_i)(w_3-w_i)=1$ give
$d_i^2A_i\omega=-1$ at $t=0$.
By Lemma~\ref{lem:d1-torque}, $t^{-(\lambda-f)/2}T_P$ tends along $\gamma$ to
$-m_3(m_1d_1+m_2d_2)$, evaluated at $t=0$, and this is nonzero because
$(m_1d_1)^2-(m_2d_2)^2=-(m_1^3+m_2^3)/(m_1A_1\omega)\ne0$.  But $T_P=0$ along $\gamma$ by
Lemma~\ref{lem:d1-syzygies}.
\end{proof}

\begin{proposition}[Coverage]\label{prop:d1-cover}
Let $m\in\Rp^5$ and let $\gamma$ be a Diagram~1 germ, normalized as above.
\begin{romanlist}
  \item $\lambda>\mu/2>0$.
  \item Every leading triangle $w$-offset $o_k$ is nonzero; hence $\ord w_{ik}=0$ and
  $\ord\delta_{ik}=\lambda/2$ for $i\in P$, $k\in T$.
  \item $\sigma\ge\frac32\lambda$ and $\mu+\rho\ge\frac12\lambda$.  In particular
  $(\lambda,\mu,\rho,\sigma)\in C$, with $\sigma-\lambda\ge\frac12\lambda$,
  $\rho+\frac32\mu\ge\frac12(\lambda+\mu)$ and $\rho+\lambda+\mu\ge\frac32\lambda$.
  \item Along $\gamma$ every chart coordinate of \eqref{eq:d1-chart} converges as $t\to0$; the
  limits of $\pi$, $\varpi$, of the $D_{jl}$, of the $\eta_k$ and of the differences
  $\zeta_k-\zeta_l$, $\eta_k-\eta_l$ $(k\ne l)$ are nonzero; and the limit point lies in the
  boundary fibre $\mathcal B_1$.
\end{romanlist}
\end{proposition}

\begin{proof}
(i) is Lemma~\ref{lem:d1-orders}(i),(ii), and (ii) follows from Lemmas~\ref{lem:d1-central}
and~\ref{lem:d1-orders}(iv).  (iii) On $\gamma$, Lemma~\ref{lem:d1-syzygies} gives
$M_PW_P=\sum m_im_k\delta_{ik}^3(w_i-w_k)$ and $M_TZ_T=\sum m_im_k\delta_{ik}^3(z_k-z_i)$.  By (ii)
and Lemma~\ref{lem:d1-orders}(iii) the terms of the first sum have order $\frac32\lambda$ and
those of the second $\frac32\lambda-\lambda=\frac12\lambda$; as $\ord Z_T=\ord Z_P=\mu+\rho$ the
two bounds follow.  {Substitution gives the three remaining inequalities.}
(iv) By the definitions of $\lambda,\mu,\rho,\sigma$, by Lemma~\ref{lem:d1-orders} and by (ii),
every chart coordinate is a Puiseux series of order $\ge0$, of order exactly $0$ for the listed
ones; so all converge.  By Lemma~\ref{lem:d1-chart} and (iii), each weighted equation equals
$c_0+\sum_{e>0}c_et^e$ along $\gamma$ with bounded $c_e$; it vanishes identically, so $c_0$
vanishes at the limit point.
\end{proof}

\begin{remark}\label{rem:d1-ak}
{Lemma~\ref{lem:d1-orders}} and Proposition~\ref{prop:d1-cover} give the clustering scheme
$z:1\ldots345\ldots2$, $w:3\ldots12\ldots5\ldots4$, which
\cite[\S8.1, p.~572]{AK} derives from Rule~2e, Estimate~2 and \cite[Prop.~4]{AK}.
Lemma~\ref{lem:d1-central} is the torque argument of the proof of \cite[Prop.~4]{AK}, carried out
in the chart with every competing term controlled.  {The chart and order estimates above suffice for the proof of
Theorem~\ref{thm:d1}; no further use of the clustering scheme of \cite[\S8.1]{AK} is needed.}
\end{remark}

\subsection{The limiting triangle and its masses}

{\begin{lemma}\label{lem:d1-shape}
Let $m_3,m_4,m_5>0$ and let $(\zeta,\eta,D_{34},D_{35},D_{45})$, with
$\sum_km_k\zeta_k=\sum_km_k\eta_k=0$, solve the zero-multiplier triangle system.  Put
$\alpha_{kl}=\zeta_l-\zeta_k$ and $\beta_{kl}=\eta_l-\eta_k$ for $k<l$ in $T$.  Then
\begin{alphalist}
  \item $\beta_{kl}=\Lambda^{-1}\alpha_{kl}$ for one $\Lambda\in\C^*$, and $(D_{kl}\alpha_{kl})^2=
  \Lambda$ for all $k<l$;
  \item $(D_{34}\alpha_{34},D_{35}\alpha_{35},D_{45}\alpha_{45})=s\,(-1,1,-1)$ with $s^2=\Lambda$;
  the other six sign patterns admit no solution;
  \item $m_5\alpha_{34}^2=m_3\alpha_{45}^2=m_4\alpha_{35}^2$ and $Q(m_3,m_4,m_5)=0$, where
  \[
    Q=2m_3m_4m_5(m_3+m_4+m_5)-m_3^2m_4^2-m_4^2m_5^2-m_5^2m_3^2 .
  \]
\end{alphalist}
\end{lemma}}
{\begin{proof}
Let $M=\operatorname{diag}(m_3,m_4,m_5)$ and let $B$ be the symmetric matrix
\[
 B_{kl}=-m_km_lD_{kl}^3\quad(k\ne l),\qquad
 B_{kk}=\sum_{l\ne k}m_km_lD_{kl}^3.
\]
Then $B\mathbf1=0$, $B\zeta=M\phi=0$ and $B\eta=M\psi=0$.
All the $D_{kl}$ are nonzero by the ties, so $B\ne0$.
The mass-centred vectors $\zeta,\eta$ cannot be independent.
If they were, $\mathbf1,\zeta,\eta$ would form a basis of $\C^3$, forcing $B$ to vanish.
Thus $\eta=\Lambda^{-1}\zeta$ with $\Lambda\ne0$, and the ties give (a).

Choose $s^2=\Lambda$ and write $D_{kl}\alpha_{kl}=s\varepsilon_{kl}$,
where $\varepsilon_{kl}\in\{1,-1\}$.  Put
$u=\alpha_{45}/\alpha_{34}$ and $v=\alpha_{35}/\alpha_{34}=1+u$.
The equations $\phi_4=\phi_5=0$ give
\[
 u^2=\frac{m_5}{m_3}\frac{\varepsilon_{45}}{\varepsilon_{34}},
 \qquad
 v^2=-\frac{m_5}{m_4}\frac{\varepsilon_{35}}{\varepsilon_{34}}.
\]
Both squares are nonzero real numbers.  Thus $u$ is real or purely imaginary.
The latter possibility is excluded because $(1+u)^2$ would then have nonzero
imaginary part.  Hence $u,v$ are nonzero real numbers.  Positivity of the masses now implies
$\varepsilon_{45}=\varepsilon_{34}=-\varepsilon_{35}$.  Changing the sign
of $s$ if necessary gives (b), and the two displayed equations give the
equalities in (c).

With $x_k=m_k^{-1/2}>0$, the equalities just obtained and
$\alpha_{35}=\alpha_{34}+\alpha_{45}$ imply that one of $x_3,x_4,x_5$
is the sum of the other two.  The factorization
\[
 Q=(m_3m_4m_5)^2(x_3+x_4+x_5)(-x_3+x_4+x_5)
       (x_3-x_4+x_5)(x_3+x_4-x_5),
\]
therefore gives $Q=0$.  For positive masses, the same factorization shows that
$Q=0$ is equivalent to relation (23).
\end{proof}

}The torus $\mathbb T=\{(\theta_1,\theta_2,\theta_3)\in(\C^*)^3:\theta_3^2=\theta_1\theta_2\}$ acts by
$(\zeta,\eta,D)\mapsto(\theta_1\zeta,\theta_2\eta,\theta_3^{-1}D)$.  It multiplies $\phi$ by
$\theta_1\theta_3^{-3}$, $\psi$ by $\theta_2\theta_3^{-3}$ and $I_T$ by $\theta_1\theta_2$, and it
fixes $\chi$.

{\begin{lemma}\label{lem:d1-mb}
Let $b,c>0$, $(m_3,m_4,m_5)=(b^2c^2,\,b^2(b+c)^2,\,c^2(b+c)^2)$ and $\kappa>0$.  Let $y^*$ be the
point with $\eta=\zeta$, $\sum_km_k\zeta_k=0$, $\zeta_4-\zeta_3=\kappa m_5^{-1/2}$,
$\zeta_5-\zeta_3=-\kappa m_4^{-1/2}$, $D_{34}=m_5^{1/2}/\kappa$, $D_{35}=m_4^{1/2}/\kappa$ and
$D_{45}=-m_3^{1/2}/\kappa$.  In the coordinates $y=(\zeta_3,\zeta_4,\eta_3,\eta_4,D_{34},D_{35},
D_{45})$ let $F=(\phi_3,\phi_4,\psi_3,\psi_4,\chi_{34},\chi_{35},\chi_{45})$.  Then $F(y^*)=0$; the
Jacobian matrix $F'(y^*)$ has rank $5$, and its kernel is spanned by the generators
$(\zeta_3,\zeta_4,0,0,-\frac12D)$ and $(0,0,\eta_3,\eta_4,-\frac12D)$ of the $\mathbb T$-action; and
\[
  I_T(y^*)=\frac{\kappa^2}{m_3+m_4+m_5}\Big(\frac{m_3m_4}{m_5}+\frac{m_3m_5}{m_4}
  +\frac{m_4m_5}{m_3}\Big)>0 .
\]
\end{lemma}}
{\begin{proof}
Write $q=\zeta=\eta$ at $y^*$.  The mass relation is
$m_3^{-1/2}=m_4^{-1/2}+m_5^{-1/2}$, and hence
$q_5-q_4=-\kappa m_3^{-1/2}$.  All three ties hold and
$D_{34}(q_4-q_3)=1$, $D_{35}(q_5-q_3)=-1$,
$D_{45}(q_5-q_4)=1$.
Substitution in the two force identities used in the proof of
Lemma~\ref{lem:d1-shape} gives $\phi_4=\phi_5=0$; the mass-centre identity
gives $\phi_3=0$, and $\psi=\phi$.  Thus $F(y^*)=0$.

Use the matrix $B$ from that proof and put $L=M^{-1}B$.
Since $B\mathbf1=Bq=0$, the vectors $\mathbf1,q$ are independent and
$B\ne0$, we have $\operatorname{rank}B=1$ and
$\ker B=\operatorname{span}\{\mathbf1,q\}$.
On the mass-centred plane
$H=\{x:\sum_km_kx_k=0\}$, the map $L$ has rank $1$ and kernel $\C q$.

The three tie differentials are independent, since
$\partial\chi_{kl}/\partial D_{kl}=2/D_{kl}\ne0$ and the other two
tie rows have zero entries in this column.  Their zero set is therefore
a smooth four-dimensional manifold, with local coordinates $(\zeta,\eta)\in H^2$.
On its tangent space at $y^*$,
\[
 dD_{kl}=-\frac{D_{kl}}{2(q_l-q_k)}
   \bigl(d\zeta_l-d\zeta_k+d\eta_l-d\eta_k\bigr),
\]
and differentiation of the forces gives
\[
 \begin{pmatrix}d\phi\\d\psi\end{pmatrix}
 =-\frac12
 \begin{pmatrix}L&3L\\3L&L\end{pmatrix}
 \begin{pmatrix}d\zeta\\d\eta\end{pmatrix}.
\]
The scalar $2\times2$ matrix here is invertible, so this map has rank $2$
and kernel $\C(q,0)\oplus\C(0,q)$.  Since either two components determine
each mass-centred force vector, the full Jacobian $F'(y^*)$ has rank
$3+2=5$.  The formula for $dD$ identifies its two kernel generators with
the stated infinitesimal torus actions.

The elementary mass-centre identity
\[
 \sum_km_kq_k^2=\frac1{m_3+m_4+m_5}
       \sum_{k<l}m_km_l(q_l-q_k)^2
\]
gives the displayed value of $I_T(y^*)$, which is positive because $q$ is
a nonzero real vector.  Conversely every positive solution of the mass
relation has the stated $(b,c)$ parametrization: take
$b/c=\sqrt{m_4/m_5}$ and $bc=\sqrt{m_3}$.
\end{proof}
By equivariance, the same rank and kernel statements hold at every point of
$O=\mathbb T\cdot y^*$, and $I_T\ne0$ on $O$.

}\begin{proposition}[The limiting mass condition]\label{prop:d1-first}
Let $m\in\Rp^5$ and let $\gamma$ be a Diagram~1 germ.  Then $Q(m_3,m_4,m_5)=0$, that is, the
triangle masses satisfy relation (23) of \cite[\S8.1]{AK}.  After renumbering $T$ so that
$m_3^{-1/2}=m_4^{-1/2}+m_5^{-1/2}$, the limit $y_0$ of
$y(t)=(\zeta_3,\zeta_4,\eta_3,\eta_4,D_{34},D_{35},D_{45})(t)$ lies on the orbit $O$ of
Lemma~\ref{lem:d1-mb}.
\end{proposition}

\begin{proof}
By Proposition~\ref{prop:d1-cover}(iv) the limit lies in $\mathcal B_1$, so its triangle part
solves the zero-multiplier triangle system.  In the notation of Lemma~\ref{lem:d1-shape},
$\phi_4=m_3\alpha_{34}D_{34}^3-m_5\alpha_{45}D_{45}^3$ and
$\phi_5=m_4\alpha_{45}D_{45}^3+m_3\alpha_{35}D_{35}^3$, and $m_3\phi_3=-m_4\phi_4-m_5\phi_5$;
likewise for $\psi$.  So $\phi=0$ and $\psi=0$ are relations (22) and (21) of \cite[\S8.1]{AK} in
cross-multiplied form, and Lemma~\ref{lem:d1-shape} gives $Q=0$.
After the renumbering, Lemma~\ref{lem:d1-shape}(c) and $\alpha_{35}=\alpha_{34}+\alpha_{45}$
leave only $(\alpha_{34},\alpha_{45},\alpha_{35})=\kappa(x_5,-x_3,-x_4)$ with $\kappa\in\C^*$, and
Lemma~\ref{lem:d1-shape}(a),(b) then determine $\eta$ and $D$ from $(\kappa,s)$.  The torus
$\mathbb T$ acts transitively on these pairs $(\kappa,s)$, and $y^*$ is the solution with
$\kappa>0$, $s=-1$.  Hence the solutions form the single orbit $O$, and $y_0\in O$.
\end{proof}

\subsection{Excluding Diagram~1}

{\begin{lemma}\label{lem:d1-virial}
For arbitrary $\zeta,\eta\in\C^3$ and $D_{34},D_{35},D_{45}\in\C$,
\begin{gather*}
  \sum_{k\in T}m_k\zeta_k\psi_k=U_T+\sum_{k<l}m_km_lD_{kl}\chi_{kl},\\
  2\,dU_T+\sum_{k\in T}m_k\big(\psi_k\,d\zeta_k+\phi_k\,d\eta_k\big)
  =\sum_{k<l}m_km_l\big(D_{kl}\,d\chi_{kl}-2\chi_{kl}\,dD_{kl}\big).
\end{gather*}
\end{lemma}}
{\begin{proof}
For a single edge put $\alpha=\zeta_l-\zeta_k$,
$\beta=\eta_l-\eta_k$, $D=D_{kl}$ and $\chi=D^2\alpha\beta-1$.
Pairing the $k$ and $l$ terms in the force sum gives
$m_km_lD^3\alpha\beta=m_km_lD(1+\chi)$, which proves the first identity
after summing over the three edges.  For the second use the exact
single-edge identity
\[
 D\,d\chi-2\chi\,dD
       =2\,dD+D^3(\beta\,d\alpha+\alpha\,d\beta).
\]
Multiplication by $m_km_l$ and summation gives the result.
\end{proof}

}{On the tie manifold $\{\chi=0\}$, these identities reduce to
$U_T=\sum_km_k\zeta_k\psi_k$ and
$dU_T=-\frac12\sum_km_k(\psi_k\,d\zeta_k+\phi_k\,d\eta_k)$.
Thus both the potential and its differential vanish on $O$.  Together with the rank
calculation in Lemma~\ref{lem:d1-mb}, this gives the quadratic estimate below.}

\begin{lemma}\label{lem:d1-estimate}
Every $y_0\in O$ has a neighbourhood $V\subset\C^7$ with a constant $c_V$ such that
$|U_T(y)|\le c_V\,|(\phi_3,\phi_4,\psi_3,\psi_4)(y)|^2$ for every $y\in V$ with $\chi(y)=0$.
\end{lemma}

\begin{proof}
{By Lemma~\ref{lem:d1-mb} and equivariance, five components $F_5$ of $F$
have independent differentials at $y_0$.  Their common zero set $F_5^{-1}(0)$ is therefore
a $2$-dimensional manifold near $y_0$.  The orbit map $\theta\mapsto\theta y_0$ immerses
a neighbourhood of $1\in\mathbb T$ onto a $2$-dimensional piece of $O$ contained in this
manifold.  Shrinking $V$, we obtain $F_5^{-1}(0)\cap V\subset O$.
The implicit function theorem then gives
$\operatorname{dist}(y,O)\le c\,|F_5(y)|\le c\,|(\phi_3,\phi_4,\psi_3,\psi_4)(y)|$
for $y\in V$ with $\chi(y)=0$.

Near $y_0$, the tie manifold $\{\chi=0\}$ is the graph of $D$ over
$(\zeta_3,\zeta_4,\eta_3,\eta_4)$, because
$\partial\chi_{kl}/\partial D_{kl}=2D_{kl}^{-1}(\chi_{kl}+1)\ne0$.
In these local coordinates, both $U_T$ and its differential vanish on $O$
(Lemma~\ref{lem:d1-virial}).  Taylor's formula at a nearest point of $O$ gives
$|U_T(y)|\le c'\operatorname{dist}(y,O)^2$.  Combining this with the preceding distance
estimate proves the assertion.}
\end{proof}

\begin{theorem}\label{thm:d1}
Let $m\in\Rp^5$.  No curve in $A_m$ has a Diagram~1 germ (Definition~\ref{def:germ}).  In
particular there is no Diagram~1 germ with an unbounded four-product, which is the case $r=1$ of
Theorem~\ref{thm:exclusion}\,(a).
\end{theorem}

\begin{proof}
{Suppose that $\gamma$ is a Diagram~1 germ.  Normalize it as above and
renumber $T$ as in Proposition~\ref{prop:d1-first}.  Along $\gamma$, the triangle ties give $\chi(y(t))=0$
identically, and $y(t)\to y_0\in O$.  We compare the order of the triangle potential obtained
from the force estimate with the order given by the virial identity.

Since $E^z_k=E^w_k=0$ along $\gamma$, Lemma~\ref{lem:d1-chart} gives
$\phi_k=t^{\mu/2}(z_k-X^z_k)$ and $\psi_k=t^{3\mu/2}(w_k-X^w_k)$, with the exponents listed
there and bounded coefficients.  By Proposition~\ref{prop:d1-cover}(iii),
$\frac32\mu+\rho\ge\frac12(\lambda+\mu)$.  Hence $|\phi(y(t))|=O(|t|^r)$ with
$r=\min(\frac32\mu,\frac12(\lambda+\mu))>0$, and $|\psi(y(t))|=O(|t|^{3\mu/2})$.
The quadratic estimate of Lemma~\ref{lem:d1-estimate} therefore gives
$|U_T(y(t))|=O(|t|^{2r})$.

On the other hand, Lemma~\ref{lem:d1-virial} and $\chi=0$ give
$U_T(y(t))=\sum_km_k\zeta_k\psi_k=t^{3\mu/2}\sum_km_k\zeta_k(w_k-X^w_k)$.
Because $w_k=W_T+\eta_k$ and $\sum_km_k\zeta_k=0$, we have
$\sum_km_k\zeta_kw_k=I_T(y(t))$.  Moreover,
$\ord X^w_k\ge\frac32\lambda$ by Proposition~\ref{prop:d1-cover}(ii).  Thus
$U_T(y(t))=t^{3\mu/2}\big(I_T(y_0)+o(1)\big)$ with $I_T(y_0)\ne0$.

Comparison with the quadratic estimate forces
$\frac32\mu\ge2r=\min(3\mu,\lambda+\mu)$.  Since $\mu>0$, it follows that
$\lambda+\mu\le\frac32\mu$, i.e.\ $\lambda\le\frac12\mu$, contradicting
Proposition~\ref{prop:d1-cover}(i).}
\end{proof}

\begin{remark}\label{rem:d1-wall}
{The strict inequality $\lambda>\mu/2$ enters only in the final
contradiction.  In the chart family \eqref{eq:d1-chart}, the outer bodies are $z$-circled
exactly when $\lambda\ge\mu/2$.  At $\lambda=\mu/2$, the chart becomes a Diagram~4 chart after
a rotation.  The orders $\frac32\mu$ and $\lambda+\mu$ then coincide, so the present
argument yields no contradiction; this case is treated in Section~\ref{sec:d45}.
For $\lambda<\mu/2$, the diagram is Diagram~5.  Since the proof requires no bound on
four-products, Theorem~\ref{thm:d1} excludes Diagram~1 germs both as first and as second
limits.}
\end{remark}



\section{Diagram 2}\label{sec:d2}

{Diagram~2 in \cite[Fig.~11]{AK} consists of the $z$-segment $12$ and the $w$-triangle $345$.
Every vertex carries a circle of the same colour as its edges \cite[\S8.2, p.~572]{AK}.
Since \eqref{eq:system} is invariant under exchanging $z$ and $w$ and renumbering the bodies,
we may assume that a Diagram~2 germ has this form.
Throughout this section, $A,A'$ denote the bodies in the pair, $\{A,A'\}=\{1,2\}$, and $B,B'$
denote bodies in the triangle $\{3,4,5\}$.  We write $M=m_1+m_2$, $M_T=m_3+m_4+m_5$ and
$m'=(m_3,m_4,m_5)$.  We reserve $t$ for a chart coordinate and use $s$ for the germ parameter
called $t$ in Lemma~\ref{lem:germ}; thus $\ord$ denotes order in $s$.
The notation $x,y,u$ agrees with \cite[\S8.2]{AK}.}

{The main difficulty is to rule out a germ approaching the boundary of the chart.
The blow-up of \eqref{eq:system} adapted to the diagram gives every Diagram~2 germ a limit
in the fibre $t=0$.  We call such boundary points \emph{labels}
(Proposition~\ref{prop:d2-cover}).  The torque identity for the pair imposes $\Phi=0$ at a
limit label and recovers the mass loci $L_j$, $K_j$, $P_j$ of \cite[\S8.2]{AK}
(Proposition~\ref{prop:d2-first}).  To determine whether a point on these loci can actually
be approached, we study the equations near its label.  When twenty selected equations
define a graph over $t$, the remaining two are equivalent to $tq(t)=0$.  Here $q(0)=\Phi$,
and $q'(0)$ is a quadratic form associated with the triangle
(Lemmas~\ref{lem:d2-isolation} and~\ref{lem:d2-second}).  Fold labels at which this graph
description fails require local reductions in one or two parameters.
The analysis in \S\ref{sec:d2-cases} excludes every label with $\Phi=0$ at positive masses;
in most regular cases, $q'(0)\ne0$ suffices.  The argument applies to complex configurations
without a reality condition.}

\subsection{\texorpdfstring{The blow-up chart and its boundary}{The blow-up chart and its boundary}}\label{sec:d2-chart}

{We use coordinates $(t,Z,W,D)\in\C^{21}$, where $Z=(Z_1,\dots,Z_5)$,
$W=(W_1,\dots,W_5)$ and $D=(D_{kl})_{k<l}$.  Their relation to the coordinates of
\eqref{eq:system} is}
\begin{equation}\label{eq:d2-chart}
\begin{gathered}
  z_A=t^{-1}Z_A,\quad w_A=tW_A,\quad z_B=tZ_B,\quad w_B=t^{-1}W_B,\\
  \delta_{12}=D_{12},\quad \delta_{BB'}=D_{BB'},\quad \delta_{AB}=tD_{AB}.
\end{gathered}
\end{equation}
{The chart coordinates $Z_k,W_k$ should be distinguished from the quantities
$Z_{kl},W_{kl}$ used to describe diagrams.  After multiplying the coordinate equations of
\eqref{eq:system} for $z_A,w_A,z_B,w_B$ by $t,t^{-1},t^{-1},t$, respectively, we obtain}
\begin{equation}\label{eq:d2-rows}
\begin{aligned}
  F^Z_A&=Z_A-m_{A'}D_{12}^3(Z_A-Z_{A'})-t^3\textstyle\sum_Bm_BD_{AB}^3(Z_A-t^2Z_B),\\
  F^W_A&=W_A-m_{A'}D_{12}^3(W_A-W_{A'})+t\textstyle\sum_Bm_BD_{AB}^3(W_B-t^2W_A),\\
  F^Z_B&=Z_B-\textstyle\sum_{B'\ne B}m_{B'}D_{BB'}^3(Z_B-Z_{B'})
         +t\textstyle\sum_Am_AD_{AB}^3(Z_A-t^2Z_B),\\
  F^W_B&=W_B-\textstyle\sum_{B'\ne B}m_{B'}D_{BB'}^3(W_B-W_{B'})
         -t^3\textstyle\sum_Am_AD_{AB}^3(W_B-t^2W_A),
\end{aligned}
\end{equation}
{The distance equations become $T_{12}=D_{12}^2(Z_2-Z_1)(W_2-W_1)-1$,
$T_{BB'}=D_{BB'}^2(Z_{B'}-Z_B)(W_{B'}-W_B)-1$ ($B<B'$) and
$T_{AB}=D_{AB}^2(t^2Z_B-Z_A)(W_B-t^2W_A)-1$.
We replace the gauge in \eqref{eq:system} by $G_1=Z_2-Z_1-1$ and $G_2=W_5-1$.
The common zero set of these $22$ polynomials is denoted by $\mathcal X_m\subset\C^{21}$,
and its boundary fibre by $\mathcal B_m=\mathcal X_m\cap\{t=0\}$.
A \emph{label} is a point of $\mathcal B_m$; it records the limiting pair, triangle and
crossed inverse distances in this chart.}

\begin{definition}[Shape factors]\label{def:d2-shapes}
For $\varsigma\in\{\pm1\}^3$ let $N_\varsigma(w_3,w_4,w_5)$ be the numerator of
$w_3/(\varsigma_3(w_5-w_4)^2)+w_4/(\varsigma_4(w_3-w_5)^2)+w_5/(\varsigma_5(w_4-w_3)^2)$.  The
quintics $S_1,S_2,S_3$ are $N_\varsigma$ for $\varsigma=(1,-1,-1),(1,-1,1),(1,1,-1)$, and
$N_{(1,1,1)}$ is a constant multiple of $S_4S_5$, where
\begin{gather*}
  S_4=w_3^3+w_4^3+w_5^3-w_3^2w_4-w_3w_4^2-w_3^2w_5-w_3w_5^2-w_4^2w_5-w_4w_5^2+3w_3w_4w_5,\\
  S_5=w_3^2+w_4^2+w_5^2-w_3w_4-w_3w_5-w_4w_5 .
\end{gather*}
The quartics $S_6,S_7,S_8$ are the factors of
$\operatorname{Res}_p\bigl((w_4-w_3)^3-p^3(w_5-w_4)^3,\ (w_4-w_3)^3-(-1-p)^3(w_3-w_5)^3\bigr)$
(which is $S_5^3S_6S_7S_8$ up to a constant) whose coefficients of $w_3^3w_4$ are $-2$, $1$ and
$-5$.  These are the factors of the complex three-body central configuration locus of
\cite[Appendix, pp.~585--586]{AK}.
\end{definition}

\begin{lemma}[computer-assisted]\label{lem:d2-fibre}
Let $m\in\Rp^5$.  A point $p=(0,Z,W,D)$ is a label if and only if
\begin{alphalist}
  \item $Z_1=-m_2/M$, $Z_2=m_1/M$, $D_{12}^3=1/M$, $W_1=-m_2D_{12}$ and $W_2=m_1D_{12}$;
  \item $F^Z_B=F^W_B=T_{BB'}=0$ at $t=0$ (system \eqref{eq:system} for the bodies $3,4,5$,
  without its gauge) and $W_5=1$;
  \item $D_{AB}^2Z_AW_B=-1$ for all $A,B$.
\end{alphalist}
At a label $\sum_Bm_BZ_B=\sum_Bm_BW_B=0$, every $W_B\ne0$, the $W_B$ are pairwise distinct, and
the \emph{triangle} $(Z_B,W_B,D_{BB'})$ is of one of two kinds.
\begin{romanlist}
  \item \emph{Lagrange:} $D_{BB'}=M_T^{-1/3}\omega^{a_{BB'}}$, $\omega=e^{2\pi i/3}$,
  $a_{BB'}\in\{0,1,2\}$ ($27$ classes), and $\kappa=(W_4-W_3)/(W_5-W_4)$ is a root of
  $\omega^{a_{45}}\kappa^2+(\omega^{a_{34}}+\omega^{a_{45}}-\omega^{a_{35}})\kappa+\omega^{a_{34}}$.
  The root is double exactly for the six classes with $a_{34},a_{35},a_{45}$ distinct; these are
  the \emph{fold} triangles.  $(W_3,W_4,W_5)$ is a zero of $S_5$ if $a_{34}=a_{35}=a_{45}$ or the
  $a_{BB'}$ are distinct, of $S_6$ if $a_{34}=a_{35}\ne a_{45}$, of $S_7$ if
  $a_{34}=a_{45}\ne a_{35}$, and of $S_8$ if $a_{35}=a_{45}\ne a_{34}$.
  \item \emph{Euler:} $Z_B=c^{-2}W_B$ for some $c\in\C^*$, $D_{BB'}=e_{BB'}c/(W_{B'}-W_B)$ for
  $B<B'$ with $e_{BB'}=\pm1$, and $(W_3,W_4,W_5)$ is a zero of $S_1$, $S_2$, $S_3$ or $S_4S_5$
  according to the class of $(e_{34},e_{35},e_{45})$ modulo a global sign.  An Euler triangle on
  $S_5$ is a Lagrange fold triangle.
\end{romanlist}
\end{lemma}

\begin{proof}
{At $t=0$, the crossed terms in \eqref{eq:d2-rows} vanish and
$T_{AB}=-D_{AB}^2Z_AW_B-1$, giving (b) and (c).  Together with $G_1$, the remaining pair
equations are $Z_A=m_{A'}D_{12}^3(Z_A-Z_{A'})$,
$W_A=m_{A'}D_{12}^3(W_A-W_{A'})$ and $T_{12}=0$, whose solutions are those in (a).
Summing $m_BF^Z_B$ and $m_BF^W_B$ gives the two centre-of-mass relations.
The distance equations also ensure that the triangle positions are nonzero and distinct:
if $W_B=0$, then $T_{1B}=-1$, and if $W_B=W_{B'}$, then $T_{BB'}=-1$.}

{Suppose first that $(1)_B,(Z_B)_B,(W_B)_B$ are linearly independent.
We claim that $D_{BB'}^3=1/M_T$ for all three edges.  The linear map
$X\mapsto(\sum_{B'\ne B}m_{B'}D_{BB'}^3(X_B-X_{B'}))_B$ on $\C^3$ fixes $(Z_B)$ and $(W_B)$
by (b), and annihilates $(1)_B$.  By the centre-of-mass relations,
$X\mapsto X-M_T^{-1}(\sum_Bm_BX_B)(1)_B$ has the same action on this basis.
The two maps therefore agree.  Comparing their entries at $(B,B')$, $B\ne B'$, gives
$m_{B'}D_{BB'}^3=m_{B'}/M_T$.
The coordinate equations now reduce to the centre-of-mass relations, and the distance
equations give $Z_{B'}-Z_B=M_T^{2/3}\omega^{a_{BB'}}/(W_{B'}-W_B)$.
Substituting these differences into $Z_5-Z_3=(Z_4-Z_3)+(Z_5-Z_4)$ gives the quadratic for
$\kappa$.}

{If the three vectors are dependent, then $Z=\alpha W+\beta$.
The centre-of-mass relation gives $\beta=0$, and the distance equations give $\alpha\ne0$.
Writing $\alpha=c^{-2}$ yields the asserted formula for $D_{BB'}$.
To identify the double roots, write $c_{BB'}=\omega^{a_{BB'}}$.
When two of the $c_{BB'}$ agree, the discriminant is $b(b-4a)\ne0$, with $|a|=|b|=1$,
where $a$ is the repeated value and $b$ the remaining one.
When all three are distinct, $c_{34}+c_{45}=-c_{35}$ and $c_{34}c_{45}=c_{35}^2$, so the
discriminant vanishes.  The identification of the corresponding shape factors is
computer-assisted.}
\end{proof}

{This identification uses exact arithmetic over $\mathbb Q(\omega)$.
For each of the $27$ Lagrange classes, the corresponding factor $S_j$, expressed as a
polynomial in $\kappa$, is reduced modulo the quadratic for $\kappa$; for fold classes it is
evaluated at the double root.  For each of the eight sign vectors $(e_{34},e_{35},e_{45})$,
the corresponding $N_\varsigma$ is identified with one of $\pm S_1$, $\pm S_2$, $\pm S_3$,
$\pm S_4S_5$.  Finally, for Euler triangles on $S_5$, the three quantities $D_{BB'}^3$ are
checked to coincide modulo $S_5$ (Appendix~\ref{app:verification}).}

\begin{definition}[Mass polynomials]\label{def:d2-mass}
For positions $W_3,W_4,W_5$ put $f_B=m_B^2/W_B$,
$\Lambda=f_3^2+f_4^2+f_5^2-2f_3f_4-2f_3f_5-2f_4f_5$ and
$\Theta=3(f_3^2+f_4^2+f_5^2)+2f_3f_5+2f_4f_5-6f_3f_4$.  Equations (25) and (28) of
\cite[\S8.2]{AK} are
\begin{equation}\label{eq:d2-25-28}
  \Lambda=0,\qquad
  \Lambda\,y^2+2(f_3-f_4+f_5)(f_3-f_4-f_5)\,y+\tfrac12(\Theta-\Lambda)=0 .
\end{equation}
For a rational expression $H$ in $W_3,W_4,W_5$ whose denominator is a monomial, let $H^\flat$ be
the polynomial in $W_4$ obtained by clearing the position denominators, substituting $W_5=1$ and
$W_3=-(m_4W_4+m_5)/m_3$ (centre of mass), and clearing the resulting power of $m_3$.  Define
$L_j(m')$ and $P_j(y)$ by
\[
  \operatorname{Res}_{W_4}(S_j^\flat,\Lambda^\flat)=c_j\,L_j,\qquad
  \operatorname{Res}_{W_4}\bigl(S_j^\flat,(28)^\flat\bigr)=m_3^{2\deg S_j}P_j(y),
\]
where $c_j$ has no zero on $\Rp^3$; for $j=5$ the first resultant is
$m_3^4M_T^2L_5L_5'$.  Write $P_j(y)=E_j(y^2)+yO_j(y^2)$ and $K_j=P_j(0)=E_j(0)$.
\end{definition}

{These definitions recover the factors $L_1,\dots,L_8,L_5'$ and the
polynomials $P_1,\dots,P_8$ in \cite[pp.~573--575]{AK} from the shape factors in its appendix.
They agree with the coefficients and symmetries displayed in \cite{AK}.
In particular, $2L_5=(m_3-m_4)^2+(m_4-m_5)^2+(m_5-m_3)^2$, while $L_5'$ has degree $12$
and $91$ positive coefficients.
Permuting $(m_3,m_4,m_5)$ permutes $\{L_1,L_2,L_3\}$ and $\{L_6,L_7,L_8\}$ up to nonzero
constants.  Exchanging bodies $3$ and $4$ sends $S_1$ to $-S_2$, exchanges
$P_1\leftrightarrow P_2$ and $P_6\leftrightarrow P_7$, and fixes $P_3,P_4,P_5,P_8$.}

{The verification is exact over $\mathbb Q$, treating the masses and
$y$ as indeterminates.  The resultants defining $L_j$ and $L_5'$ are computed and factored.
Those defining $P_j$ are determined by evaluations on unisolvent grids modulo three primes,
with a priori bounds on the degrees and coefficients.  Expansion and substitution then
verify the decompositions $P_j=E_j+yO_j$, the leading coefficients and the symmetries
(Appendix~\ref{app:verification}).}
\subsection{\texorpdfstring{Limits of Diagram~2 germs in the chart}{Limits of Diagram 2 germs in the chart}}\label{sec:d2-cover}

\begin{proposition}[Existence of a boundary limit]\label{prop:d2-cover}
Let $m\in\Rp^5$ and let $\gamma$ be a Diagram~2 germ on a curve in $A_m$.  Along $\gamma$ put
$t=(z_{12}w_5)^{-1/2}$ \textup{(}either branch\textup{)} and $\lambda=(tz_{12})^{-1}$, and let
$(Z,W,D)$ be given by \eqref{eq:d2-chart} applied to the rotated configuration
$(\lambda z,\lambda^{-1}w,\delta)$.  Then $t,Z,W,D$ are Puiseux series in $s$ with $\ord t>0$ and
nonnegative orders, $(t,Z,W,D)\in\mathcal X_m$ for $s\ne0$, and the limit $p$ as $s\to0$ is a label at
which $Z_1,Z_2$, $W_3,W_4,W_5$ and all $D_{kl}$ are nonzero.
\end{proposition}

\begin{proof}
{The equations in \eqref{eq:system}, apart from the gauge, are invariant
under $(z,w)\mapsto(\lambda z,\lambda^{-1}w)$.  The chart equations are obtained from them by
multiplying by powers of $t$, and the chosen normalization gives $G_1=\lambda tz_{12}-1=0$
and $G_2=\lambda^{-1}tw_5-1=t^2z_{12}w_5-1=0$.
Thus the chart point belongs to $\mathcal X_m$.

To prove convergence, we estimate the order of each coordinate.  A rotation of $\gamma$
before this construction changes neither $t$ nor the resulting chart point.
We may therefore use the normalization of Lemma~\ref{lem:germ} and \cite[p.~547]{AK}, in
which the maximal components of $\mathcal Z$ and $\mathcal W$ have the same order.
We measure orders in units for which this common order is $-2$.
Fact~\ref{fact:rules} then gives the following estimates.}
\begin{enumerate}
  \item The pair is an isolated component of the $z$-diagram with $z$-circles, so its $z$-stroke
  is maximal (Rule~1e): $\ord z_{12}=-2$; by Estimate~2, $r_{12}\approx1$, so $\ord w_{12}=2$ and
  $\ord\delta_{12}=0$.
  \item The triangle is an isolated $w$-component with $w$-circles; it has a maximal $w$-stroke
  (Rule~1e), and by Rule~2e all three are maximal: $\ord w_{BB'}=-2$, $\ord z_{BB'}=2$,
  $\ord\delta_{BB'}=0$.
  \item The circles give $\ord z_A=\ord w_B=-2$.  The bodies $3,4,5$ carry no $z$-stroke and
  $1,2$ no $w$-stroke, hence no circle of that colour (Rule~1a): $\ord z_B>-2$, $\ord w_A>-2$.
  \item Hence $\ord z_{AB}=\ord w_{AB}=-2$ and $\ord\delta_{AB}=2$.
  \item In $w_A=\sum_{l\ne A}m_l\delta_{Al}^3(w_A-w_l)$ the pair term has order $2$ and the
  crossed terms order $4$, so $\ord w_A\ge2$; likewise $\ord z_B\ge2$.
  \item So $\ord t=2$ and $\ord\lambda=0$; $Z_A=\lambda tz_A$, $W_B=\lambda^{-1}tw_B$, $D_{12}$,
  $D_{BB'}$ and $D_{AB}=\delta_{AB}/t$ have order exactly $0$, and $W_A=\lambda^{-1}w_A/t$,
  $Z_B=\lambda z_B/t$ have order $\ge0$.
\end{enumerate}
{All chart coordinates consequently converge, with $t\to0$.
Those of order exactly $0$ have nonzero limits.
Since $\mathcal X_m$ is closed, the limit belongs to it and satisfies $t=0$.}
\end{proof}

\subsection{\texorpdfstring{The torque identity and the mass conditions}{The torque identity and the mass conditions}}\label{sec:d2-first}

\begin{lemma}[Syzygies]\label{lem:d2-syzygy}
Identically in $(t,Z,W,D)$,
\begin{gather*}
  \textstyle\sum_{k=1}^5m_k(W_kF^Z_k-Z_kF^W_k)=0,\qquad
  \sum_Am_A(W_AF^Z_A-Z_AF^W_A)=t\,Q,\\
  Q=\textstyle\sum_{A,B}m_Am_BD_{AB}^3(t^4W_AZ_B-Z_AW_B).
\end{gather*}
At $t=0$ the only rows with a nonzero $t$-derivative are $F^W_A$ and $F^Z_B$, with
$\partial_tF^W_A=\sum_Bm_BD_{AB}^3W_B$ and $\partial_tF^Z_B=\sum_Am_AD_{AB}^3Z_A$.  At a label,
$Q=\Phi:=\sum_Bm_B\varphi_B$ with $\varphi_B=m_1D_{1B}+m_2D_{2B}$, and
$\partial_tF^Z_B=-\varphi_B/W_B$.
\end{lemma}

\begin{proof}
Since $W_AF^Z_A=w_AF^z_A$, $Z_AF^W_A=z_AF^w_A$, $W_BF^Z_B=w_BF^z_B$ and $Z_BF^W_B=z_BF^w_B$ for
the coordinate equations $F^z_k,F^w_k$ of \eqref{eq:system}, the first identity is the torque
identity $\sum_km_k(w_kF^z_k-z_kF^w_k)=0$, which holds by antisymmetry.  In the second sum the
terms with $D_{12}$ cancel by antisymmetry and the crossed terms give $tQ$.  The derivatives are
read off from \eqref{eq:d2-rows}, and Lemma~\ref{lem:d2-fibre}(c) gives
$D_{AB}^3Z_AW_B=-D_{AB}$ and $D_{AB}^3Z_A=-D_{AB}/W_B$.
\end{proof}

{To express the torque condition in terms of the triangle, we first describe
the crossed inverse distances and their signs.}  At a label fix $v_B$ with $v_B^2=1/W_B$.  By Lemma~\ref{lem:d2-fibre}(a),(c),
$D_{1B}^2=M/(m_2W_B)$ and $D_{2B}^2=-M/(m_1W_B)$, so $D_{1B}=s_B\sqrt{M/m_2}\,v_B$ and
$D_{2B}=is'_B\sqrt{M/m_1}\,v_B$ with $s_B,s'_B\in\{\pm1\}$ (positive square roots); these are the
$64$ crossed-sign assignments.  Put $\sigma_B=s_Bs'_B$,
$c_\pm=\sqrt M\,(m_1/\sqrt{m_2}\pm i\,m_2/\sqrt{m_1})$ and $x=i(m_2/m_1)^{3/2}$.  Then
$c_+c_-=M(m_1^3+m_2^3)/(m_1m_2)>0$, $\varphi_B=c_{\sigma_B}s_Bv_B$, and
$m_2D_{2B}/(m_1D_{1B})=\sigma_Bx$ with $x^2=-m_2^3/m_1^3$.  The label is \emph{same-sign} if
$\sigma_3=\sigma_4=\sigma_5$ and \emph{one-different} otherwise (two and six of the eight
patterns); then the \emph{odd body} is the $B$ whose $\sigma_B$ differs from the majority sign
$\sigma$.  In both cases put $\tilde v_B=\varphi_B/c_\sigma$, so that $\tilde v_B=s_Bv_B$ if
$\sigma_B=\sigma$ and $\tilde v_k=\rho\,s_kv_k$ for the odd body $k$, where
$\rho=c_{-\sigma}/c_\sigma$ satisfies $|\rho|=1$, $\rho\ne\pm1$.  Thus
$\Phi=c_\sigma\sum_Bm_B\tilde v_B$.
\begin{proposition}[Necessary mass conditions]\label{prop:d2-first}
Let $m\in\Rp^5$ and let $p$ be the limit label of a Diagram~2 germ
\textup{(}Proposition~\ref{prop:d2-cover}\textup{)}.  Then $\Phi(p)=0$, and:
\begin{alphalist}
  \item if $p$ is same-sign, its triangle satisfies \textup{(25)}, $\Lambda=0$, and
  $L_j(m')=0$ for every $j\ne5$ with $S_j(W_3,W_4,W_5)=0$ \textup{(}$L_5L_5'=0$ if $j=5$\textup{)};
  \item if $p$ is one-different with odd body $5$ and majority sign $\sigma$, its triangle satisfies
  \textup{(28)} with $y=\sigma(x+x^{-1})/2$, and $u:=y^2=-(T-1)^2/(4T)$, $T=(m_1/m_2)^3$.  Hence
  $P_j(y)=0$ for every $j$ with $S_j(W_3,W_4,W_5)=0$: if $m_1=m_2$ then $K_j(m')=0$, and if
  $m_1\ne m_2$ then $u<0$ and $E_j(u)=O_j(u)=0$.
\end{alphalist}
\end{proposition}

\begin{proof}
{All equations vanish along the germ, and $t\not\equiv0$.
Lemma~\ref{lem:d2-syzygy} therefore gives $Q\equiv0$ along it, and continuity implies
$\Phi(p)=Q(p)=0$.  This argument requires no rank condition on the Jacobian and applies
also at fold labels.}

{In case (a), $\Phi=c_\sigma\sum_Bm_Bs_Bv_B$ with $c_\sigma\ne0$.
The product of $m_3v_3\pm m_4v_4\pm m_5v_5$ over the four sign choices equals $\Lambda$,
because $(m_Bv_B)^2=f_B$.  Hence the vanishing of the torque gives (25).
In case (b), put $r_B=m_1m_BD_{1B}$, so that $r_B^2=(m_1^2M/m_2)f_B$, and $x'=\sigma x$.
The torque becomes $\Phi=(r_3+r_4)(1+x')+r_5(1-x')$, as in (26) of \cite[\S8.2]{AK}.
Its product over the four sign changes of $r_3,r_4$ is $4x'^2(m_1^2M/m_2)^2$ times the
left-hand side of (28), with $y=(x'+1/x')/2$.}

{Thus $W_4$ is a common root of $S_j^\flat$ and $\Lambda^\flat$ in
case (a), or of the shape polynomial and $(28)^\flat$ in case (b).
The resultants in Definition~\ref{def:d2-mass} vanish, and their nonvanishing prefactors
$c_j\ne0$, $m_3\ne0$ give the asserted mass conditions.
Finally, $y=iv$ with $v=\sigma\bigl((m_2/m_1)^{3/2}-(m_1/m_2)^{3/2}\bigr)/2\in\R$,
which vanishes if and only if $m_1=m_2$.
Since $E_j,O_j$ have real coefficients, separating the real and imaginary parts of
$P_j(iv)=E_j(-v^2)+ivO_j(-v^2)=0$ proves the last assertion.}
\end{proof}

\subsection{\texorpdfstring{Local exclusion of boundary limits}{Local exclusion of boundary limits}}\label{sec:d2-isolation}

Let $F_{\mathrm{sel}}$ denote the twenty rows other than $F^Z_2$ and $F^Z_5$.  A label $p$ is
\emph{excluded} if no holomorphic curve germ $c\colon(\C,0)\to(\mathcal X_m,p)$ has
$t\circ c\not\equiv0$; an isolated point of $\mathcal X_m$ is excluded.

\begin{lemma}[Local analytic lemma]\label{lem:d2-isolation}
Let $m\in\Rp^5$ and $p\in\mathcal B_m$.
\begin{alphalist}
  \item If $p$ is excluded, it is not the limit label of a Diagram~2 germ.
  \item Suppose $dF_{\mathrm{sel}}(p)$ has rank $20$, so that near $p$ the set
  $\{F_{\mathrm{sel}}=0\}$ is a smooth curve $C$.  Then near $p$ the set $\mathcal X_m$ consists of the
  points of $C$ at which $tQ=0$, and $p$ is excluded if $Q|_C\not\equiv0$.  If moreover
  $\partial F_{\mathrm{sel}}/\partial(Z,W,D)(p)$ is invertible, then $C$ is the graph of an
  analytic map $t\mapsto\phi(t)$, and $p$ is an isolated point of $\mathcal X_m$ as soon as
  $q(t)=Q(t,\phi(t))\not\equiv0$, in particular if $q(0)=\Phi(p)\ne0$ or $q'(0)\ne0$.
  \item Suppose some twenty of the twenty-two polynomials have a differential of rank $20$ at
  $p$, so that near $p$ their common zeros form a smooth curve $C'$.  If one of the other two
  polynomials does not vanish identically on $C'$, then $p$ is an isolated point of $\mathcal X_m$.
\end{alphalist}
\end{lemma}

\begin{proof}
(a) By Proposition~\ref{prop:d2-cover}, a Diagram~2 germ with limit $p$ gives, after the
substitution $s=\tau^N$ that clears the denominators of the Puiseux exponents, a holomorphic germ
$\tau\mapsto(t,Z,W,D)$ in $\mathcal X_m$ at $p$ with $t\not\equiv0$.
(b) On $C$ all rows except $F^Z_2,F^Z_5$ vanish and $W_5=1$, so Lemma~\ref{lem:d2-syzygy} gives
$m_2W_2F^Z_2=tQ=-m_5F^Z_5$ on $C$; and $W_2(p)=m_1D_{12}(p)\ne0$.  This describes $\mathcal X_m$
near $p$.  A germ $c$ in $\mathcal X_m$ at $p$ lies in $C$, and if $t\circ c\not\equiv0$ then
$Q\circ c\equiv0$; when
$Q|_C\not\equiv0$ its zeros on the curve $C$ are isolated, so $c$ is constant, which contradicts
$t\circ c\not\equiv0$.  In the graph case $\mathcal X_m$ is, near $p$, the set of points $(t,\phi(t))$ with
$tq(t)=0$, which is $\{p\}$ when $q\not\equiv0$.
(c) Near $p$, $\mathcal X_m$ lies in $C'$ and in the zero set of the other polynomial, whose restriction
to the smooth curve $C'$ is analytic and not identically zero, hence has an isolated zero at $p$.
\end{proof}

{At the labels where the Jacobian of the twenty-two equations has rank
$19$, the implicit function theorem leaves two local parameters.  The following lemma
reduces exclusion to the Taylor coefficients of the three remaining equations.
Its proof follows the possible tangent directions of a formal arc by a Newton--Puiseux
argument after Lyapunov--Schmidt reduction.}

\begin{lemma}[Corank-two isolation]\label{lem:d2-corank-two}
Let $m\in\Rp^5$ and $p\in\mathcal B_m$.  Let $F_R$ be nineteen of the twenty-two chart polynomials,
$F_{R'}$ the other three, and $y\ne t$ a chart coordinate such that the derivative of
$(F_R,t,y)\colon\C^{21}\to\C^{21}$ at $p$ is invertible.  By the implicit function theorem,
near $p$ the set $\{F_R=0\}$ is the image of an injective analytic map $\phi$, defined near
$0\in\C^2$, with $\phi(0)=p$ and $(y,t)\circ\phi(a,t)=(y(p)+a,t)$, and near $p$
\[
  \mathcal X_m=\phi\bigl(g^{-1}(0)\bigr),\qquad g=F_{R'}\circ\phi\colon(\C^2,0)\to(\C^3,0).
\]
Let $g^{(j)}$ be the homogeneous part of degree $j$ of $g$, and for $\xi\in\C^2$ let
$\Pi_\xi=Dg^{(2)}(\xi)$, a $3\times2$ matrix.  Suppose:
\begin{alphalist}
  \item $g^{(1)}=0$; equivalently, the Jacobian at $p$ of all twenty-two polynomials has
  rank~$19$;
  \item the components of $g^{(2)}$ have only finitely many common zeros $[\xi_1],\dots,[\xi_r]$
  in $\mathbb P^1$, and $\Pi_{\xi_j}\ne0$ for every $j$;
  \item for every $j$ there is $N\ge3$ such that the vector $r_N(\xi_j)$ defined below exists
  and does not lie in $\Pi_{\xi_j}\C^2$.
\end{alphalist}
Then no nonconstant formal arc $x\in\C[[\tau]]^2$ with $x(0)=0$ satisfies $g(x)=0$.  Hence $p$
is an isolated point of $\mathcal X_m$; in particular $p$ is excluded.  The same holds if
\textup{(a)--(c)} hold for $Kg(L\,\cdot)$, $K\in GL_3(\C)$, $L\in GL_2(\C)$, in place
of $g$.

The vectors $r_N(\xi)$ are defined for a common zero $\xi$ of $g^{(2)}$ and a fixed
$e\in\C^2\setminus\C\xi$.  As $\Pi_\xi\xi=2g^{(2)}(\xi)=0$ and $\Pi_\xi\ne0$, the vector
$\varpi=\Pi_\xi e$ spans the line $\Pi_\xi\C^2$.  Put $\beta_1=0$.  For $n\ge1$, if $\beta_n$ is
defined, let $r_{n+2}(\xi)$ be the coefficient of $\lambda^{n+2}$ in
$g(\lambda\xi+\beta_n(\lambda)e)$; if moreover $r_{n+2}(\xi)=-b_{n+1}\varpi$ with $b_{n+1}\in\C$,
put $\beta_{n+1}(\lambda)=\beta_n(\lambda)+b_{n+1}\lambda^{n+1}$.  Thus
\[
  r_3(\xi)=g^{(3)}(\xi),\qquad
  r_4(\xi)=g^{(4)}(\xi)+b_2\,Dg^{(3)}(\xi)e+b_2^2\,g^{(2)}(e).
\]
Condition (c) with $N=3$ says that the $3\times3$ matrix $\bigl(\Pi_{\xi_j}\,|\,r_3(\xi_j)\bigr)$
has rank two; with $N=4$ it says that $\bigl(\Pi_{\xi_j}\,|\,r_3(\xi_j)\bigr)$ has rank one and
$\bigl(\Pi_{\xi_j}\,|\,r_4(\xi_j)\bigr)$ has rank two.
\end{lemma}

\begin{proof}
The points of $\{F_R=0\}$ near $p$ are the points $\phi(a,t)$, and $\phi(a,t)\in\mathcal X_m$ if
and only if $g(a,t)=0$.  If $p$ were not isolated in $\mathcal X_m$, the curve selection lemma
\cite[Thm.~2.5.5]{BCR}, applied to the real algebraic set $\mathcal X_m\subset\C^{21}=\R^{42}$,
would give a nonconstant semialgebraic arc $\gamma\colon[0,\varepsilon)\to\mathcal X_m$ with
$\gamma(0)=p$.  Its
coordinates are convergent Puiseux series, so $x=(y\circ\gamma-y(p),t\circ\gamma)$ is, after a
substitution $s=\tau^{N'}$, a nonconstant arc in $\C[[\tau]]^2$ with $g(x)=0$.  So the first
assertion implies the second.  The zero set of $Kg(L\,\cdot)$ is $L^{-1}g^{-1}(0)$, which gives
the last assertion.

Let $x$ be a nonconstant formal arc with $x(0)=0$ and $g(x)=0$, and write
$x=\tau^k\xi'+O(\tau^{k+1})$ with $\xi'\ne0$.  By (a),
$g(x)=\tau^{2k}g^{(2)}(\xi')+O(\tau^{2k+1})$, so $\xi'\in\C^*\xi$ for one of the $\xi=\xi_j$
of (b).  Write $x=u\xi+\tilde\theta e$; then $\ord u=k<\ord\tilde\theta$.  Take $\lambda=u(\tau)$ as
the new parameter; then $\tau$ is a power series in $\lambda^{1/k}$ with nonzero linear term,
and $x=\lambda\xi+\theta e$ with $\theta\in\C[[\lambda^{1/k}]]$ and $\ord\theta>1$ ($\theta=0$ is
allowed).

We show by induction on $n\ge1$ that $\beta_n$ is defined, that
$g(\lambda\xi+\beta_ne)=O(\lambda^{n+2})$ and that $o_n=\ord(\theta-\beta_n)>n$.  For $n=N-2$
this contradicts (c).  For $n=1$ the claim holds by (a), $g^{(2)}(\xi)=0$ and $\ord\theta>1$.
Taylor's formula in the direction $e$ gives
\[
  0=g(\lambda\xi+\theta e)
   =g(\lambda\xi+\beta_ne)+(\theta-\beta_n)\bigl(\lambda\varpi+O(\lambda^2)\bigr)
    +O(\lambda^{2o_n}),
\]
because $Dg(\lambda\xi+\beta_ne)e=\lambda\Pi_\xi e+\beta_n\Pi_ee+O(\lambda^2)
=\lambda\varpi+O(\lambda^2)$ by (a) and $\ord\beta_n\ge2$.  If $o_n<n+1$, the term
$(\theta-\beta_n)\lambda\varpi$ has order $o_n+1<n+2$, every other term has larger order
($2o_n>o_n+1$), and its leading coefficient is a nonzero multiple of $\varpi\ne0$; this is
impossible.  Hence $o_n\ge n+1$, and the coefficient of $\lambda^{n+2}$ gives
$r_{n+2}(\xi)+b\varpi=0$, where $b$ is the coefficient of $\lambda^{n+1}$ in $\theta$.  So
$\beta_{n+1}$ is defined, $b_{n+1}=b$ since $\varpi\ne0$, and $o_{n+1}>n+1$.  The same expansion
with $\beta_{n+1}$ in place of $\theta$ gives
$g(\lambda\xi+\beta_{n+1}e)=(r_{n+2}(\xi)+b_{n+1}\varpi)\lambda^{n+2}+O(\lambda^{n+3})
=O(\lambda^{n+3})$.
\end{proof}

\begin{remark}\label{rem:d2-corank-two}
The $\xi_j$ are the points where the strict transform of $g^{-1}(0)$ could meet the exceptional
curve of the blow-up of $0\in\C^2$, and the induction is the Newton--Puiseux algorithm in the
blow-up chart at $\xi_j$.  Arcs inside the boundary fibre $\{t=0\}$ need no separate argument:
they are tangent to the line $t=0$, which is then one of the $\xi_j$, and the same induction
excludes them.
\end{remark}

{We next compute the first derivative of the reduced torque equation
using the triangle block and the quadratic form $Q_T$.
For a label $p$, let $J_T(p)$ be the
Jacobian at $p$ of the nine equations
$F^Z_3,F^W_3,F^Z_4,F^W_4,F^W_5,T_{34},T_{35},T_{45},G_2$ with respect to the nine triangle
coordinates $Z_B,W_B,D_{BB'}$.  If $J_T(p)$ is invertible and $\xi\in\C^3$, let
$(\zeta,\eta,\vartheta)$ be the solution of $J_T(p)(\zeta,\eta,\vartheta)=r(\xi)$.
Here $r(\xi)$ has entries $\xi_3/W_3$ and $\xi_4/W_4$ in the rows $F^Z_3$, $F^Z_4$ and $0$
elsewhere.  Define}
\[
  Q_T(\xi)=-\tfrac12\textstyle\sum_Bm_B\,\xi_B\,\eta_B/W_B .
\]
{This quadratic form $Q_T$ depends only on $m'$ and the triangle of $p$.}

\begin{lemma}[Second-order coefficient]\label{lem:d2-second}
Let $p$ be a label with $\Phi(p)=0$ and $J_T(p)$ invertible.  Then
$\partial F_{\mathrm{sel}}/\partial(Z,W,D)(p)$ is invertible and, in the notation of
Lemma~\ref{lem:d2-isolation} and of the crossed data,
\[
  q(0)=0,\qquad q'(0)=c_\sigma^2\,Q_T(\tilde v).
\]
\end{lemma}

\begin{proof}
{At $t=0$ the pair rows involve only $Z_A,W_A,D_{12}$, the triangle rows only the triangle
coordinates, and $T_{AB}$ involves no crossed coordinate other than $D_{AB}$, with
$\partial T_{AB}/\partial D_{AB}=2/D_{AB}$; hence the Jacobian is block lower triangular.
To check its pair block directly, write
\[
 d=D_{12},\qquad u=Z_2-Z_1,\qquad h=W_2-W_1.
\]
At the label $u=1$, $Md^3=1$ and $h=Md$. The five selected pair rows are
$F^Z_1,F^W_1,F^W_2,T_{12},G_1$, and at $t=0$ they satisfy
\[
\begin{gathered}
 F^Z_1=Z_1+m_2d^3u,\qquad F^W_2-F^W_1=(1-Md^3)h,\\
 m_1F^W_1+m_2F^W_2=m_1W_1+m_2W_2,\qquad T_{12}=d^2uh-1.
\end{gathered}
\]
A vector in their differential kernel first has $\dot u=0$ from the gauge, then
$\dot d=0$ from the difference of the two $W$-rows, whose $\dot d$ coefficient is
$-3Md^2h\ne0$. The tie gives $\dot h=0$, the weighted $W$-sum gives
$\dot W_1=\dot W_2=0$, and the $Z$-row and gauge give
$\dot Z_1=\dot Z_2=0$. Thus the pair block is invertible.
Differentiate $F_{\mathrm{sel}}(t,\phi(t))=0$ at $t=0$.
By Lemma~\ref{lem:d2-syzygy}, the pair $Z$-row, tie and gauge have no explicit
$t$-forcing, whereas
\[
 \partial_tF^W_A=-Z_A^{-1}\sum_Bm_BD_{AB},\qquad
 \partial_tF^W_1-\partial_tF^W_2=\frac{M\Phi}{m_1m_2}=0.
\]
The same difference, tie and $Z$-row calculation therefore gives
$D_{12}'(0)=Z_1'(0)=Z_2'(0)=0$.

}The triangle rows give
$J_T(p)(Z',W',D')_T(0)=r(\varphi)$, hence $W_B'(0)=c_\sigma\eta_B(\tilde v)$.  On $C$ the ties
$T_{AB}=0$ give $Q=\sum_{A,B}m_Am_BD_{AB}+O(t^2)$ and
$D_{AB}'=-\frac12D_{AB}(Z_A'/Z_A+W_B'/W_B)$.  Therefore
$q'(0)=\sum_{A,B}m_Am_BD'_{AB}(0)=-\frac12\sum_Bm_B\varphi_BW_B'(0)/W_B=c_\sigma^2Q_T(\tilde v)$.
\end{proof}

\begin{definition}\label{def:d2-psi}
For a label $p$ with $J_T(p)$ invertible put $\Psi_T(p)=Q_T(\tilde v)$ if $p$ is same-sign and
$\Psi_1(p)=Q_T(\tilde v)$ if $p$ is one-different.  The label is \emph{isolated at second order}
if $\Phi(p)=0$, $J_T(p)$ is invertible and $Q_T(\tilde v)\ne0$; by
Lemmas~\ref{lem:d2-isolation} and~\ref{lem:d2-second} it is then an isolated point of $\mathcal X_m$, and
the omitted rows $F^Z_2,F^Z_5$ vanish to order exactly two along the curve $C$.
\end{definition}

At a label with $\Phi(p)=0$ the vector $\tilde v$ is determined by the triangle and the signs
$s_B$ (for an odd body $k$, $m_k\tilde v_k=-\sum_{B\ne k}m_B\tilde v_B$), so $\Psi_T$ and
$\Psi_1$ do not involve $m_1,m_2$.  For a one-different label with odd body $5$ the condition
$\Phi=0$ requires $|m_3\tilde v_3+m_4\tilde v_4|=m_5|v_5|$, i.e.\ the vanishing of the
\emph{sheet function} $G=|m_3\tilde v_3+m_4\tilde v_4|^2-m_5^2$ (recall $W_5=1$), which again
depends on the triangle only; then $\tilde v_5=\rho s_5v_5$ fixes $\rho$ up to sign, hence
$m_1/m_2$ up to exchanging $m_1$ and $m_2$.

\subsection{\texorpdfstring{Exclusion on the mass loci}{Exclusion on the mass loci}}\label{sec:d2-cases}

{It remains to exclude every label $p$ with $\Phi(p)=0$ at positive
masses, by Proposition~\ref{prop:d2-first} and Lemma~\ref{lem:d2-isolation}(a).
The torque condition divides the problem according to the relative signs of the crossed
inverse distances and, in the one-different case, according to whether the pair masses
are equal.  Table~\ref{tab:d2-cases} records the mass loci and the local argument used in
each case.  The required computer-assisted results are stated below.
The ball arithmetic used here is IEEE arithmetic with proved rounding bounds
(Appendix~\ref{app:verification}).}

\begin{table}[ht]
\centering
\footnotesize
\setlength{\tabcolsep}{3pt}
\begin{tabular}{p{2.5cm}p{2.5cm}p{3.3cm}p{3.4cm}}
\hline
Labels & Mass locus & Method & Computation \\
\hline
$\Phi\ne0$ (any) & all $m$ & {$Q(p)=\Phi(p)\ne0$, $tQ=0$} & exact identities \\
same-sign, $S_4$, or $S_5$ with $L_5'=0$ & $L_4=0$, $L_5'=0$ & no point in $\Rp^3$ & exact \\
same-sign, $S_5$ & $L_5=0$: $m_3=m_4=m_5$ & order three; corank-two blow-up & exact jets \\
same-sign, $S_1$--$S_3$, $S_6$--$S_8$ & rays of $L_1$--$L_3$, $L_6$--$L_8$ & second order,
{$\Psi_T\ne0$} & Bernstein subdivision, Krawczyk method; interval arithmetic \\
\hline
one-different, $m_1=m_2$, $S_3,S_6$--$S_8$ & $K_3,K_6,K_7,K_8=0$ & no point in $\Rp^3$ &
Bernstein coefficients \\
one-different, $m_1=m_2$, $S_1,S_2$ & $K_1=0$, $K_2=0$ & second order & exact; interval cover \\
one-different, $m_1=m_2$, $S_4$ & {$K_4=0$: cones over four arcs} & second order & exact \\
one-different, $m_1=m_2$, $S_5$ & $K_5=0$: $\Rp(1,1,\rho_*)$ & orders one to four &
exact; interval arithmetic \\
\hline
one-different, $m_1\ne m_2$, not fold & $E_j(u)=O_j(u)=0$ & second order, $\Psi_1\ne0$: ball
covers, exact lemmas & ball and interval arithmetic; exact \\
one-different, $m_1\ne m_2$, fold & $E_5(u)=O_5(u)=0$ & corank one; $Q|_C$ has a simple zero &
exact \\
\hline
\end{tabular}
\caption{The Diagram~2 cases.  A label with $\Phi=0$ lies on $S_j$ for the $j$ given in the
first column (Lemma~\ref{lem:d2-fibre}); one-different labels are taken with odd body~$5$.}
\label{tab:d2-cases}
\end{table}

\begin{lemma}[computer-assisted]\label{lem:d2-empty}
\textup{(a)} The polynomials $L_4$, $L_5'$, $K_3$, $K_6$, $K_7$ and $K_8$ have no zero in $\Rp^3$.
{More precisely, at positive masses every normalized $W$-shape on $S_4$ is real
\textup{(}the centre-of-mass section of $S_4$ is a cubic with positive discriminant\textup{)}
and violates \textup{(25)}.  Here real refers to the coordinates $W_B$, not to the condition
$W_B=\overline{Z_B}$ for the full complex configuration.}
\textup{(b)} If $m_3=m_4=m_5$, no $S_j$ with $j\ne5$ has a root satisfying \textup{(25)}.
\end{lemma}
{The verification uses exact arithmetic over $\mathbb Q$.
For $L_4$, a rational parametrization of $S_4$ gives the obstruction:
at positive masses, a real normalized $W$-shape has positions $W_B$ of different signs.
Then $\Lambda$ is a sum of nonnegative terms, at least one of which is positive.
The coefficients of $L_5'$ are positive.  For $K_3$, $K_6$, $K_7$, $K_8$, nonnegative
Bernstein coefficients are checked on twelve projective charts, three for each
polynomial.  Each chart is split into the two sectors obtained by blowing up a corner
of the cone.  For (b), the verification shows that $S_j^\flat$ and $\Lambda^\flat$ are
coprime when $m_3=m_4=m_5$ and $j\ne5$ (Appendix~\ref{app:verification}).}

\begin{lemma}[computer-assisted]\label{lem:d2-L5}
Let $m\in\Rp^5$ with $m_3=m_4=m_5$.  Every label at $m$ whose triangle lies on $S_5$ \textup{(}an
all-equal Lagrange class or a fold triangle\textup{)} and which satisfies $\Phi=0$ or is a fold
label is an isolated point of $\mathcal X_m$.  At each such $m$ these are $72$ labels of the all-equal
classes and $1152$ fold labels, $1224$ in all.
\end{lemma}

{We first normalize the masses.  Common mass scaling gives an exact
isomorphism of the gauged chart, so we may take $m_3=m_4=m_5=1$ and hence
$\delta_*:=M_T^{-1/3}=3^{-1/3}$.
Every positive pair $(m_1,m_2)$ has a unique parametrization by
$(\nu,M)\in(0,1)\times\Rp$ of the form}
\[
  m_1=\frac{4M\nu^2}{(1+\nu^2)^2},\qquad m_2=\frac{M(1-\nu^2)^2}{(1+\nu^2)^2},
\]
{Indeed, $M=m_1+m_2$ and $\nu=1/(\chi+\sqrt{1+\chi^2})$, where
$\chi=\sqrt{m_2/m_1}=(1-\nu^2)/(2\nu)$.}

{Consider first the $72$ labels in the all-equal classes.
They satisfy $\Phi=0$, have invertible $J_T(p)$, and satisfy $q(0)=q'(0)=0$.
The next coefficient is obtained from $F^Z_5=\kappa_3t^3+O(t^4)$ along $C$, which gives
$q''(0)=-2\kappa_3$.
For the nine labels of one sign class with $D_{BB'}=\omega^k\delta_*$,
$(W_3,W_4,W_5)=(\omega,\omega^2,1)$ and $D_{12}=\omega^jM^{-1/3}$, the exact computation gives}
\[
  \kappa_3=\omega^jM^{5/3}R_1(1-i\chi)+3\omega^k\delta_*MR_2(1-i\chi^5),\qquad
  R_1=\frac{18\nu^2(1-\nu^2)}{(1+\nu^2)^3},\quad R_2=\frac{6\nu^2(1+\nu^2)}{(1-\nu^2)^3}.
\]
{To obtain these nine coefficient identities from the case $j=k=0$, first regard
$\nu,b$ as algebraically independent indeterminates, set $M=b^3$, and work over
$\mathbb Q(\omega,i,\delta_*)(\nu,b)$, where $\delta_*^3=1/3$.
The commuting automorphisms
\[
  \delta_*\mapsto\omega^k\delta_*,\qquad b\mapsto\omega^{-j}b
\]
fix $\nu$, $\omega$, $i$ and $b^3$, hence the masses and the chart polynomials.
They carry the representative label and its formal implicit solution to those with
$D_{BB'}=\omega^k\delta_*$ and $D_{12}=\omega^jb^{-1}$.
In particular, $b^5\mapsto\omega^jb^5$, giving the displayed identity for $\kappa_3$.

We now specialize $0<\nu<1$ and $b=M^{1/3}>0$.
The label coordinates are regular throughout this range, and the selected-coordinate
Jacobian is invertible: its triangle block is $J_T$, its pair block is invertible by
Lemma~\ref{lem:d2-second}, and its crossed block has diagonal entries $2/D_{AB}\ne0$.
At each Taylor order, the implicit equations determine the next coefficient by applying
the inverse of this same Jacobian to expressions in the preceding coefficients.
Thus the Taylor coefficients are regular at every such specialization.
Clearing the nonvanishing denominators in the formal identities and then specializing
proves the formula at all positive masses.
The remaining $63$ labels are obtained by exchanging bodies $3$ and $4$, conjugating,
and reversing $t$.

It remains to see that the two terms in the displayed coefficient cannot cancel.
We have $R_1,R_2>0$, and $(1-i\chi)/(1-i\chi^5)$ is a positive multiple of
$1+\chi^6+i(\chi^5-\chi)$.
Moreover, $|\chi^5-\chi|\le\chi+\chi^5\le1+\chi^6$, since
$(1+\chi^6)-(\chi+\chi^5)=(1-\chi)(1-\chi^5)\ge0$.
The quotient of the two terms of $\kappa_3$ is therefore $\omega^{j-k}$ times a number
whose argument has absolute value at most $\pi/4$, and cannot equal $-1$.
Thus $\kappa_3\ne0$, and Lemma~\ref{lem:d2-isolation}(b) gives isolatedness.}

{We treat the fold labels of rank $20$ first.
Among the fold labels, $1080$ have a Jacobian of rank $20$ for the
twenty-two polynomials.  At each such label, twenty equations together with
$h=W_4-W_4(p)$ have an invertible derivative.
For $864$ labels these equations are $F_{\mathrm{sel}}$; for the remaining $216$ they are
all equations other than $F^Z_2$ and $F^Z_3$.
Along the resulting local curve, an omitted equation has a nonzero coefficient of
$h^3$.  Lemma~\ref{lem:d2-isolation}(c) therefore gives isolatedness for these labels.}

{The remaining $72$ fold labels have rank $19$ and form one orbit
under the symmetries described below.
We choose a representative $p_*=p_*(\nu,M)$ whose pair coordinates are as in
Lemma~\ref{lem:d2-fibre}(a), with $D_{12}=M^{-1/3}$, and whose fold triangle is}
\[
  (W_3,W_4,W_5)=(\omega,\omega^2,1),\qquad Z_B=-\delta_*W_B,\qquad
  (D_{34},D_{35},D_{45})=\delta_*(1,\omega,\omega^2).
\]
{The six crossed inverse distances of this representative are
\[
\begin{aligned}
 (D_{13},D_{14},D_{15})&=-\frac{1+\nu^2}{1-\nu^2}(\omega,\omega^2,1),\\
 (D_{23},D_{24},D_{25})&=\frac{i(1+\nu^2)}{2\nu}(\omega,\omega^2,1).
\end{aligned}
\]
This sign choice fixes the coordinate $t$ relative to the Taylor coefficients below;
simultaneously reversing all crossed signs would reverse this coordinate.}
{For this representative, we use Lemma~\ref{lem:d2-corank-two} with
$F_{R'}=(F^W_2,F^W_4,F^W_5)$, with $F_R$ consisting of the other nineteen equations, and
with $y=Z_3$.
The two kernel coordinates are the chart parameter $t$ and $a=Z_3-Z_3(p_*)$.
At $t=0$, the derivative of $(F_R,t,Z_3)$ is block lower triangular, as in the proof of
Lemma~\ref{lem:d2-second}.
Its diagonal blocks are $(1)$ for $dt$; the triangle block}
(rows $F^Z_3,F^W_3,F^Z_4,F^Z_5,T_{34},T_{35},T_{45},G_2,dZ_3$, columns $Z_B,W_B,D_{BB'}$), with
determinant $18\delta_*$; the pair block (rows $F^Z_1,F^W_1,F^Z_2,T_{12},G_1$, columns
$Z_A,W_A,D_{12}$), with determinant $-3M^{-1/3}$; and the diagonal crossed block $(2/D_{AB})$,
with determinant $-512i\,\nu^3(1-\nu^2)^3/(1+\nu^2)^6$.
{All these determinants are nonzero, so the required invertibility
holds for every $(\nu,M)$.
The differentials of the three omitted equations vanish on the two-dimensional kernel
of $dF_R(p_*)$, establishing (a).
We write the reduced equations as $g=(\psi_2,\psi_4,\psi_5)$, where
$\psi_k=F^W_k\circ\phi$.}

At $M=1$ the jet of $g$ through degree four was computed exactly over
$\mathbb Q(\nu)(\sqrt3,i,\delta_*)$, and the nineteen rows of $F_R$ vanish through degree four on
the computed jet of $\phi$.  In the quadratic part, $\psi_2^{(2)}=0$, there is no $a^2$ term,
and the $3\times3$ coefficient matrix has rank one.  So $g^{(2)}=\psi_4^{(2)}\,n$ for a fixed
vector $n=(0,1,n_5)$, where $\psi_4^{(2)}=t(\gamma_{11}a+\gamma_{02}t)$ and
\[
  \gamma_{11}=\frac{9i\delta_*(\nu-i)(\nu+i)^3\,\Upsilon(\nu)}{4\nu(\nu^2-1)(\nu^2+1)^2},\qquad
  \Upsilon(\nu)=\nu^4-2i\nu^3-6\nu^2+2i\nu+1=\tfrac14\bigl(3(\nu-i)^4+(\nu+i)^4\bigr).
\]
For real $\nu$, $|3(\nu-i)^4|=3(\nu^2+1)^2>|(\nu+i)^4|$, so $\Upsilon(\nu)\ne0$ and
$\gamma_{11}\ne0$ for $0<\nu<1$.  Hence the common zeros of $g^{(2)}$ are two lines: the
\emph{boundary line} $\xi_\partial=(1,0)$, tangent to $\{t=0\}$, and the \emph{slanted line}
$\xi_s=(-\gamma_{02},\gamma_{11})$.  Moreover $\Pi_\xi=n\,d\psi_4^{(2)}(\xi)$, with
$d\psi_4^{(2)}(\xi_\partial)=(0,\gamma_{11})$ and
$d\psi_4^{(2)}(\xi_s)=\gamma_{11}(\gamma_{11},\gamma_{02})$, both nonzero.  So (b) holds.
On the slanted line take $N=3$: the matrix $(\Pi_{\xi_s}\,|\,r_3(\xi_s))$ has the $2\times2$
minor
\[
  \mu_3=-\frac{3^{10}(\nu-i)^4(\nu+i)^{16}\,\Upsilon(\nu)^5}
             {2^{12}\delta_*\,\nu^6(\nu^2-1)^6(\nu^2+1)^8}\ne0 .
\]
On the boundary line take $e=(0,1)$, so that $\varpi=\gamma_{11}n$, and $N=4$.  Identically in
$\nu$, $r_3(\xi_\partial)=g^{(3)}(1,0)=-b_2\varpi$, and the matrix
$(\Pi_{\xi_\partial}\,|\,r_4(\xi_\partial))$ has the minor
\[
  \mu_4=-\frac{27i\delta_*(\nu-i)^3(\nu+i)^5\,\Upsilon(\nu)}{64\,\nu^3(\nu^2-1)^3}\ne0 .
\]
So (c) holds.  Every coefficient used is regular for $0<\nu<1$: the jets are polynomials in the
coordinates of $p_*$, the masses and the entries of the inverses of the three diagonal blocks,
and the only further division is by $\gamma_{11}$.  Hence the identities above hold at every
$\nu\in(0,1)$.

For general $M$, $p_*(\nu,M)$ is obtained from $p_*(\nu,1)$ by multiplying $W_1,W_2$ by $M^{2/3}$
and $D_{12}$ by $M^{-1/3}$.  The map
$\tilde g(a,\tilde t)=\bigl(M^{-2/3}\psi_2,\psi_4,\psi_5\bigr)(a,\tilde t/M)$ was computed in
exact Laurent arithmetic in $M^{1/3}$, and the selected rows vanish along both lines through the
orders used.  The quadratic part of $\tilde g$ does not depend on $M$, the compatibility
$r_3(\xi_\partial)\in\C\varpi$ persists, and the two minors become $M^{-3}\mu_3(\nu)$ and
$M^{-5/3}\mu_4(\nu)$.  By Lemma~\ref{lem:d2-corank-two}, applied to $\tilde g$, $p_*(\nu,M)$ is
an isolated point of $\mathcal X_m$ for every $0<\nu<1$ and $M>0$.

{To transfer isolatedness from this representative to the other labels,
we use the symmetries of the chart equations.
Because the twenty-two polynomials have real coefficients and $m_3=m_4$, they are
preserved by two commuting involutions: the biholomorphic reversal
$(t,D_{AB})\mapsto(-t,-D_{AB})$ and the antiholomorphic map obtained by conjugating and
exchanging bodies $3$ and $4$.
Both fix $(\nu,M)$ and together carry $p_*$ to the other three sign choices in its orbit.
Multiplying $t$, $Z_B$, $W_A$, $D_{12}$ and $D_{BB'}$ by $\omega^j$ multiplies each
polynomial by a cube root of unity and preserves the gauges.
A permutation of bodies $3,4,5$, followed by the chart-torus action restoring the
gauges, gives a local biholomorphism of $\mathcal X_m$.
The implicit function theorem applies because the transversality determinant is
$2(Z_2-Z_1)W_5=2$.
These $3\cdot6$ maps produce eighteen labels for each of the four sign choices.
The resulting $72$ labels form precisely one orbit under these discrete symmetries and
are disjoint from the $1080$ labels already treated.
Every map used here preserves isolatedness.}

The computations in the proof of Lemma~\ref{lem:d2-L5}, namely the coefficients $\kappa_3$, the
ranks and the coefficients of $h^3$ at the $1080$ fold labels, the determinants, jets and minors
at $p_*$, and the orbit of $p_*$, are exact, over number fields containing $\omega$, $i$ and
$\delta_*$ and over rational function fields over them in $\nu$ and $M^{1/3}$
(Appendix~\ref{app:verification}).

\begin{lemma}[computer-assisted]\label{lem:d2-rays}
\textup{(a)} The zero set of $L_8$ in $\Rp^3$ is $\Rp(p_*,q_*,1)\cup\Rp(q_*,p_*,1)$ with
$p_*=0.67987758\ldots$ and $q_*=0.82343785\ldots$.
\textup{(b)} The zero set of $L_1$ in $\Rp^3$ is $\Rp(a_*,b_*,1)\cup\Rp(a_*,1,b_*)$ with
$a_*=1.02227400\ldots$ and $b_*=0.90559635\ldots$; at $(a_*,b_*,1)$ the zeros of $S_1$ that
satisfy \textup{(25)} form exactly one pair of complex conjugate non-real shapes.
\textup{(c)} Consequently the zero sets of $L_1L_2L_3$ and of $L_6L_7L_8$ in $\Rp^3$ are the rays
through the six permutations of $(a_*,b_*,1)$ and of $(1,p_*,q_*)$, respectively.
\end{lemma}
{For (a), exact Bernstein subdivision of the mass simplex locates the
possible zeros, and the Krawczyk method in rational interval arithmetic encloses them.
For (b), subresultants eliminate the shape variable.
Bernstein subdivision and the signs of resultants locate the mass zeros, which are
again enclosed by the Krawczyk method.  The quadratic subresultant has negative
discriminant in these enclosures, giving the conjugate non-real shapes
(Appendix~\ref{app:verification}).
Part (c) follows from (a), (b) and the permutation symmetry following
Definition~\ref{def:d2-mass}.}

\begin{lemma}[computer-assisted]\label{lem:d2-Lray}
Let $m\in\Rp^5$ with $m'$ on one of the rays of Lemma~\ref{lem:d2-rays}(c), and let $p$ be a
same-sign label at $m$ with $\Phi(p)=0$.  Then $J_T(p)$ is invertible and $\Psi_T(p)\ne0$, so $p$
is isolated at second order.
\end{lemma}

{The enumeration includes all $48$ Lagrange triangles at each
permutation of $(1,p_*,q_*)$ and the six Euler triangles on the conjugate shape pair at
each permutation of $(a_*,b_*,1)$.
At these mass vectors, only the factor defining the ray vanishes among
$L_1,\dots,L_8,L_5'$.  Proposition~\ref{prop:d2-first}(a) therefore excludes any other
triangle from carrying a same-sign label with $\Phi=0$.
For $36+36$ triangles, there is a sign vector satisfying $\sum_Bm_Bs_Bv_B=0$, unique up
to a simultaneous sign change.  Their $J_T$ are invertible, and complex interval
arithmetic gives $|\Psi_T|>0.36$ ($L_6$--$L_8$) and $|\Psi_T|>0.3$ ($L_1$--$L_3$) at the
normalized representatives.

Common mass scaling preserves isolatedness and the nonvanishing of $\Psi_T$, so the
argument extends to each entire ray.  The displayed numerical lower bounds apply only
to the normalized representatives.  Since $\Psi_T$ does not involve $m_1,m_2$, the pair
masses impose no further restriction here.
The Krawczyk method encloses the shapes and the square and cube roots; the solutions of
the linear systems are also enclosed rigorously (Appendix~\ref{app:verification}).}

\begin{lemma}[computer-assisted]\label{lem:d2-K}
Let $m\in\Rp^5$ with $m_1=m_2$, let $p$ be a one-different label at $m$ with odd body $5$ and
$\Phi(p)=0$, and let $S_j$ carry its triangle, so that $K_j(m')=0$.  Then $p$ is an isolated
point of $\mathcal X_m$.  More precisely:
\begin{alphalist}
  \item $j\notin\{3,6,7,8\}$, by Lemma~\ref{lem:d2-empty};
  \item if $j\in\{1,2\}$: at every zero of $K_1$ in $\Rp^3$ the $S_1$ shapes compatible with
  \textup{(28)} at $y=0$ reduce to one simple common root, which is real and lies on one of two
  explicit real branches; $J_T(p)$ is invertible and $q'(0)\ne0$.  The case $j=2$ is the image of
  $j=1$ under the exchange of bodies $3$ and $4$;
  \item {if $j=4$: $m'$ lies in one of the positive cones over the four arcs
  $k\mapsto\mathrm{II}(k),\mathrm{III}(k)$, $5^{-1/4}<k<1$, and their images under
  $m_3\leftrightarrow m_4$.  Thus $m'=c\,\mathrm{II}(k)$ or $c\,\mathrm{III}(k)$, or an
  exchanged vector, for some $c>0$.  Here}
  $\alpha=4k(1+k^4)$, $\beta=(1-k^2)(1+6k^2+k^4)$, $\gamma=(1+k^2)(6k^2-1-k^4)$,
  $\mathrm{II}=((\gamma^3-\beta^3)\alpha,(\alpha^3-\gamma^3)\beta,(\alpha^3-\beta^3)\gamma)$ and
  $\mathrm{III}=((\gamma^3-\beta^3)\alpha,(\gamma^3+\alpha^3)\beta,(\alpha^3+\beta^3)\gamma)$;
  $J_T(p)$ is invertible and $q'(0)\ne0$;
  \item if $j=5$: $m'\in\Rp(1,1,\rho_*)$, where $\rho_*\in(1.5736409,1.573641)$ is the unique
  positive root of $\rho^8+2\rho^7+3\rho^6-4\rho^5-5\rho^4-6\rho^3-3\rho^2-12\rho-12$, and every
  label at $m$ with a Lagrange triangle is isolated, for all $27$ classes and all $64$
  crossed-sign assignments.
\end{alphalist}
\end{lemma}

{For (b) and (c), $60$ of the $64$ crossed-sign assignments satisfy
$\Phi\ne0$.  The other four form one orbit under conjugation and $t\mapsto-t$ and have
$q'(0)\ne0$.
In (d), the two-equal classes have $\Phi\ne0$.
Among the all-equal classes, $60$ assignments have $\Phi\ne0$ and $4$ have $q'(0)\ne0$.
For the fold classes, an augmented system using $h=W_4-W_4(p)$ gives a nonzero coefficient
of order three for $60$ assignments and of order four for the remaining $4$.
The exact grading makes these computations uniform in the common pair mass $m_1=m_2$.}

{In (b), Bernstein coefficients establish smoothness of the zero set
of $K_1$ in $\Rp^3$.  A linear subresultant over $\mathbb Q$ then determines the common
shape root.  A rational parametrization of $S_1$, together with Sturm sequences,
identifies its two real branches.  An exact reduction and an interval cover of the
branch parameter give $q'(0)\ne0$.

For (c), exact computation on the rational parametrizations of the four projective mass
arcs classifies the real normalized $W$-shapes.  Lemma~\ref{lem:d2-empty}(a) shows that
these include every $S_4$ shape at positive masses.
Common mass scaling, followed by restoration of the chart gauges, is a local analytic
isomorphism.  Isolatedness therefore extends from the arcs to their positive cones.

For (d), exact resultants and a Sturm count determine the zero set of $K_5$ in $\Rp^3$.
Complex rational interval arithmetic classifies the labels, using one triangle from
each orbit under conjugation and exchange of bodies $3$ and $4$
(Appendix~\ref{app:verification}).
Renumbering the bodies as at the start of this section gives the same conclusions for
every choice of the pair and of the odd body.}

\begin{lemma}[computer-assisted]\label{lem:d2-P}
Let $m\in\Rp^5$ with $m_1\ne m_2$, and let $p$ be a one-different label at $m$ with odd body $5$,
$\Phi(p)=0$, whose triangle is not a fold triangle.  Then $J_T(p)$ is invertible and
$\Psi_1(p)\ne0$, so $p$ is isolated at second order.
\end{lemma}

{The verification covers the mass simplex $m_3+m_4+m_5=1$.
For $S_1,S_3,S_4$, it includes every Euler shape, enclosed by the Krawczyk method or
Pellet discs.  For $S_5,S_6,S_8$, it includes every Lagrange shape up to conjugation and
the cube-root action.  Both relative signs $s_3s_4$ are considered.
On each box, either $G$ has a strict sign or ball arithmetic encloses $\Psi_1$ away from
$0$.  On boxes meeting $m_3=0$ or $m_4=0$, where $\Psi_1$ vanishes, the corresponding
transversal derivative is instead enclosed away from zero.}

{The remaining boundary and degeneracy cases require separate arguments.
At the vertices, $G\ne0$.
At the ends $(m_3,m_4)=(0,\frac12)$ for $S_1$ and $S_4$, and $(\frac12,0)$ for $S_4$, a
majority body becomes massless.  An exact edge identity gives interval margins
$0.787$, $0.588$, $0.551$.
At the corner $(\frac12,\frac12,0)$ for $S_3$, an exact weighted blow-up gives weight
two in $m_5$ to the Euler shape parameter.  The resulting determinant and response
numerator are $m_5^{13}$ and $m_5^6$ times factors whose moduli are bounded below by
$6.96\cdot10^8$ and $1.63\cdot10^6$, respectively, in interval arithmetic.
For the six real points of $S_4$ where the Euler label degenerates, and for the corner
$(\frac12,\frac12,0)$ of $S_4$, an exact formula for $\Psi_1$ on the Euler shapes cancels
the degenerating factor.
The remaining factor $S_5$ is treated by exact algebra (Appendix~\ref{app:verification}).}

{On $S_5$, let $g$ be the $W_4$-component of the solution of the
triangle linear system with row $F^Z_5$ omitted and $dW_5=0$, when the right-hand side is
one in row $F^Z_4$ and zero elsewhere.  Put $X=\tilde v_4/W_4-\tilde v_3/W_3$.
Then $\Psi_1=-\frac12m_4gX^2$.
The numerator of $g$ is a nonzero constant times
$m_3(m_3^2+m_3m_4+m_4^2)^2(P_0+\omega P_1)$, where the real polynomials $P_0,P_1$ have
only one common zero in the simplex, at $m_3=m_4=\frac13$.
Also, $X\ne0$ except at this point.
There $m_3=m_4=m_5$, and a label with $G=0$ would have $\rho=-1$, which is impossible.
The cases $S_2$ and $S_7$ follow from $S_1$ and $S_6$ by exchanging bodies $3$ and $4$.}

\begin{lemma}[computer-assisted]\label{lem:d2-fold}
Let $m\in\Rp^5$ and let $p$ be a one-different label at $m$ with odd body $5$, $\Phi(p)=0$ and a
fold triangle.  Then $dF_{\mathrm{sel}}(p)$ has rank $20$ and $Q|_C$ has a simple zero at $p$; in
particular $p$ is excluded.
\end{lemma}

{The triangle Jacobian $J_T(p)$ has corank exactly one.
Its kernel vector and left null vector $\ell$ are computed explicitly.
The $t$-column has entries $-\varphi_B/W_B$ in the triangle $Z$-rows.
Applying $\ell$ to this column gives $c_\sigma X$, with $X$ as in
Lemma~\ref{lem:d2-P}.
Here $X\ne0$ at every label under consideration: the exceptional point $m_3=m_4=m_5$
has no one-different label with $\Phi=0$ at positive masses.
Thus the additional parameter resolves the corank, and the curve $C$ exists with local
parameter $h=W_4-W_4(p)$.
At $p$, it satisfies $dt/dh=0$ and $dQ/dh=-\frac12c_\sigma m_4X\ne0$, which gives the
required simple zero.

These statements are verified exactly over $\mathbb Q(m_3,m_4)(\omega)$, treating
$m_3,m_4$ as indeterminates.
The computation gives the kernel and left null vectors, an order-eight minor of
$J_T(p)$ that is nonzero throughout the simplex, and the zero set of $X$
(Appendix~\ref{app:verification}).}

\begin{theorem}\label{thm:d2}
{For $m\in\Rp^5$, no curve in $A_m$ has a Diagram~2 germ.
In particular, Theorem~\ref{thm:exclusion}\textup{(a)} holds for $r=2$.}
\end{theorem}

\begin{proof}
{Suppose that $\gamma$ is a Diagram~2 germ on a curve in $A_m$, and
let $p$ be its limit label given by Proposition~\ref{prop:d2-cover}.
For each $B$, the quantity $\sigma_Bx$ is the limit of
$m_2\delta_{2B}/(m_1\delta_{1B})$ along $\gamma$.
This limit is independent of the chart, so the crossed-sign pattern of $p$ is
determined by $\gamma$.
Renumbering bodies $3,4,5$ before applying Proposition~\ref{prop:d2-cover} permutes the
pattern.  We may therefore assume that $p$ is same-sign or that its odd body is $5$.

By Lemma~\ref{lem:d2-isolation}(a), it suffices to exclude $p$.
Proposition~\ref{prop:d2-first} gives $\Phi(p)=0$.
Choose a shape factor $S_j$ vanishing on the triangle of $p$, as in
Lemma~\ref{lem:d2-fibre}.}

{Suppose first that $p$ is same-sign.
Proposition~\ref{prop:d2-first}(a) gives $L_j(m')=0$, or $L_5L_5'=0$ if $j=5$.
Lemma~\ref{lem:d2-empty} rules out $j=4$ and $L_5'=0$.
If $L_5=0$, then $m_3=m_4=m_5$, and Lemma~\ref{lem:d2-empty}(b) gives $j=5$;
Lemma~\ref{lem:d2-L5} excludes these labels.
In all other cases, $j\in\{1,2,3,6,7,8\}$ and $m'$ lies on one of the rays in
Lemma~\ref{lem:d2-rays}(c), so Lemma~\ref{lem:d2-Lray} applies.}

{For a one-different label with $m_1=m_2$, exclusion follows from
Lemma~\ref{lem:d2-K}.
When $m_1\ne m_2$, Lemma~\ref{lem:d2-fold} applies to fold triangles and
Lemma~\ref{lem:d2-P} to all other triangles.}

{Thus every possible $p$ is excluded; for those isolated at second
order this follows from Definition~\ref{def:d2-psi}.
This contradicts Lemma~\ref{lem:d2-isolation}(a).}
\end{proof}

\begin{remark}\label{rem:d2-scope}
{(i) Theorem~\ref{thm:d2} excludes every Diagram~2 germ, regardless
of whether a four-product is unbounded.  It applies to complex configurations and is
therefore stronger than the assertion needed in Theorem~\ref{thm:exclusion}(a).
The input from \cite{AK} consists of Fact~\ref{fact:rules}, used in
Proposition~\ref{prop:d2-cover}, and Lemma~\ref{lem:germ}.
The mass conditions (25)--(28) of \cite[\S8.2]{AK} follow from the torque identity in
Proposition~\ref{prop:d2-first}.}

{(ii) Two further computations show that the layer $u=\infty$ of
the $P_j$ misses the positive cone, and that a non-degenerate triple of positions with
positive centre-of-mass masses cannot lie on two curves $S_i\ne S_j$
(Appendix~\ref{app:verification}).
Neither statement is needed in the proof above.}
\end{remark}



\section{Diagram 3}\label{sec:d3}

{Diagram~3 of \cite[Fig.~11]{AK} occurs at infinity on Roberts' continuum
\cite[\S7.2]{AK}. We use the numbering of \cite[\S5.2 and \S8.3]{AK}, in which the clusters are
$z:12\ldots34$, $w:14\ldots23$. The diagram is the quadrilateral $1234$, with $z$-edges $14$, $23$ and
$w$-edges $12$, $34$; neither the diagonals nor body $5$ carry a stroke. Exchanging $z$ and $w$
has the same effect as exchanging bodies $2$ and $4$. Thus, after renumbering, every Diagram-3 germ
has the form considered below. Our argument works near this first limit; the second singular
limit used in \cite[proof of Thm.~6(iii)]{AK} is unnecessary.}

\subsection{The crossed-cluster chart}
{Apart from the gauge, the equations in \eqref{eq:system} are invariant under
$(z,w,\delta)\mapsto(\mu z,\mu^{-1}w,\delta)$, $\mu\in\C^*$. We therefore impose
$z_{13}=w_{13}$, following \cite[p.~543]{AK} for Roberts' continuum. To resolve the two pairs
of clusters, introduce a parameter $t$ and coordinates $A,B,x,y,\hat z_5,C,D,u,v,\hat w_5$ and $D_{kl}$
by the following substitutions:}
\begin{equation}\label{eq:d3-chart}
\begin{gathered}
z_1=\tfrac{A}{t},\quad z_2=\tfrac{A}{t}+tx,\quad z_3=\tfrac{B}{t},\quad z_4=\tfrac{B}{t}+ty,
\quad z_5=t^2\hat z_5,\\
w_1=\tfrac{C}{t},\quad w_2=\tfrac{D}{t},\quad w_3=\tfrac{D}{t}+tv,\quad w_4=\tfrac{C}{t}+tu,
\quad w_5=t^2\hat w_5,\\
\delta_{kl}=D_{kl}\ \ (kl=12,14,23,34),\qquad
\delta_{kl}=t\,D_{kl}\ \ (kl=13,24,15,25,35,45).
\end{gathered}
\end{equation}
{Multiply the $z$- and $w$-equations for bodies $1,\dots,4$, together with the gauge
equation, by $t$, and those for body $5$ by $t^{-2}$. The resulting $21$ equations are polynomial:
$F=F^0+t^2F^2+\dots+t^6F^6=0$. In particular, there is no term linear in $t$.
Setting $t=0$ gives the \emph{boundary equations} $F^0=0$, whose left-hand sides are}
\begin{equation}\label{eq:d3-boundary}
\begin{gathered}
R^z_1=A-m_4D_{14}^3(A-B),\qquad R^z_4=B-m_1D_{14}^3(B-A),\\
R^z_2=A-m_3D_{23}^3(A-B),\qquad R^z_3=B-m_2D_{23}^3(B-A),\\
R^w_1=C-m_2D_{12}^3(C-D),\qquad R^w_2=D-m_1D_{12}^3(D-C),\\
R^w_3=D-m_4D_{34}^3(D-C),\qquad R^w_4=C-m_3D_{34}^3(C-D),\\
\hat z_5+(m_1D_{15}^3+m_2D_{25}^3)A+(m_3D_{35}^3+m_4D_{45}^3)B,\\
\hat w_5+(m_1D_{15}^3+m_4D_{45}^3)C+(m_2D_{25}^3+m_3D_{35}^3)D,\\
D_{12}^2x(D-C)-1,\quad D_{34}^2y(C-D)-1,\quad D_{14}^2u(B-A)-1,\\
D_{23}^2v(B-A)-1,\quad D_{13}^2(B-A)(D-C)-1,\quad D_{24}^2(B-A)(C-D)-1,\\
D_{15}^2AC-1,\quad D_{25}^2AD-1,\quad D_{35}^2BD-1,\quad D_{45}^2BC-1,\quad (B-A)-(D-C);
\end{gathered}
\end{equation}
{Here $R^z_k$, $R^w_k$ are the boundary equations for body $k$.
We restrict to the slice $B=A+1$. Proposition~\ref{prop:d3-cover} justifies this choice
by showing that every germ can be represented with $t=1/z_{13}$.}

\subsection{Representation of germs in the chart}
{We first verify that this chart contains every germ with the prescribed diagram.
We write $s$ for the parameter of a germ.}

\begin{proposition}[Coverage]\label{prop:d3-cover}
Let $\gamma$ be a Diagram-3 germ, renormalized as in Lemma~\ref{lem:germ}.
\begin{romanlist}
  \item Bodies $1,\dots,4$ are circled in both colours, body $5$ in neither.
  \item The sides are maximal edges: $r_{12},r_{14},r_{23},r_{34}\approx1$,
  $z_{14},z_{23},w_{12},w_{34}\approx\epsilon^{-2}$ and $w_{14},w_{23},z_{12},z_{34}\approx\epsilon^2$.
  \item $z_{13},w_{13},z_{24},w_{24},z_{k5},w_{k5}\approx\epsilon^{-2}$ for $k\le4$, and
  $z_5,w_5\preccurlyeq\epsilon^4$.
  \item Renormalize instead by $\lambda=(w_{13}/z_{13})^{1/2}$, so that $z_{13}=w_{13}$, and put
  $t=1/z_{13}\approx\epsilon^2$.  Then $B-A=1$ and every chart coordinate converges as $s\to0$; the
  limit $X_0$ satisfies \eqref{eq:d3-boundary}, and $A,B,C,D,x,y,u,v$ and all $D_{kl}$ are
  nonzero at $X_0$.
\end{romanlist}
\end{proposition}

\begin{proof}
{Lemma~\ref{lem:germ} allows us to apply the order relations of \cite[\S3]{AK}
along $\gamma$, and hence the rules in Fact~\ref{fact:rules}.

For (i), each side is the only stroke of its colour incident to either endpoint.
Rule~1a \cite[pp.~545--546]{AK} therefore circles $1,\dots,4$ in both colours.
Since every circle requires a stroke of the same colour, $5$ has no circle.

For (ii), the isolated component $\{1,4\}$ of the $z$-diagram has $z$-circles.
Rule~1e \cite[p.~546]{AK} makes its stroke $14$ maximal. The same argument applies to $23$,
and to $12$, $34$ in the $w$-diagram. The right equality case of Estimate~2
\cite[p.~548]{AK} gives the asserted orders.

For (iii), (ii) gives $z_{13}=z_{14}-z_{34}$, $w_{13}=w_{12}+w_{23}$, $z_{24}=z_{23}+z_{34}$ and
$w_{24}=w_{14}-w_{12}$, all of order $\approx\epsilon^{-2}$.
Part (i) also gives $z_k\approx Z=\epsilon^{-2}\succ z_5$ for $k\le4$, and the corresponding
relations for $w$. In the equation
$z_5(1-\sum_{k\le4}m_k\delta_{k5}^3)=-\sum_{k\le4}m_k\delta_{k5}^3z_k$ for body $5$,
we have $\delta_{k5}\approx\epsilon^2$. It follows that $z_5\preccurlyeq\epsilon^4$.
The same argument applies to $w_5$.

For (iv), (iii) shows that $\lambda$ differs from the normalization factor in
Lemma~\ref{lem:germ} by a Puiseux series of order $0$. Thus the orders in (i)--(iii) remain
unchanged. We now have $B-A=tz_{13}=1$, and the coordinates $A=z_1/z_{13}$, $B=z_3/z_{13}$,
$C=w_1/w_{13}$, $D=w_2/w_{13}$, $x=z_{12}z_{13}$, $y=z_{34}z_{13}$, $u=w_{14}w_{13}$,
$v=w_{23}w_{13}$ and $D_{kl}$ ($=\delta_{kl}$ for the sides, $\delta_{kl}z_{13}$ otherwise) are
$\approx1$. The remaining coordinates satisfy $\hat z_5=z_5z_{13}^2$,
$\hat w_5=w_5w_{13}^2\preccurlyeq1$. Every chart coordinate therefore has a finite limit,
and the limiting point satisfies $F=0$.}
\end{proof}

By Proposition~\ref{prop:d3-cover}(ii) and~(iii), every Diagram-3 germ has an unbounded
four-product: $p_{12}=r_{12}^2r_{34}^2r_{35}^2r_{45}^2\approx\epsilon^{-8}$.

\subsection{The first mass relation}

\begin{proposition}[First condition]\label{prop:d3-first}
At every point of $\{t=0\}$ satisfying \eqref{eq:d3-boundary}, $m_1m_3=m_2m_4$.  Hence every
Diagram-3 germ satisfies (29) of \cite[\S8.3]{AK}.
\end{proposition}

\begin{proof}
Since $m_1R^z_1+m_4R^z_4=m_1A+m_4B$ and $m_2R^z_2+m_3R^z_3=m_2A+m_3B$ identically,
\begin{equation}\label{eq:d3-syzygy}
(B-A)(m_1m_3-m_2m_4)=(m_1+m_4)(m_2R^z_2+m_3R^z_3)-(m_2+m_3)(m_1R^z_1+m_4R^z_4),
\end{equation}
and likewise
$(D-C)(m_1m_3-m_2m_4)=(m_1+m_2)(m_3R^w_3+m_4R^w_4)-(m_3+m_4)(m_1R^w_1+m_2R^w_2)$.  With
$T_{13}=D_{13}^2(B-A)(D-C)-1$ (the distance equation of $13$) and $g=(B-A)-(D-C)$ (the gauge),
$D_{13}^2(B-A)^2-1=T_{13}+D_{13}^2(B-A)g$.  At a boundary point $T_{13}=g=0$, so
$D_{13}^2(B-A)^2=1$, $B-A\ne0$, and \eqref{eq:d3-syzygy} gives $m_1m_3=m_2m_4$.  The second
statement follows from Proposition~\ref{prop:d3-cover}(iv).
\end{proof}

{These identities express in chart coordinates the argument of \cite[\S5.2]{AK}:
Rule~1c makes $m_1z_1+m_4z_4$ and $m_2z_2+m_3z_3$ non-maximal, whereas
$z_1\sim z_2$ and $z_3\sim z_4$.}

\subsection{Boundary solutions and analytic continuation}
{Assume henceforth that $m\in\Rp^5$ satisfies (29), and set
$M=m_1+m_2+m_3+m_4$. After eliminating $B=A+1$, denote the remaining $19$ chart coordinates
by $X$. We will solve nineteen independent equations by the implicit function theorem,
leaving two equations that must vanish along any germ.

Let $F_{z_k},F_{w_k}$ denote the rescaled equations for body $k$.
Define $\Lambda$ to be the $19\times21$ matrix selecting the combinations
$m_1F_{z_1}+m_4F_{z_4}$, $F_{z_4}-F_{z_1}$, $F_{z_3}-F_{z_2}$,
$m_1F_{w_1}+m_2F_{w_2}$, $F_{w_2}-F_{w_1}$, $F_{w_3}-F_{w_4}$,
together with the other $13$ equations. The two remaining linear combinations are
$\ell_zF=F_{z_1}-F_{z_2}-\frac{m_3}{m_2}F_{z_3}+\frac{m_4}{m_1}F_{z_4}$ and
$\ell_wF=F_{w_1}+\frac{m_2}{m_1}F_{w_2}-\frac{m_2}{m_1}F_{w_3}-F_{w_4}$.}

{\begin{lemma}\label{lem:d3-fibre}
\begin{alphalist}
  \item On the slice, \eqref{eq:d3-boundary} has exactly $3^4\cdot2^6=5184$ solutions, which we call
  \emph{labels}. They are obtained as follows: $D_{12},D_{34},D_{14},D_{23}$ are arbitrary cube roots of $\frac1{m_1+m_2}$,
  $\frac1{m_3+m_4}$, $\frac1{m_1+m_4}$, $\frac1{m_2+m_3}$; $D_{13},D_{24},D_{15},D_{25},D_{35},
  D_{45}$ are arbitrary square roots of $1$, $-1$, $\frac M{m_3}$, $-\frac M{m_4}$, $\frac M{m_1}$,
  $-\frac M{m_2}$; $A=-\frac{m_4}{m_1+m_4}$, $C=-\frac{m_2}{m_1+m_2}$, $D=\frac{m_1}{m_1+m_2}$,
  $x=(m_1+m_2)D_{12}$, $y=-(m_3+m_4)D_{34}$, $u=(m_1+m_4)D_{14}$, $v=(m_2+m_3)D_{23}$, and
  $\hat z_5,\hat w_5$ are determined by the equations for body $5$. These boundary solutions
  are independent of the mass $m_5$.
\item At every label the matrix $\partial_X(\Lambda F^0)$ is invertible.

\end{alphalist}
\end{lemma}}

{\begin{proof}
On the slice $B=A+1$, the three selected combinations of the $z$-equations are
\[
 m_1R^z_1+m_4R^z_4=(m_1+m_4)A+m_4,\qquad
 R^z_4-R^z_1=1-(m_1+m_4)D_{14}^3,
\]
\[
 R^z_3-R^z_2=1-(m_2+m_3)D_{23}^3.
\]
The gauge gives $D-C=1$, while the weighted combination of the $w$-equations gives
$m_1C+m_2D=0$. The other two selected combinations are
\[
 R^w_2-R^w_1=(D-C)\bigl(1-(m_1+m_2)D_{12}^3\bigr),\qquad
 R^w_3-R^w_4=(D-C)\bigl(1-(m_3+m_4)D_{34}^3\bigr).
\]
These equations determine $A,C,D$ and the four cubes listed in (a).
The four distance equations for the sides then determine $x,y,u,v$.
The other six distance equations determine the six squares, and the equations for body $5$
determine $\hat z_5,\hat w_5$.
The relation $m_1m_3=m_2m_4$ ensures that the omitted boundary coordinate rows
also vanish. Since all roots are nonzero and simple, each choice
of the four cube roots and six square roots gives exactly one label.
This accounts for all $3^4 2^6$ labels.

The same elimination proves (b). Let a vector belong to the kernel of the differential.
The linear equations force its variations of $A,C,D$ to vanish.
The four cube equations then force the variations of the side coordinates $D_{kl}$ to vanish,
because their coefficients are nonzero multiples of $D_{kl}^2$.
The four side distance equations next force the variations of $x,y,u,v$ to vanish.
Each remaining distance equation does the same for its own $D_{kl}$, since its
derivative in that coordinate is $2/D_{kl}\ne0$.
Finally, the two equations for body $5$ force the variations of $\hat z_5,\hat w_5$ to vanish.
The square matrix $\partial_X(\Lambda F^0)$ therefore has trivial kernel.
\end{proof}

}
{For all $X$, we have $\ell_zF^0=-B(m_1m_3-m_2m_4)/(m_1m_2)$ and
$\ell_wF^0=-(D-C)D_{34}^3(m_1m_3-m_2m_4)/m_1$. Thus $\ell_z,\ell_w$ annihilate
$J_0=\partial_XF^0(X_0)$. Since $\Lambda J_0$ is invertible, $J_0$ has rank $19$,
$\ell_z,\ell_w$ span its left kernel, and the $21$ rows of $\Lambda,\ell_z,\ell_w$ are independent.}

\begin{lemma}\label{lem:d3-graph}
Let $X_0$ be a label.  There is a unique analytic $t\mapsto X(t)$ near $0$ with $X(0)=X_0$ and
$\Lambda F(t,X(t))=0$; put $G(t)=(\ell_zF(t,X(t)),\ell_wF(t,X(t)))=\sum_{k\ge1}G_kt^k$.  If the
chart coordinates of a Diagram-3 germ converge to $X_0$, then along the germ they are $X(t)$, and
$G\equiv0$.  Every analytic $Y(t)$ with $Y(0)=X_0$ and $F(t,Y(t))\equiv0$ equals $X(t)$.
\end{lemma}

\begin{proof}
{The analytic implicit function theorem applies by Lemma~\ref{lem:d3-fibre}(b)
and gives the unique analytic graph. After a reparametrization $s\mapsto s^N$,
the chart coordinates $Y(s)$ and $t(s)=1/z_{13}(s)$ of the germ are convergent power series
(Lemma~\ref{lem:germ} and Proposition~\ref{prop:d3-cover}), with $Y(0)=X_0$, $t(0)=0$ and
$t\not\equiv0$. The identity $\Lambda F(t(s),Y(s))=0$ and uniqueness imply
$Y(s)=X(t(s))$. Consequently, $F(t,X(t))$ vanishes on the image of $t(s)$,
which contains a neighbourhood of $0$. Thus $F(t,X(t))\equiv0$ and $G\equiv0$.
The last assertion follows from the same uniqueness.}
\end{proof}

{The lemma gives a necessary condition at every label: $G\equiv0$,
even when the germ is ramified. Since $F$ has no term linear in $t$, its analytic graph has
the expansion $X(t)=X_0+X_2t^2+X_3t^3+\cdots$. Moreover, $\ell J_0=0$, where
$\ell=(\ell_z,\ell_w)$, so the coefficient $G_k$ is independent of $X_k$.
The first coefficients, evaluated at $X_0$, are $G_2=\ell F^2$, $G_3=\ell F^3$,
$G_4=\ell(F^4+\partial F^2\cdot X_2+\tfrac12\partial^2F^0[X_2,X_2])$ and
$G_5=\ell(F^5+\partial F^3\cdot X_2+\partial F^2\cdot X_3+\partial^2F^0[X_2,X_3])$.}

\subsection{Orders three and four}

\begin{lemma}[computer-assisted]\label{lem:d3-order3}
At every label $G_2=0$ and
$G_3=\bigl(\frac{m_1+m_2}{m_1^2m_2}\Omega,\frac{m_1+m_4}{m_1^2m_4}\Omega\bigr)$, where
$\Omega=\Omega_{\mathrm{re}}+\Omega_{\mathrm{im}}$,
\[
\begin{aligned}
\Omega_{\mathrm{re}}&=m_1m_2m_4D_{13}+m_5m_1^2D_{15}+m_5m_2m_4D_{35},\\
\Omega_{\mathrm{im}}&=m_1m_2m_4D_{24}+m_5m_1m_2D_{25}+m_5m_1m_4D_{45}.
\end{aligned}
\]
\end{lemma}

{This is an exact symbolic computation over $\mathbb Q(m_1,m_2,m_4,m_5)$,
with $m_3=m_2m_4/m_1$. It evaluates $\ell F^2$ and $\ell F^3$ at the labels and reduces the
results modulo the cubic and quadratic equations for $D_{kl}$.
The verification is described in Appendix~\ref{app:verification}.}

\begin{corollary}\label{cor:d3-order3}
If $G_3=0$ at some label, then $N_{\mathrm{re}}(m)=N_{\mathrm{im}}(m)=0$, where $Q=m_5^2M$ and
\begin{align*}
N_{\mathrm{re}}&=\bigl((m_1m_2m_4)^3-Q(m_1^6+m_2^3m_4^3)\bigr)^2-4Q^2m_1^6m_2^3m_4^3,\\
N_{\mathrm{im}}&=\bigl((m_1m_2m_4)^3-Qm_1^3(m_2^3+m_4^3)\bigr)^2-4Q^2m_1^6m_2^3m_4^3 .
\end{align*}
As $N_{\mathrm{re}}-N_{\mathrm{im}}+QS=0$ with
$S=(m_1^3-m_2^3)(m_1^3-m_4^3)\bigl(2(m_1m_2m_4)^3-Q(m_1^3+m_2^3)(m_1^3+m_4^3)\bigr)$, also
$S(m)=0$.
\end{corollary}

\begin{proof}
At positive masses $D_{13},D_{15},D_{35}$ are real and $D_{24},D_{25},D_{45}$ purely imaginary
(Lemma~\ref{lem:d3-fibre}), so $\Omega_{\mathrm{re}}=\Omega_{\mathrm{im}}=0$.  If
$\alpha+\beta+\gamma=0$ then $(\alpha^2-\beta^2-\gamma^2)^2-4\beta^2\gamma^2=\prod_\pm(\alpha+\beta
\pm\gamma)\prod_\pm(\alpha-\beta\pm\gamma)=0$.  Inserting the squares $(m_1m_2m_4)^2$,
$m_5^2m_1^4M/m_3$, $m_5^2m_2^2m_4^2M/m_1$ of the terms of $\Omega_{\mathrm{re}}$, using
$m_3=m_2m_4/m_1$ and multiplying by $(m_1m_2m_4)^2$ gives $N_{\mathrm{re}}=0$; the terms of
$\Omega_{\mathrm{im}}$ are $i$ times reals with squares $(m_1m_2m_4)^2$, $m_5^2m_1^2m_2^2M/m_4$,
$m_5^2m_1^2m_4^2M/m_2$, which gives $N_{\mathrm{im}}=0$.  Finally $Q>0$.
\end{proof}

\begin{lemma}[computer-assisted]\label{lem:d3-order4}
At every label $G_4=-\frac{\Pi}{m_1^2}\bigl(m_4(m_1+m_2),m_2(m_1+m_4)\bigr)$ with
$\Pi=(m_1D_{12}-m_4D_{34})(m_1D_{14}-m_2D_{23})$.
\end{lemma}

{As in Lemma~\ref{lem:d3-order3}, the computation uses exact arithmetic and reduction
modulo the equations for $D_{kl}$. After solving for $X_2$, it evaluates the preceding formula
for $G_4$ (Appendix~\ref{app:verification}). At every label,}
$(m_1D_{12})^3-(m_4D_{34})^3=m_1(m_1+m_4)(m_1-m_4)/(m_1+m_2)$ and
$(m_1D_{14})^3-(m_2D_{23})^3=m_1(m_1+m_2)(m_1-m_2)/(m_1+m_4)$, so $\Pi=0$ forces $m_1=m_4$ or
$m_1=m_2$; with (29), $m=(p,q,q,p,m_5)$ or $(p,p,q,q,m_5)$.

{\begin{lemma}\label{lem:d3-curves}
Let $m\in\Rp^5$ satisfy (29), $N_{\mathrm{re}}(m)=N_{\mathrm{im}}(m)=0$ and $m_1=m_2$ (resp.\
$m_1=m_4$).  Write $m_1=a^2$, $m_3=b^2$ (resp.\ $m_2=b^2$) with $a,b>0$, and
$c=(2(a^2+b^2))^{1/2}$.  Then
\[
m_5=\frac{a^3b^3}{c\,|a^3+b^3|}\qquad\text{or}\qquad m_5=\frac{a^3b^3}{c\,|a^3-b^3|}\quad(a\ne b).
\]
Conversely $N_{\mathrm{re}}$, $N_{\mathrm{im}}$ and $S$ vanish on these four families, called
$E^\pm_{12}$ (if $m_1=m_2$) and $E^\pm_{14}$ (if $m_1=m_4$).  They meet only at $a=b$, i.e.\ at
$m$ proportional to $(1,1,1,1,\frac14)$, which lies on $E^+_{12}\cap E^+_{14}$.
\end{lemma}}

{\begin{proof}
For either mass pattern, the equations $N_{\mathrm{re}}=0$ and
$N_{\mathrm{im}}=0$ coincide. Their common polynomial factors as
\[
 N_{\mathrm{re}}=N_{\mathrm{im}}
 =a^{12}\bigl(a^6b^6-m_5^2c^2(a^3+b^3)^2\bigr)
          \bigl(a^6b^6-m_5^2c^2(a^3-b^3)^2\bigr).
\]
Positivity of $a,b,c,m_5$ shows that the first factor gives exactly the plus
formula, and the second gives the minus formula, with $a\ne b$.
One of the first two factors of $S$ vanishes, so $S=0$ as well.
The two mass patterns coincide only when all four outer masses are equal.
In that case the minus formula is impossible, while the plus formula gives $m_5=a^2/4$.
For each mass pattern, the plus and minus formulas are distinct when $a,b>0$.
\end{proof}

\begin{lemma}\label{lem:d3-census}
Let $m$ lie on one of the four families.
\begin{alphalist}
  \item If $a\ne b$, exactly $108$ labels have $G_3=G_4=0$: two of the eight sign choices for
  $(D_{13},D_{15},D_{35})$ give $\Omega_{\mathrm{re}}=0$, two of those for $(D_{24},D_{25},D_{45})$
  give $\Omega_{\mathrm{im}}=0$, and $\Pi=0$ exactly when $D_{14}=D_{23}$ (on $E^\pm_{12}$) or
  $D_{12}=D_{34}$ (on $E^\pm_{14}$), for $27$ of the $81$ choices of cube roots.
  \item If $a=b$, exactly $180$ labels have $G_3=G_4=0$: those with $D_{15}=D_{35}=-2D_{13}$,
  $D_{25}=D_{45}=-2D_{24}$, and $D_{12}=D_{34}$ or $D_{14}=D_{23}$ ($45$ choices of cube roots).
\end{alphalist}
\end{lemma}}

{\begin{proof}
Exchanging $2$ and $4$ together with $z$ and $w$ reduces the proof to
$E_{12}^{\epsilon}$, $\epsilon\in\{1,-1\}$.
Write $A_0=a^3$, $B_0=b^3$, $\nu=\operatorname{sgn}(A_0+\epsilon B_0)$,
and set $D_{13}=\sigma$, $D_{15}=s_1c/b$, $D_{35}=s_2c/a$.
After division by a nonzero real factor, $\Omega_{\mathrm{re}}=0$ becomes
\[
 (\sigma\nu+s_1)A_0+(\epsilon\sigma\nu+s_2)B_0=0.
\]
For $\epsilon=1$, any nonzero coefficients have the same sign, so the equation
forces both coefficients to vanish. For $\epsilon=-1$, cancellation between two nonzero
coefficients would require $A_0=B_0$, which is excluded on the minus family.
There are therefore exactly two real sign triples. Applying the same argument
to $D_{24}/i,D_{25}/i,D_{45}/i$ gives exactly two imaginary sign triples.
If $a\ne b$, the first factor of $\Pi$ cannot vanish because the cubes of
$a^2D_{12}$ and $b^2D_{34}$ are $a^4/2$ and $b^4/2$.
Consequently $\Pi=0$ exactly when $D_{14}=D_{23}$, leaving $3^3=27$ cube-root
choices and hence $4\cdot27=108$ labels.

At $a=b$, the sign conditions reduce to
$D_{15}=D_{35}=-2D_{13}$ and $D_{25}=D_{45}=-2D_{24}$, again giving four choices.
Each equality $D_{12}=D_{34}$ or $D_{14}=D_{23}$ admits $27$ cube-root choices,
and their intersection admits $9$. Thus their union gives
$4(27+27-9)=180$ labels; the intersection gives $4\cdot9=36$.
\end{proof}

}\subsection{Orders five and six}

\begin{lemma}[computer-assisted]\label{lem:d3-order5}
Let $m$ lie on one of the four families with $a\ne b$.  Then $G_5\neq0$ at each of the $108$
labels of Lemma~\ref{lem:d3-census}(a).  At $a=b$, $G_5=0$ at all $180$ labels of
Lemma~\ref{lem:d3-census}(b).
\end{lemma}

{The fifth-order calculation gives the following identities. On $E^+_{12}$,
where $D_{14}=D_{23}$, the $z$-component of $G_5$ is a nonzero factor times
$(a-b)L_+(\xi_1,\xi_2)$, where $\xi_1=D_{12}/D_{14}$,
$\xi_2=D_{34}/D_{14}$ and
\[
 L_\epsilon(\xi_1,\xi_2)
 =a^2F_{4,\epsilon}\xi_1-a^{-2}b^4F_{4,\epsilon}\xi_2
       -b^2(a^2-b^2)F_{2,\epsilon},
\]
\[
 F_{4,\epsilon}=a^4+\epsilon a^3b+a^2b^2+\epsilon ab^3+b^4,
 \qquad F_{2,\epsilon}=a^2+\epsilon ab+b^2,
 \qquad \epsilon\in\{1,-1\},
\]
with $\xi_1^3=(a^2+b^2)/(2a^2)$ and
$\xi_2^3=(a^2+b^2)/(2b^2)$.
The product of the nine values of $L_+$ is, up to sign and a nonzero monomial
factor, $(a-b)^3(a+b)^3P_{66}(a,b)$.
On $E^-_{12}$, the calculation gives $G_5=0\Rightarrow P_{66}(a,-b)=0$
at each of the $108$ labels. All factors removed in obtaining this implication
are regular and nonzero for $a,b>0$, $a\ne b$.
These identities are the computer-assisted part of the argument.
Their required nonvanishing follows from the elementary estimates below;
no coefficient-sign or Sturm computation is needed.

We treat both signs together. Interchanging $a,b$ and $\xi_1,\xi_2$ multiplies $L_\epsilon$
by the nonzero factor $-a^2/b^2$. We may therefore assume $a>b>0$ and write
\[
 r=b/a\in(0,1),\quad c_0=\left(\frac{1+r^2}{2}\right)^{1/3},\quad
 f_{4,\epsilon}=1+\epsilon r+r^2+\epsilon r^3+r^4,\quad
 f_{2,\epsilon}=1+\epsilon r+r^2.
\]
Writing $\xi_1=c_0\omega^j$ and $\xi_2=c_0r^{-2/3}\omega^k$,
$\omega=e^{2\pi i/3}$, gives
\[
 a^{-6}L_\epsilon=A_\epsilon\omega^j-B_\epsilon\omega^k-C_\epsilon,
 \quad A_\epsilon=c_0f_{4,\epsilon},\quad
 B_\epsilon=A_\epsilon r^{10/3},\quad
 C_\epsilon=r^2(1-r^2)f_{2,\epsilon}.
\]
All three coefficients are positive, with $A_\epsilon>B_\epsilon$.
The imaginary part can vanish only for $j=k=0$.
Indeed, if both roots have imaginary parts of the same sign, cancellation would require
$A_\epsilon=B_\epsilon$; if their imaginary parts have opposite signs, the two contributions
reinforce each other. Nor can a single nonreal root have its imaginary part cancelled.
It remains to compare the real coefficients. For the plus sign, $c_0>r^2$ and
$f_{4,+}>f_{2,+}$. For the minus sign, we have
\[
 \frac{f_{4,-}}{f_{2,-}}=\frac{1+r^5}{1+r^3}>r^{4/3},\qquad c_0>r^{2/3},
\]
since
\[
 1+r^5-r^{4/3}-r^{13/3}
 =(1-r^{4/3})(1-r^{13/3})+r^5-r^{17/3}>0.
\]
In both cases, therefore, $c_0f_{4,\epsilon}>r^2f_{2,\epsilon}$, and
\[
 A_\epsilon-B_\epsilon
 =c_0f_{4,\epsilon}(1-r^{10/3})
 >r^2f_{2,\epsilon}(1-r^2)=C_\epsilon.
\]
This excludes the remaining real-root case, so all nine values of $L_\epsilon$ are nonzero.
The product identity for $L_+$ gives $P_{66}(a,b)\ne0$.
Replacing $b$ by $-b$ in the same identity gives $P_{66}(a,-b)\ne0$ for $a\ne b$.
The fifth-order identities thus exclude both families when $a\ne b$.

}
{Exchanging bodies $2$ and $4$ together with $z$ and
$w$ maps the chart at $m$ onto the chart at $(m_1,m_4,m_3,m_2,m_5)$, $E^\pm_{12}$ onto
$E^\pm_{14}$ and the $108$ labels onto the $108$ labels; the relabelling $(13)(24)$ exchanges $a$
and $b$, accounting for the minus families with $a<b$.
At $a=b$, the calculation gives zero for both components of $G_5$ at the labels of $E^+_{12}$.
The exchange then gives vanishing at all $180$ labels.}

\begin{lemma}[computer-assisted]\label{lem:d3-order6}
Let $m=(\frac12,\frac12,\frac12,\frac12,\frac18)$.  Of the $180$ labels of
Lemma~\ref{lem:d3-census}(b), the $144$ labels with $D_{12}\ne D_{34}$ or $D_{14}\ne D_{23}$ have
$G_6\neq0$.
\end{lemma}

{The lemma is proved by computing the jets modulo two primes
(Appendix~\ref{app:verification}). All data lie in $R=\mathbb Z[\frac12,i,\omega]$,
where $\omega=e^{2\pi i/3}$. Choose a prime $\mathfrak p$ of $R$ above $p\in\{13,37\}$.
Since $p\equiv1\pmod{12}$, we have $R/\mathfrak p=\mathbb F_p$.

If a germ passed through a label, Lemma~\ref{lem:d3-graph} would give $F(t,X(t))\equiv0$.
At each order, $X_k$ would then be uniquely determined by the order-$k$ equations,
using any $19$ rows with an invertible minor, and all $21$ residuals would vanish.
We use $19$ rows whose minor is a unit modulo $\mathfrak p$.
The calculation produces a nonzero order-$6$ residual at each of the $144$ labels.
Changing $\mathfrak p$ acts on the labels through $\mathrm{Gal}(\mathbb Q(i,\omega)/\mathbb Q)$,
with $\omega\mapsto\omega^2$ and $i\mapsto-i$. This action preserves the set of $144$ labels,
and the two primes give the same excluded set.}

{\begin{lemma}\label{lem:d3-twisted}
Let $m=(\frac12,\frac12,\frac12,\frac12,\frac18)$, $\alpha,\beta\in\{1,\omega,\omega^2\}$ and
$\varepsilon,\eta\in\{\pm1\}$.  Let $(P,Q,Y)$ be the analytic solution near $\tau=0$ of
\[
(1-P)(1+Q)=\alpha\tau^2,\qquad (1+P)(1-Q)=\beta\tau^2,\qquad 4PQY^2+1=0
\]
with $(P,Q,Y)(0)=(1,1,\eta i/2)$; the Jacobian there is $\pm16i$.  Then
\[
z=\tfrac1\tau(1,P,-1,-P,0),\qquad w=\tfrac1\tau(1,-Q,-1,Q,0),
\]
\[
\delta_{12}=\delta_{34}=\alpha,\ \ \delta_{14}=\delta_{23}=\beta,\ \
\delta_{13}=\tfrac{\varepsilon\tau}2,\ \ \delta_{15}=\delta_{35}=-\varepsilon\tau,\ \
\delta_{24}=\tau Y,\ \ \delta_{25}=\delta_{45}=-2\tau Y
\]
solves \eqref{eq:system} with the gauge $z_{13}=w_{13}$.  In the chart, with $t=-\tau/2$, it is an
analytic solution of $F=0$ through the label with $D_{12}=D_{34}=\alpha$, $D_{14}=D_{23}=\beta$,
$D_{13}=-\varepsilon$ and $D_{24}=-\eta i$.  These $36$ labels are the labels of
Lemma~\ref{lem:d3-census}(b) with $D_{12}=D_{34}$ and $D_{14}=D_{23}$.  Along these germs
\begin{equation}\label{eq:d3-identities}
\delta_{15}+2\delta_{13}=\delta_{35}+2\delta_{13}=\delta_{25}+2\delta_{24}=\delta_{45}+2\delta_{24}=0 .
\end{equation}
\end{lemma}}

{\begin{proof}
The derivative of the three defining equations with respect to $(P,Q,Y)$ at
$(\tau,P,Q,Y)=(0,1,1,\eta i/2)$ is
\[
 \begin{pmatrix}-2&0&0\\0&-2&0\\-1&-1&4\eta i\end{pmatrix},
\]
with determinant $16\eta i\ne0$. The implicit function theorem gives the
claimed analytic solution. Since the equations depend on $\tau$ only through
$\tau^2$, uniqueness makes $P,Q,Y$ even, and $P-1,Q-1=O(\tau^2)$.

We verify the body equations first. For either coordinate $q=z$ or $q=w$,
the two adjacent side forces at outer body $k$ sum to $q_k$:
the cubes of their inverse distances are $1$, the adjacent masses are $1/2$,
and the adjacent bodies have opposite positions.
At body $1$ or $3$, the diagonal force and the force from body $5$ cancel,
since
\[
 \tfrac12(\varepsilon\tau/2)^3(2q_k)
       +\tfrac18(-\varepsilon\tau)^3q_k=0.
\]
At body $2$ or $4$ the same cancellation is
\[
 \tfrac12(\tau Y)^3(2q_k)+\tfrac18(-2\tau Y)^3q_k=0.
\]
Pairing opposite bodies also proves the equation for body $5$.
For the distance equations, the pairs $12,34$ reduce to
$\alpha^2(1-P)(1+Q)/\tau^2=1$, the pairs $14,23$ to
$\beta^2(1+P)(1-Q)/\tau^2=1$, and $24,25,45$ to $-4PQY^2=1$;
the remaining distance equations hold identically. The gauge follows from
$z_{13}=w_{13}=-2/\tau$.
Finally $t=-\tau/2$ makes the chart coordinates analytic, since the only
apparent quotients include $(P-1)/\tau^2$ and $(Q-1)/\tau^2$.
The displayed inverse distances give \eqref{eq:d3-identities} identically
and yield precisely the $9\cdot4=36$ labels stated in the lemma.
\end{proof}

}\subsection{Exclusion of curves containing a positive point}

\begin{theorem}\label{thm:d3}
Let $m\in\Rp^5$, let $\Gamma\subset A_m$ be an irreducible curve, and let $\gamma$ be a Diagram-3
germ on $\Gamma$, with the bodies numbered as in Diagram~3.  Then $m_1=m_2=m_3=m_4=4m_5$, and
\eqref{eq:d3-identities} holds identically on $\Gamma$.  In particular $\Gamma$ contains no
positive point, and Theorem~\ref{thm:exclusion}(c) holds.
\end{theorem}

\begin{proof}
{Exchanging $z$ and $w$ and renumbering if necessary puts $\gamma$ in the form
described at the start of this section. These operations preserve the conclusions.
Proposition~\ref{prop:d3-cover} gives a limiting chart point $X_0$ satisfying
\eqref{eq:d3-boundary} and $B-A=1$. Relation (29) holds by Proposition~\ref{prop:d3-first},
so Lemma~\ref{lem:d3-fibre} identifies $X_0$ as a label. Lemma~\ref{lem:d3-graph} then gives
$G\equiv0$ at $X_0$.

The third-order condition, from Lemma~\ref{lem:d3-order3} and
Corollary~\ref{cor:d3-order3}, gives $N_{\mathrm{re}}(m)=N_{\mathrm{im}}(m)=0$.
The fourth-order condition in Lemma~\ref{lem:d3-order4} gives $m_1=m_4$ or $m_1=m_2$.
Thus $m$ belongs to one of the four families in Lemma~\ref{lem:d3-curves}.
Lemmas~\ref{lem:d3-census}(a) and~\ref{lem:d3-order5} exclude every point with $a\ne b$,
leaving $a=b$, or equivalently $m_1=m_2=m_3=m_4=4m_5$.

The covariance of \eqref{eq:system} under
$(m,z,w,\delta)\mapsto(\kappa^3m,\kappa z,\kappa w,\kappa^{-1}\delta)$, $\kappa>0$,
preserves both Diagram-3 germs and \eqref{eq:d3-identities}.
We may therefore normalize to $m=(\frac12,\frac12,\frac12,\frac12,\frac18)$.
Lemmas~\ref{lem:d3-census}(b) and~\ref{lem:d3-order6} now imply $D_{12}=D_{34}$ and
$D_{14}=D_{23}$ at $X_0$.
By Lemmas~\ref{lem:d3-twisted} and~\ref{lem:d3-graph}, the germ $\gamma$ coincides in the chart
with the twisted Roberts germ through $X_0$. Since renormalization leaves $\delta_{kl}$ unchanged,
\eqref{eq:d3-identities} holds along $\gamma$.

Each expression in \eqref{eq:d3-identities} is a regular function on the irreducible curve
$\Gamma$. It vanishes at infinitely many points of $\gamma$, and hence vanishes identically
on $\Gamma$. A positive point would have $\delta_{15}+2\delta_{13}>0$, a contradiction.}
\end{proof}

{Thus the only possible masses for a Diagram-3 germ are $(4,4,4,4,1)$ up to scale.
Up to the preceding symmetries, every such germ is one of the $36$ twisted Roberts germs
constructed in Lemma~\ref{lem:d3-twisted}. The following example explains why positivity
of the point, rather than merely its reality, is needed.}

\begin{remark}[The modified Roberts continuum]\label{rem:d3-roberts}
Let $m=(\frac12,\frac12,\frac12,\frac12,\frac18)$, proportional to $(4,4,4,4,1)$
\cite[Rem.~8, p.~583]{AK}.  The curve
\[
z=(a,b,-a,-b,0),\quad w=(a,-b,-a,b,0),\quad b^2=a^2-1,
\]
\[
\delta_{12}=\delta_{23}=\delta_{34}=\delta_{14}=1,\ \ \delta_{13}=\tfrac1{2a},\ \
\delta_{24}=\tfrac{i}{2b},\ \ \delta_{15}=\delta_{35}=-\tfrac1a,\ \ \delta_{25}=\delta_{45}=-\tfrac ib
\]
is the modified Roberts continuum \cite[p.~543]{AK}.  It solves \eqref{eq:system} with the gauge
$z_{13}=w_{13}$, $U=I=1$ on it, and \eqref{eq:d3-identities} holds, so it has no positive point.
Since $w_{12}/z_{12}=(a+b)^2$ on it, the rotation by $\mu=a+b$ gives an irreducible curve
$\Gamma_R\subset A_m$.  For $0<a<1$ and $b\in i\R$, $|\mu|=1$ and the points are real, with
$\delta_{13}>0>\delta_{15}=\delta_{35}$.
{Before this rotation, the end $a\to\infty$ with $b\sim a$ is precisely the twisted Roberts
germ of Lemma~\ref{lem:d3-twisted} with
$(\alpha,\beta,\varepsilon,\eta)=(1,1,1,1)$, $\tau=1/a$ and $P=Q=b/a$.
In this balanced $z_{13}=w_{13}$ representative its strokes are
$z_{14}$, $z_{23}$, $w_{12}$, $w_{34}$, so the corresponding end of $\Gamma_R$ is a Diagram-3 germ}
(Roberts' continuum at infinity \cite[\S7.2]{AK}); the sides keep length $1$,
$r_{13},r_{24},r_{k5}\to\infty$, $p_{12}=a^2(1-a^2)$ is unbounded, and (29) holds.
{The curve $\Gamma_R$ therefore contains real points, has the positive constant
$U=I=1$, and carries a Diagram-3 germ with an unbounded four-product.
Consequently, curves containing a real point cannot be excluded in Diagram~3.
In Theorem~\ref{thm:d3}, positivity enters only when the identity
$\delta_{15}+2\delta_{13}=0$, valid on all of $\Gamma$, is evaluated at a positive point.
This is the distinction required in Remark~\ref{rem:levels}.}
\end{remark}



\section{The triangle in Diagrams 4 and 5}\label{sec:d45}

{Both Diagrams 4 and 5 contain an isolated $zw$-triangle $345$ with uncircled vertices.
We first extract the mass relation common to these two diagrams.}

\begin{proposition}[Relation (23)]\label{prop:d45-first}
Let $m\in\Rp^5$ and let $\gamma$ be a Diagram 4 or Diagram 5 germ, renormalized as in
Lemma~\ref{lem:germ} so that $\mathcal Z$ and $\mathcal W$ have the same least order $-2\nu<0$.
Let $z_c$ and $w_c$ be the centres of mass of $T=\{3,4,5\}$.  Then the limits
\[
  \zeta_k=\lim_{t\to0}t^{-\nu}(z_k-z_c),\qquad \eta_k=\lim_{t\to0}t^{-\nu}(w_k-w_c),\qquad
  D_{kl}=\lim_{t\to0}t^{\nu}\delta_{kl}\qquad(k,l\in T)
\]
exist and solve the zero-multiplier triangle system of Section~\ref{sec:d1}, and
$\sum_km_k\zeta_k=\sum_km_k\eta_k=0$.  Hence $(\zeta,\eta,D)$ is as in
Lemma~\ref{lem:d1-shape}, and $Q(m_3,m_4,m_5)=0$, that is, the masses satisfy (23).
\end{proposition}

\begin{proof}
{The sides of $T$ are $zw$-edges. Estimate~1 and the left equality case of Estimate~2,
as stated in Fact~\ref{fact:rules} and justified for germs by Lemma~\ref{lem:germ}, give
$\ord z_{kl}=\ord w_{kl}=\nu$ and $\ord\delta_{kl}=-\nu$ for $k,l\in T$.
Since $z_k-z_c$ and $w_k-w_c$ are linear combinations of these differences, the limits exist.
The choice of the centres of mass makes both sums $\sum_km_k\zeta_k$, $\sum_km_k\eta_k$ vanish.

For $k\in T$, consider the equation $z_k=\sum_{l\ne k}m_l\delta_{kl}^3(z_k-z_l)$.
The terms with $l\in T$ have order $-2\nu$. Every other term has larger order:
$z_k$ because $k$ is uncircled, and the terms with $l\notin T$ because $T$ is isolated
and no $z$-stroke leaves it. Multiplication by $t^{2\nu}$ followed by the limit $t\to0$
therefore gives $\phi_k=0$. The same argument gives $\psi_k=0$.
The distance equations within $T$ give $\chi_{kl}=0$.
Lemma~\ref{lem:d1-shape} applies, and for positive masses the relation $Q=0$ is equivalent to (23).}
\end{proof}

\subsection{Diagram 4}\label{sec:d4}

Diagram~4 consists of the isolated $zw$-triangle $345$, whose vertices are not circled, and the
isolated $z$-stroke $12$ with $z$-circles at $1$ and $2$; there are no other strokes or circles
\cite[\S7.4, p.~567, and Fig.~11]{AK}.
For this diagram \cite[\S8.4]{AK} record only relation (23).  We use the notation of
Section~\ref{sec:d1}: $P=\{1,2\}$, $T=\{3,4,5\}$, $i\in P$, $k,l\in T$, $M_P$, $M_T$, $Z_P$, $Z_T$,
$W_P$, $W_T$, the equations $E^z_j,E^w_j$, the triangle functions $\phi_k,\psi_k,\chi_{kl}$, the
torque $T_P$ of Lemma~\ref{lem:d1-syzygies} and the orbit $O$ of Lemma~\ref{lem:d1-mb}.
Renumbering inside $P$ or inside $T$ does not change the diagram.

\begin{theorem}\label{thm:d4}
Let $m\in\Rp^5$.  No curve in $A_m$ carries a Diagram~4 germ.
\end{theorem}

{This proves part~(a) of Theorem~\ref{thm:exclusion} for $r=4$ without any condition
on the four-products. The proof has two steps.

First, we show that the chart contains every germ. A torque identity rules out a triangle
body at the triangle's $w$-centre (Lemma~\ref{lem:d4-central}).
Every germ then has a limit in the boundary fibre $\mathcal B_4$,
which is smooth near that limit (Proposition~\ref{prop:d4-cover} and Lemma~\ref{lem:d4-fibre}).

Second, we derive two necessary balances of the leading forces.
The Lyapunov--Schmidt reduction in Lemma~\ref{lem:d4-ls} expresses them as successive
obstructions to the existence of a germ. The first obstruction requires equal leading
$w$-forces on the two outer bodies. The second vanishes identically, and the third requires
equal leading $z$-forces (Lemma~\ref{lem:d4-obstr}).
{Positivity makes these two balances incompatible}
(Lemma~\ref{lem:d4-census}).}

\subsubsection{The chart}\label{d4-chart}
{The invariance of \eqref{eq:system}, apart from the gauge, under
$(z,w,\delta)\mapsto(cz,c^{-1}w,\delta)$, $c\in\C^*$ allows us to use
$z_{34}=w_{34}$ in place of $z_{12}=w_{12}$. With parameter $t$, write}
\[
  (\zeta_3,\zeta_4,\zeta_5)=(m_5u,\ m_5v,\ -m_3u-m_4v),\qquad
  (\eta_3,\eta_4,\eta_5)=(m_5p,\ m_5q,\ -m_3p-m_4q),
\]
so that $\sum_km_k\zeta_k=\sum_km_k\eta_k=0$, and
\begin{equation}\label{eq:d4-chart}
\begin{gathered}
  z_i=t^{-4}A_i,\qquad w_i=t^{4}B_i\qquad(i\in P),\\
  z_k=t^{-1}\hat Z+t^{2}\zeta_k,\qquad w_k=t^{5}\hat W+t^{2}\eta_k\qquad(k\in T),\\
  \delta_{12}=D_{12},\qquad\delta_{ik}=t\,D_{ik},\qquad\delta_{kl}=t^{-2}D_{kl}.
\end{gathered}
\end{equation}
{The chart has twenty coordinates $x$: $A_1,A_2,B_1,B_2,\hat Z,\hat W,u,v,p,q$
and the ten $D_{jl}$. The substitution is invertible for $t\ne0$.
For convenience, set}
\[
\begin{gathered}
  \Sigma_i=\sum_km_kD_{ik}^3,\qquad f_i=\sum_km_kD_{ik}^3\eta_k,\\
  R^z_i=A_i-m_{i'}D_{12}^3(A_i-A_{i'}),\qquad R^w_i=B_i-m_{i'}D_{12}^3(B_i-B_{i'}),
\end{gathered}
\]
where $\{i,i'\}=P$.  Substituting \eqref{eq:d4-chart} gives
\begin{equation}\label{eq:d4-rows}
\begin{aligned}
  t^{4}E^z_i&=R^z_i-\textstyle\sum_km_kD_{ik}^3\,(t^3A_i-t^6\hat Z-t^9\zeta_k),\\
  t^{-4}E^w_i&=R^w_i+t\,f_i-\textstyle\sum_km_kD_{ik}^3\,(t^3B_i-t^4\hat W),\\
  t^{4}E^z_k&=-\phi_k+t^3\hat Z+t^6\zeta_k-\textstyle\sum_im_iD_{ik}^3\,(t^6\hat Z+t^9\zeta_k-t^3A_i),\\
  t^{4}E^w_k&=-\psi_k+t^6\eta_k+t^9\hat W-\textstyle\sum_im_iD_{ik}^3\,(t^9\eta_k+t^{12}\hat W
  -t^{11}B_i),
\end{aligned}
\end{equation}
{The distance equations have left-hand sides
$D_{12}^2(A_1-A_2)(B_1-B_2)-1$, the \emph{crossed ties}
$D_{ik}^2(t^3\hat Z+t^6\zeta_k-A_i)(\eta_k-t^2B_i+t^3\hat W)-1$, and $\chi_{kl}$.
The gauge becomes $t^{-2}(z_{34}-w_{34})=(\zeta_4-\zeta_3)-(\eta_4-\eta_3)$.
Together these give $21$ polynomials, denoted by $F(t,x)=\sum_{j\ge0}t^jF_j(x)$.
Their zero set at the boundary is the \emph{boundary fibre} $\mathcal B_4=\mathcal B_4(m)$,
the zero set of $F_0$. Equation~\eqref{eq:d4-rows} shows that $F_1$ vanishes except in the two
rows $t^{-4}E^w_i$, where it equals $f_i$. Likewise, $F_2$ vanishes except in the crossed
distance equations, where it equals $D_{ik}^2A_iB_i$.}

\subsubsection{Representation of germs in the chart}\label{d4-cover}
{Let $\gamma$ be a Diagram~4 germ at $m\in\Rp^5$.
Normalize it as in Lemma~\ref{lem:germ}, so that $Z\approx W\approx\epsilon^{-2}$.
After a base change of the branch parameter $\tau$, all coordinates are Laurent series
in $\tau$, with $\ord_\tau Z=\ord_\tau W=-4q$ for an integer $q\ge1$.
Set $t=\tau^q$ and measure orders in units of $\ord t$.
The maximal components of $\mathcal Z$ and $\mathcal W$ then have order $-4$,
and $\ord\epsilon=2$.}

\begin{lemma}\label{lem:d4-val}
Along $\gamma$, for $i\in P$ and $k\ne l$ in $T$:
\begin{romanlist}
  \item $\ord z_k>-4$ and $\ord z_{ik}=-4$;
  \item $\ord z_1=\ord z_2=\ord z_{12}=-4$, $\ord\delta_{12}=0$, $\ord w_{12}=4$,
  $\ord w_1,\ord w_2\ge4$, and hence $\ord W_T=\ord W_P\ge4$;
  \item $\ord z_{kl}=\ord w_{kl}=2$ and $\ord\delta_{kl}=-2$.
\end{romanlist}
\end{lemma}

\begin{proof}
{Part (i) follows because the bodies in $T$ have no $z$-circles, whereas $1$ and $2$ do.

For (ii), the $z$-circles at bodies $1,2$ give $\ord z_1=\ord z_2=-4$.
The isolated component $\{1,2\}$ of the $z$-diagram has circled vertices, so Rule~1e
\cite[p.~546]{AK} makes its stroke $12$ maximal. The right equality case of Estimate~2
\cite[p.~548]{AK} gives $r_{12}\approx1$, $z_{12}\approx\epsilon^{-2}$ and
$w_{12}\approx\epsilon^2$. The triangle edges are not maximal: part (iii), which is independent
of (ii), gives $z_{kl}\approx\epsilon$ on them. Thus $12$ is the only maximal $z$-edge.
Estimate~1 \cite[p.~548]{AK}, with $z$ and $w$ exchanged, now gives $w_{12}\prec w_{ik}$
for $i\in P$, $k\in T$.

Apply \cite[Prop.~1, p.~559]{AK}, with $z$ and $w$ exchanged, to bodies $1$ and $2$,
as in the proof of \cite[Prop.~4, p.~564]{AK}. It places the origin in the $w$-cluster
formed by $w_1$ and $w_2$, whose size is $w_{12}$. Hence $w_1,w_2\preccurlyeq w_{12}$.
The relation $M_PW_P+M_TW_T=0$ gives the assertion about the two centres.

For (iii), the triangle edges are $zw$-edges. Estimate~1 and the left equality case
of Estimate~2 give $r_{kl}\approx z_{kl}\approx w_{kl}\approx\epsilon$.}
\end{proof}

{\begin{lemma}\label{lem:d4-torque}
Let $0<e<2$.  In \eqref{eq:d4-chart} replace $\hat Z$ by $t^{-3e/2}\hat Z$, $\hat W$ by
$t^{-e/2}\hat W$, $(\eta_3,\eta_4,\eta_5)$ by $(t^e\omega,\ m_5q,\ -m_3m_5^{-1}t^e\omega-m_4q)$ and
$D_{i3}$ by $t^{-e/2}d_i$ $(i\in P)$.  Then
\[
  T_P=t^{1-e/2}\Big(m_3\sum_{i\in P}m_id_i\,\big(d_i^2A_i\omega\big)+\sum_{g\in G}c_g\,t^g\Big),
\]
where the $c_g$ are polynomials in the chart coordinates and
$G=\{\tfrac12e,\ \tfrac32e,\ 3,\ 3-\tfrac32e,\ 5-\tfrac52e,\ 5-e,\ 8-e,\ 8+\tfrac12e\}$; all
elements of $G$ are positive.
\end{lemma}
\begin{proof}
After the stated substitutions,
\[
 z_iw_k-w_iz_k
 =t^{-2}A_i\eta_k+t^{1-e/2}A_i\hat W
       -t^{3-3e/2}B_i\hat Z-t^6B_i\zeta_k.
\]
For $k=3$, we have $\eta_3=t^e\omega$ and
$\delta_{i3}^3=t^{3-3e/2}d_i^3$.  The four terms, divided by
$t^{1-e/2}$, have exponents
$0$, $3-3e/2$, $5-5e/2$ and $8-e$, and the first coefficient is
$d_i^3A_i\omega$.
For $k=4,5$, multiply by $t^3D_{ik}^3$ instead.  The constant and
$t^e$ parts of $\eta_k$ give the exponents $e/2$ and $3e/2$;
the last three terms give $3$, $5-e$ and $8+e/2$.
These exponents give the eight elements of $G$, all of which are positive
because $0<e<2$.
\end{proof}

}\begin{lemma}\label{lem:d4-central}
For $k\in T$ let $o_k$ be the coefficient of $t^2$ in $w_k-W_T$ along $\gamma$.  Then
$o_3,o_4,o_5\ne0$.
\end{lemma}

\begin{proof}
{Lemma~\ref{lem:d4-val}(iii) identifies $o_l-o_k$ as the leading coefficient
of $w_{kl}$. Thus at most one}
{$o_k$ vanishes.  Renumbering $T$, suppose that $o_3=0$.  Write
$\ord(w_3-W_T)=2+e$ with $e>0$ when $w_3-W_T\not\equiv0$, and put $e=+\infty$ otherwise.}
By Lemma~\ref{lem:d4-val}(ii), $\ord(W_T-w_i)\ge4$.

{There are two cases. If $e\ge2$, then $\ord w_{i3}\ge4$, so
$\ord\delta_{i3}=-\frac12(\ord z_{i3}+\ord w_{i3})\le0$ and
$\ord Z_{i3}=\ord z_{i3}+3\ord\delta_{i3}\le-4$. This would make $Z_{i3}$ a maximal component
of $\mathcal Z$, producing a $z$-stroke between $i$ and $3$, contrary to Diagram~4.

Suppose instead that $0<e<2$. We have $\ord w_{i3}=2+e$ and $\ord\delta_{i3}=1-\frac e2$,
whereas $\ord w_{ik}=2$ and $\ord\delta_{ik}=1$ for $k=4,5$.
The first two identities of Lemma~\ref{lem:d1-syzygies} express $M_PW_P$ and $M_TZ_T$
as sums of crossed terms $m_im_k\delta_{ik}^3(w_i-w_k)$ and $m_im_k\delta_{ik}^3(z_k-z_i)$.
These have orders at least $5-\frac e2$ and $-1-\frac32e$, respectively.
Together with $M_PW_P+M_TW_T=0$, this gives $\ord W_T\ge5-\frac e2$ and
$\ord Z_T\ge-1-\frac32e$.

The modified chart in Lemma~\ref{lem:d4-torque} therefore represents $\gamma$ with all coordinates
of nonnegative order, and with $A_i,\omega,d_i$ of order $0$.
At $t=0$, the distance equations $\delta_{i3}^2z_{i3}w_{i3}=1$ give $d_i^2A_i\omega=-1$,
while $m_1A_1+m_2A_2=t^4(m_1z_1+m_2z_2)=-t^4M_TZ_T$ vanishes.
Lemma~\ref{lem:d4-torque} now shows that $t^{e/2-1}T_P$ tends to $-m_3(m_1d_1+m_2d_2)$.
This limit is nonzero, since}
\[
  (m_1d_1)^2-(m_2d_2)^2=-\frac{m_1^2}{A_1\omega}+\frac{m_2^2}{A_2\omega}
  =-\frac{m_1^3+m_2^3}{m_1A_1\omega}\ne0 .
\]
{This contradicts $T_P=0$ on $A_m$, proved in Lemma~\ref{lem:d1-syzygies}.}
\end{proof}

\begin{proposition}[Coverage]\label{prop:d4-cover}
Let $m\in\Rp^5$ and let $\gamma$ be a Diagram~4 germ.  Then $m$ satisfies (23), and
$\ord Z_T\ge-1$, $\ord W_T\ge5$ along $\gamma$.  Renormalize by $(z,w)\mapsto(cz,c^{-1}w)$ with
$c=(w_{34}/z_{34})^{1/2}$, so that $z_{34}=w_{34}$.  After a base change of $\tau$, every chart
coordinate of \eqref{eq:d4-chart} is a power series in $\tau$; the value $y$ of $x$ at $\tau=0$
lies in $\mathcal B_4(m)$, and $\eta_k(y)\ne0$ for $k\in T$.
\end{proposition}

\begin{proof}
Relation (23) is Proposition~\ref{prop:d45-first}.  By Lemmas~\ref{lem:d4-val}
and~\ref{lem:d4-central}, $\ord w_{ik}=2$ for all $i\in P$, $k\in T$, hence $\ord\delta_{ik}=1$.
The first two identities of Lemma~\ref{lem:d1-syzygies} then give $\ord Z_T\ge3-4=-1$ and
$\ord W_P\ge3+2=5$, and $\ord W_T=\ord W_P$.  The factor $c$ has order $0$ by
Lemma~\ref{lem:d4-val}(iii), so all orders persist after the renormalization.  Inverting
\eqref{eq:d4-chart}, $A_i=t^4z_i$, $B_i=t^{-4}w_i$, $\hat Z=tZ_T$, $\hat W=t^{-5}W_T$,
$\zeta_k=t^{-2}(z_k-Z_T)$, $\eta_k=t^{-2}(w_k-W_T)$, $D_{12}=\delta_{12}$, $D_{ik}=t^{-1}\delta_{ik}$
and $D_{kl}=t^2\delta_{kl}$; by Lemma~\ref{lem:d4-val} all have nonnegative order, and
$\eta_k\to o_k\ne0$.  The rows of $F$ are monomial multiples of the equations, so they vanish
along $\gamma$, and $F_0(y)=0$.
\end{proof}

\subsubsection{The boundary fibre}\label{d4-fibre}
On $\C^{20}$ consider the scalings
\[
\begin{gathered}
  \alpha\cdot x:\ (A,B,D_{ik})\mapsto(\alpha^2A,\ \alpha^{-2}B,\ \alpha^{-1}D_{ik}),\\
  \beta\cdot x:\ (\zeta,\eta,D_{kl},D_{ik})\mapsto(\beta^2\zeta,\ \beta^2\eta,\ \beta^{-2}D_{kl},\
  \beta^{-1}D_{ik}),
\end{gathered}
\]
{Here $\alpha,\beta\in\C^*$, and coordinates not displayed remain fixed.
Together with translations of $(\hat Z,\hat W)$, these transformations preserve the zero set
of $F_0$: each row changes only by a nonzero factor.
They therefore act on $\mathcal B_4(m)$, although they are not symmetries of $F$.}

\begin{lemma}\label{lem:d4-fibre}
Let $m\in\Rp^5$ satisfy (23) and let $y\in\mathcal B_4(m)$.
\begin{alphalist}
  \item $D_{12}^3=M_P^{-1}$, $A_2=\rho A_1$ and $B_2=\rho B_1$ with $\rho=-m_1/m_2$,
  $D_{12}^2(A_1-A_2)(B_1-B_2)=1$ and $D_{ik}^2A_i\eta_k=-1$; in particular $A_i,B_i,\eta_k\ne0$.
  The point $(\zeta_3,\zeta_4,\eta_3,\eta_4,D_{34},D_{35},D_{45})$ solves the zero-multiplier
  triangle system, so, after renumbering $T$, it lies on the orbit $O$ of Lemma~\ref{lem:d1-mb}.
  \item $dF_0(y)$ has rank $16$.  Its kernel is spanned by $\partial_{\hat Z}$, $\partial_{\hat W}$
  {and the generators of the two scalings.  The annihilator of its image, equivalently the
dual of its cokernel, is spanned by the linear form}
  \[
    v_P(r)=(B_1-B_2)\,(r^z_1-r^z_2)-(A_1-A_2)\,(r^w_1-r^w_2),
  \]
  where $r^z_i,r^w_i$ are the components of $r\in\C^{21}$ in the rows $t^4E^z_i$, $t^{-4}E^w_i$, and
  by four linear forms that involve only the rows $t^4E^z_k$, $t^4E^w_k$, $\chi_{kl}$ and the gauge.
  \item Near $y$, $\mathcal B_4(m)$ coincides with the orbit of $y$ under the scalings and
  the translations; it is a complex submanifold of dimension $4$ with tangent space
  $\ker dF_0(y)$ at $y$.
\end{alphalist}
\end{lemma}

\begin{proof}
{For (a), the identities $R^z_1-R^z_2=(A_1-A_2)(1-M_PD_{12}^3)$ and
$m_1R^z_1+m_2R^z_2=m_1A_1+m_2A_2$ hold everywhere, as do the corresponding identities for
$R^w$. The distance equation for $12$ gives $A_1\ne A_2$ and $B_1\ne B_2$.
Thus $F_0(y)=0$ implies $D_{12}^3=M_P^{-1}$, $m_1A_1+m_2A_2=0$ and $m_1B_1+m_2B_2=0$.
The crossed distance equations are $D_{ik}^2A_i\eta_k=-1$.
The triangle equations in $F_0$ are $\phi_k,\psi_k,\chi_{kl}$ together with the gauge.
The proof of Proposition~\ref{prop:d1-first} gives the last assertion.

For (b), order the coordinates as $(D_{ik};\,A,B,D_{12};\,u,v,p,q,D_{kl};\,\hat Z,\hat W)$.
Order the rows as the crossed distance equations, the pair equations $R^z_i,R^w_i$
and the distance equation for $12$, followed by the triangle equations.
The pair and triangle equations involve only their respective coordinates.
Each $D_{ik}$ occurs only in its own crossed distance equation, where
$\partial/\partial D_{ik}=-2D_{ik}A_i\eta_k\ne0$.}
{So $dF_0(y)$ is block upper triangular, and its rank is $6+\operatorname{rank}J_P+
\operatorname{rank}J_T$}, {where $J_P$, $J_T$ are the pair and triangle blocks.
We compute their ranks separately.

At $y$, the}
differentials of $R^z_1-R^z_2$ and $R^w_1-R^w_2$ are $-3M_PD_{12}^2$ times
$(A_1-A_2)\,dD_{12}$ and $(B_1-B_2)\,dD_{12}$, those of $m_1R^z_1+m_2R^z_2$ and $m_1R^w_1+m_2R^w_2$
are $m_1dA_1+m_2dA_2$ and $m_1dB_1+m_2dB_2$, and the tie of $12$ adds $D_{12}^2$ times
$(B_1-B_2)(dA_1-dA_2)+(A_1-A_2)(dB_1-dB_2)$ modulo $dD_{12}$.
{These relations give $\operatorname{rank}J_P=4$, with kernel spanned by
$(A,-B,0)$. The linear form $v_P$, obtained by differentiating at $y$ the identity}
{$(B_1-B_2)(R^z_1-R^z_2)-(A_1-A_2)(R^w_1-R^w_2)=0$, spans the annihilator of its image.}

{By Lemma~\ref{lem:d1-mb}, the rows $\phi_3,\phi_4,\psi_3,\psi_4,\chi_{kl}$ have rank $5$
at $y$, and their kernel is spanned by the generators $(\zeta,0,-\frac12D)$ and
$(0,\eta,-\frac12D)$ of $\mathbb T$.
The rows $\phi_5$ and $\psi_5$ are redundant because $\sum_km_k\phi_k=\sum_km_k\psi_k=0$.
On the first generator, the gauge row takes the value $\zeta_4-\zeta_3$,
which is nonzero by $\chi_{34}=0$; it takes the value $0$ on $(\zeta,\eta,-D)$.
It follows that $\operatorname{rank}J_T=6$, with kernel spanned by $(\zeta,\eta,-D)$.

The total rank is therefore $16$. Completing the four stated vectors in their $D_{ik}$
components by the crossed distance equations gives a basis of the kernel.
Any linear form annihilating the image has zero coefficients on the crossed distance rows,
as is seen from the columns $D_{ik}$. It is consequently the}
{sum of a form annihilating the image of $J_P$ and one annihilating the image of $J_T$;
the latter annihilator has dimension $10-6=4$.}
{For (c), the differential of the orbit map is onto $\ker dF_0(y)$.
The orbit is thus a submanifold of dimension $4$ near $y$, contained in $\mathcal B_4(m)$.
Conversely, (b) and the rank theorem place $\mathcal B_4(m)$ near $y$ inside a submanifold
of dimension $20-16=4$. The two submanifolds therefore coincide locally.}
\end{proof}

\subsubsection{A Lyapunov--Schmidt lemma}\label{d4-ls}
{The next lemma makes precise why a nonzero obstruction excludes even a ramified
germ. For a polynomial curve $x(t)$, denote by $[F]_j$ the coefficient of $t^j$ in
$F(t,x(t))$.}

\begin{lemma}[Lyapunov--Schmidt]\label{lem:d4-ls}
Let $F(t,x)=\sum_{j\ge0}t^jF_j(x)$ be a polynomial map $\C\times\C^n\to\C^p$, let
$Z=F_0^{-1}(0)$, and let $y\in Z$ be a point near which $Z$ is a complex submanifold with
$T_{y'}Z=\ker dF_0(y')$ for all $y'\in Z$.  For $y'\in Z$ let $Q_{y'}$ be the cokernel of
$dF_0(y')$.  Let $N\ge1$ and suppose that every $y'\in Z$ near $y$ is the value at $t=0$ of a
polynomial curve $x(t)$ with $[F]_j=0$ for $j<N$.
\begin{alphalist}
  \item The class $\Omega_N(y')\in Q_{y'}$ of $[F]_N$ does not depend on the choice of such a curve.
  \item If $\Omega_N(y)\ne0$, no holomorphic germ $c\colon(\C,0)\to(F^{-1}(0),(0,y))$ has
  $t\circ c\not\equiv0$.
\end{alphalist}
For $N=1$ the hypothesis is empty and $\Omega_1(y')$ is the class of $F_1(y')$.
\end{lemma}

\begin{proof}
{Let $K=\ker dF_0(y)$ and $R=dF_0(y)(\C^n)$.
Choose complements $\C^n=K\oplus K'$ and $\C^p=R\oplus R'$, and let $\pi$ project onto $R$
along $R'$. The map $\pi\,dF_0(y)$ sends $K'$ isomorphically onto $R$.
The implicit function theorem therefore gives a holomorphic function $h(t,k)\in K'$,
defined near $(0,0)\in\C\times K$, with $h(0,0)=0$.
Near $(0,y)$, the equation $\pi F(t,y+k+h')=0$, for small $h'\in K'$, holds precisely when
$h'=h(t,k)$. This uniqueness also holds for formal power series in $t$.

The remaining equations are $G(t,k)=F(t,y+k+h(t,k))\in R'$.
Thus, near $(0,y)$, we have $F^{-1}(0)=\{(t,y+k+h(t,k)):G(t,k)=0\}$.
At the boundary, the map $k\mapsto y+k+h(0,k)$ parametrizes $\{\pi F_0=0\}$.
This submanifold contains $Z$ and has the same dimension $\dim K=\dim Z$;
it therefore coincides locally with $Z$. In particular, $G(0,k)=0$.
Write $G=\sum_{j\ge1}t^jG_j(k)$.
For every $y'\in Z$ near $y$, the space $R'$ remains a complement of $dF_0(y')(\C^n)$,
and hence $R'\to Q_{y'}$ is an isomorphism.

We now compare coefficients along a polynomial curve.
Fix $j\le N$ and $y'=y+k'+h(0,k')$, and suppose that $G_1,\dots,G_{j-1}$ vanish near $k'$.
Let $x(t)=y+\kappa(t)+\eta(t)$, with $\kappa(t)\in K$ and $\eta(t)\in K'$,
be a polynomial curve satisfying $x(0)=y'$ and $[F]_i=0$ for $i<j$.
Since $\pi F(t,x(t))=O(t^j)$, formal uniqueness gives
$\eta(t)-h(t,\kappa(t))=e\,t^j+O(t^{j+1})$ for some $e$.
Taylor's formula yields $F(t,x(t))=F(t,y+\kappa+h(t,\kappa))+dF_0(y')(e)\,t^j+O(t^{j+1})$,
while $F(t,y+\kappa+h(t,\kappa))=G(t,\kappa(t))=t^jG_j(k')+O(t^{j+1})$.
Consequently, $[F]_j\equiv G_j(k')$ modulo $dF_0(y')(\C^n)$.

For $j<N$, the curves in the hypothesis give $G_j(k')\in R'\cap dF_0(y')(\C^n)=0$
for all $k'$ near $0$. Induction shows that $G_1,\dots,G_{N-1}$ vanish near $0$.
The case $j=N$ proves (a): $\Omega_N(y')$ is the image of $G_N(k')$.

If $\Omega_N(y)\ne0$, then $G=t^NH$ for a holomorphic $H$ with $H(0,0)=G_N(0)\ne0$.
Thus, near $(0,0)$, the equation $G=0$ forces $t=0$.
Every germ $c$ in (b) has the form $(t,y+k+h(t,k))$ and satisfies $G(t,k)=0$,
so $t\circ c\equiv0$.}
\end{proof}

\subsubsection{The first three obstructions}\label{d4-obstructions}
{For $y\in\mathcal B_4(m)$, define $S_i=\sum_km_kD_{ik}$ and}
\[
  \Theta=\sum_km_k\,\frac{D_{1k}-D_{2k}}{\eta_k}=A_2\Sigma_2-A_1\Sigma_1 ;
\]
{The second expression, as well as $f_i=-S_i/A_i$, follows from
$1/\eta_k=-A_iD_{ik}^2$. Under the boundary scalings, $S_i/A_i$ acquires the factors
$\alpha^{-3}$ and $\beta^{-1}$, while $\Theta$ acquires $\alpha^{-1}$ and $\beta^{-3}$.
Neither expression involves $\hat Z$, $\hat W$, $D_{12}$ or $m_1,m_2$.

These quantities describe the leading forces from the triangle.
On body $i$, the crossed forces are $\sum_km_k\delta_{ik}^3(z_i-z_k)=t^{-1}A_i\Sigma_i+O(t^2)$
and $\sum_km_k\delta_{ik}^3(w_i-w_k)=-t^5f_i+O(t^7)$.
Thus $f_1=f_2$ and $\Theta=0$ express equality of the leading forces on the two outer bodies
in the $w$- and $z$-coordinates, respectively.}

\begin{lemma}\label{lem:d4-obstr}
Let $m\in\Rp^5$ satisfy (23), let $y\in\mathcal B_4(m)$, and let $\Omega_N$ be the classes of
Lemma~\ref{lem:d4-ls} for $F$ and $Z=\mathcal B_4(m)$.
{We regard $v_P$ as a linear functional on the cokernel $Q_y$.}
\begin{alphalist}
  \item $\Omega_1(y)=0$ if and only if $f_1=f_2$, that is, $S_1/A_1=S_2/A_2$.
  \item If $f_1=f_2$ at $y$, then the hypothesis of Lemma~\ref{lem:d4-ls} holds at $y$ for
  $N=3$, and
  \[
    v_P\big(\Omega_3(y)\big)=\tfrac32\,(A_1-A_2)\,\frac{B_1}{A_1}\,\Theta .
  \]
\end{alphalist}
\end{lemma}

\begin{proof}
{For (a), Lemma~\ref{lem:d4-fibre}(b) shows that $\Omega_1(y)$, the class of $F_1(y)$,
vanishes precisely when $F_1(y)$ is annihilated by $v_P$ and the four triangle forms.
The triangle components of $F_1$ vanish, while
$v_P(F_1(y))=-(A_1-A_2)(f_1-f_2)$.

For (b), homogeneity implies $f_1=f_2$ throughout the orbit of $y$.
By Lemma~\ref{lem:d4-fibre}(c), this orbit is a neighbourhood of $y$ in $\mathcal B_4(m)$.
Choose a point $y'$ in it and $N\in\{2,3\}$.
Consider a polynomial curve $x(t)$ with $x(0)=y'$ and $F(t,x(t))=O(t^N)$;
part (a) supplies such a curve for $N=2$.
Write $A_i(t),\dots$ for its components and $A_i,\dots$ for their values at $y'$,
and set $\Phi(t)=1-M_PD_{12}(t)^3$.

We first examine the pair equations. Equation~\eqref{eq:d4-rows} gives
$t^4E^z_i=R^z_i+O(t^3)$. Since $R^z_1-R^z_2=(A_1(t)-A_2(t))\Phi(t)$ and
$m_1R^z_1+m_2R^z_2=m_1A_1(t)+m_2A_2(t)$, we obtain $\Phi=O(t^N)$ and
$A_i(t)=A_i(1+a(t))+O(t^N)$ for a common series $a(t)=O(t)$.

Next, the triangle rows of $F$ agree with those of $F_0$ up to $O(t^3)$.
The triangle components of $x(t)$ therefore solve the triangle block of $F_0$ modulo $t^N$.
The proof of Lemma~\ref{lem:d4-fibre}(b),(c) identifies the zero set of that block near $y'$
with the curve $\beta\mapsto\beta\cdot y'$, whose dimension is $\dim\ker J_T$.
Consequently, the triangle components of $x(t)$ agree modulo $t^N$ with a curve on this orbit.
In particular, $\eta_k(t)=\theta(t)\eta_k+O(t^N)$, where $\theta(0)=1$.
For $N=2$, the triangle components of $[F]_2$ therefore belong to the image of $J_T$.

The crossed distance equations give $D_{ik}(t)^2=-1/\big(A_i(t)(\eta_k(t)-t^2B_i)\big)+O(t^N)$,
and hence}
\[
  D_{ik}(t)^3\eta_k(t)=-\frac{D_{ik}(t)}{A_i(t)}\Big(1+t^2\frac{B_i}{\eta_k}\Big)+O(t^N),\qquad
  D_{ik}(t)=\frac{D_{ik}}{\sqrt{(1+a)\theta}}\Big(1+\frac{t^2B_i}{2\eta_k}\Big)+O(t^N),
\]
and, with $f_i(t)=\sum_km_kD_{ik}(t)^3\eta_k(t)$,
\[
  f_i(t)=-\frac{S_i/A_i}{(1+a)^{3/2}\theta^{1/2}}-\frac32\,t^2\,\frac{B_i}{A_i}\sum_km_k\frac{D_{ik}}{\eta_k}
  +O(t^N).
\]
Since $S_1/A_1=S_2/A_2$ and $B_1/A_1=B_2/A_2$, $f_1(t)-f_2(t)=-\frac32\,t^2\,\frac{B_1}{A_1}\,\Theta+O(t^N)$.

{We can now evaluate the obstruction.
Applying $v_P$ to $F(t,x(t))$, with its coefficients fixed at $y'$, gives, by \eqref{eq:d4-rows},}
\[
  \Phi(t)\big[(B_1-B_2)(A_1(t)-A_2(t))-(A_1-A_2)(B_1(t)-B_2(t))\big]
  -(A_1-A_2)\,t\,\big(f_1(t)-f_2(t)\big)-t^3\,{\Delta_4}+O(t^4),
\]
with ${\Delta_4}=(B_1-B_2)(A_1\Sigma_1-A_2\Sigma_2)-(A_1-A_2)(B_1\Sigma_1-B_2\Sigma_2)$, which vanishes
because $A_2=\rho A_1$ and $B_2=\rho B_1$.  The bracket vanishes at $t=0$ and $\Phi=O(t^N)$.

{For $N=2$, this gives $v_P([F]_2)=0$.
Every linear form annihilating the image of $dF_0(y')$ has zero coefficients on the
crossed distance rows. Together with the preceding result for the triangle part,
this gives $\Omega_2(y')=0$. We can therefore correct the curve at order $2$ to obtain}
{a curve with $F(t,x(t))=O(t^3)$.
This proves the hypothesis of the Lyapunov--Schmidt lemma for $N=3$.
For $N=3$, the displayed expression gives
$v_P([F]_3)=-(A_1-A_2)\cdot\big(-\frac32\frac{B_1}{A_1}\Theta\big)$, as claimed.}
\end{proof}

{\subsubsection{Incompatibility of the two force balances}\label{d4-census}
After renumbering $T$, write
$(\sqrt{m_3},\sqrt{m_4},\sqrt{m_5})=\lambda(bc,b(b+c),c(b+c))$
with $\lambda,b,c>0$. Exchanging $4$ and $5$ if necessary allows us to assume
$b>c$. The case $b=c$ would give a zero triangle offset, contrary to
Proposition~\ref{prop:d4-cover}. The triangle offsets are proportional to
\[
 (o_3,o_4,o_5)=\Bigl(-\frac{b^2-c^2}{bc},
 \frac{c(2b+c)}{b(b+c)},-\frac{b(b+2c)}{c(b+c)}\Bigr).
\]
Put $a=\sqrt{-o_3/o_4}>0$, $d=\sqrt{o_3/o_5}>0$, and
$(g_3,g_4,g_5)=(1,ia,d)$.  If $x=b/c>1$, then
\[
 a^2=\frac{(x-1)(x+1)^2}{2x+1},\qquad
 d^2=\frac{(x-1)(x+1)^2}{x^2(x+2)}.
\]
In particular $0<d<a$ and $d<1$: the differences of the relevant
denominators are
$x^2(x+2)-(2x+1)=(x-1)(x^2+3x+1)>0$ and
$x^2(x+2)-(x-1)(x+1)^2=x^2+x+1>0$.
The mass-centre relation $\sum_km_ko_k=0$ becomes
\begin{equation}\label{eq:d4-positive-centre}
 \frac{m_4}{a^2}=m_3+\frac{m_5}{d^2}.
\end{equation}
For $s=(s_{13},s_{14},s_{15},s_{23},s_{24},s_{25})\in\{1,-1\}^6$
and $j=1,3$, define
\[
 X_j=\sum_{k\in T}s_{1k}m_kg_k^j,\qquad
 Y_j=\sum_{k\in T}s_{2k}m_kg_k^j.
\]

\begin{lemma}\label{lem:d4-census}
For every sign vector $s$ and every $\nu>0$, the two equations
\[
 \nu^3X_1+iY_1=0,\qquad \nu X_3=iY_3
\]
cannot hold simultaneously.
\end{lemma}

\begin{proof}
Put $r=m_3$, $u=m_4a$, $v=m_5d$ and
$\varepsilon=s_{13}s_{15}$, $\delta=s_{23}s_{25}$.
By \eqref{eq:d4-positive-centre} and $a>d$,
\[
 u=a^3r+\frac{m_5a^3}{d^2}>m_5d=v.
\]
The real and imaginary parts of the first equation give
\[
 \nu^3s_{13}(r+\varepsilon v)=s_{24}u,
 \qquad \nu^3s_{14}u=-s_{23}(r+\delta v).
\]
Thus $|(r+\varepsilon v)(r+\delta v)|=u^2$.
Both factors $r+\varepsilon v$ and $r+\delta v$ must be positive.
Indeed, if a factor is nonpositive then $v\ge r$ and at least one sign
is negative; the absolute product is then either $(v-r)^2$ or
$v^2-r^2$, in both cases strictly less than $v^2<u^2$.
It follows that $s_{24}=s_{13}$ and $s_{14}=-s_{23}$, and hence
\begin{equation}\label{eq:d4-first-real}
 \nu^3(r+\varepsilon v)=u,\qquad
 \nu^3u=r+\delta v.
\end{equation}
Since $g_4^3=-ia^3$, the real and imaginary parts of the second equation
now give
\begin{equation}\label{eq:d4-third-real}
 \nu(r+\varepsilon vd^2)=ua^2,\qquad
 \nu ua^2=r+\delta vd^2.
\end{equation}

If $\varepsilon=\delta=1$, then \eqref{eq:d4-first-real} gives
$\nu=1$ and $u=r+v$.  Put
$w_0=r/(r+m_5/d^2)$ and $w_1=(m_5/d^2)/(r+m_5/d^2)$.
These weights are positive and sum to $1$.
Equations \eqref{eq:d4-positive-centre}, \eqref{eq:d4-first-real} and
\eqref{eq:d4-third-real} give, respectively,
\[
 a^3=w_0+w_1d^3,\qquad a^5=w_0+w_1d^5.
\]
But strict convexity of $t\mapsto t^{5/3}$ on $(0,\infty)$, applied
to $1$ and $d^3\ne1$, gives
$w_0+w_1d^5>(w_0+w_1d^3)^{5/3}=a^5$, a contradiction.

In the other cases \eqref{eq:d4-first-real} gives either
$u=r-v<r$ or $u^2=r^2-v^2<r^2$, so $u<r$.
This forces $a<1$: otherwise \eqref{eq:d4-positive-centre} would give
$r<m_4\le m_4a=u$.  Thus $d<a<1$.
If $\varepsilon=\delta=-1$, then $\nu=1$ and the two balance equations
give $r=u+v=ua^2+vd^2<u+v$, again a contradiction.
Finally, if $\varepsilon=-\delta$, adding the two relations in each
of \eqref{eq:d4-first-real} and \eqref{eq:d4-third-real} yields
\[
 2r=u(\nu^3+\nu^{-3})=ua^2(\nu+\nu^{-1}).
\]
Consequently
$a^2=(\nu^3+\nu^{-3})/(\nu+\nu^{-1})
=\nu^2+\nu^{-2}-1\ge1$, contradicting $a<1$.
\end{proof}

}\subsubsection{Proof of Theorem~\ref{thm:d4}}
{Suppose that $\gamma$ is a Diagram~4 germ at $m\in\Rp^5$.
Proposition~\ref{prop:d4-cover} gives relation (23) and a holomorphic chart representation
$\tau\mapsto(t(\tau),x(\tau))$, with $t=\tau^q$, in $F^{-1}(0)$.
It satisfies $t\not\equiv0$ and $x(0)=y\in\mathcal B_4(m)$.
Lemma~\ref{lem:d4-fibre}(c) permits the application of Lemma~\ref{lem:d4-ls} at $y$,
with $Z=\mathcal B_4(m)$. Taking $N=1$ gives $\Omega_1(y)=0$,
or $f_1=f_2$ by Lemma~\ref{lem:d4-obstr}(a).
Lemma~\ref{lem:d4-obstr}(b) then permits the same argument for $N=3$,
giving $\Omega_3(y)=0$. Since $(A_1-A_2)B_1/A_1\ne0$, it follows that $\Theta=0$ at $y$.}

{Both conditions are symmetric in the bodies of $T$ and homogeneous under
the boundary scalings. We may therefore renumber $T$ as in \S\ref{d4-census} with $b>c$.
Indeed, $b=c$ would contradict $\eta_k(y)\ne0$ in Proposition~\ref{prop:d4-cover},
because $o_3=0$ when $b=c$.
The boundary scalings allow us to normalize $A_1=1$ and $\eta_3=-1$.
Then $\eta_k=-o_k/o_3$ and $D_{1k}^2=-1/\eta_k=g_k^2$.
With $\nu=\sqrt{m_1/m_2}>0$ we have $A_2=-\nu^2$ and
$D_{2k}^2=-D_{1k}^2/\nu^2$, so
$D_{1k}=s_{1k}g_k$ and $D_{2k}=s_{2k}ig_k/\nu$ for a sign vector $s$.
With the original triangle masses in the definitions of $X_j,Y_j$, we obtain
\[
 S_1=X_1,\qquad S_2=iY_1/\nu,\qquad
 \Sigma_1=X_3,\qquad \Sigma_2=-iY_3/\nu^3.
\]
Thus $S_1/A_1=S_2/A_2$ and $A_2\Sigma_2-A_1\Sigma_1=0$ become
$\nu^3X_1+iY_1=0$ and $\nu X_3=iY_3$.
Lemma~\ref{lem:d4-census} gives the contradiction. \qed}

\begin{remark}\label{rem:d4-witnesses}
(a) Positivity of the masses enters through (23) and the parametrization of the triangle,
through $m_1^3+m_2^3\ne0$, and through $\nu>0$; no reality of the germ is used, so
Theorem~\ref{thm:d4} excludes all complex germs.
(b) The two obstructions involve the outer masses only through $\rho=-m_1/m_2$, and not through
the ratio $M_P/M_T$.
\end{remark}


{
\subsection{Diagram 5}\label{sec:d5}

Diagram~5 consists of the isolated $zw$-triangle $345$, with no circled vertex and no
other stroke \cite[\S7.4 and Fig.~11]{AK}.  We write $T=\{3,4,5\}$ and call $1,2$
the outer bodies.  The triangle collapses along the germ, but some of its
distances to the outer bodies may grow.  To control these competing effects,
we estimate the product of all ten squared distances.  Its vanishing is the
local conclusion needed in the global argument.

\begin{theorem}\label{thm:d5}
Let $m\in\Rp^5$ and let $\gamma$ be a Diagram~5 germ on an irreducible curve
$\Gamma\subset A_m$.  Then
\[
 \Delta:=\prod_{1\le i<j\le5}r_{ij}^{2}\longrightarrow0
 \quad\hbox{along }\gamma.
\]
No assumption that $\Gamma$ contains a real point is needed.
\end{theorem}

\subsubsection{Coordinates relative to the triangle}\label{d5-frame}
We normalize as in Lemma~\ref{lem:germ} and measure orders in a Puiseux scale $t$ for
which $\epsilon\approx|t|$ and the least orders of both $\mathcal Z$ and
$\mathcal W$ are $-2$.  If initially $\epsilon\approx|s|^\nu$, we take
$t=cs^\nu$, $c\ne0$, after a finite ramification as necessary.  This choice only
multiplies all orders by a positive factor.  As usual, $\ord0=+\infty$.

Let $M_T=\sum_{k\in T}m_k$, $M=\sum_{k=1}^5m_k$, and let $z_c,w_c$ be the
mass centres of the triangle.  We distinguish the internal coordinates of
the triangle from the positions of the outer bodies relative to its centre:
\[
 \begin{aligned}
 \xi_k&=z_k-z_c,\quad \eta_k=w_k-w_c\quad(k\in T),\\
 X_j&=z_j-z_c,\quad Y_j=w_j-w_c\quad(j=1,2).
 \end{aligned}
\]
Proposition~\ref{prop:d45-first} and Lemma~\ref{lem:d1-shape} give relation~(23)
and show that the limiting triangle has zero multiplier.  After a constant
change of $t$ and a constant rotation $(z,w)\mapsto(az,a^{-1}w)$, its internal
coordinates have the form
\begin{equation}\label{d5-eq-offsets}
 \xi_k=t(s_k+o(1)),\qquad \eta_k=t(s_k+o(1)),\qquad
 s_k\in\R,\quad \sum_{k\in T}m_ks_k=0,
\end{equation}
where the three $s_k$ are pairwise distinct.  Every squared distance, as well
as $U$ and $I$, is preserved by these changes.

The rules and estimates in Fact~\ref{fact:rules} now give the orders that we
shall use.  Each coordinate difference within the triangle has order $1$ in
both colours, so each of its three squared distances has order $2$.
Since there is no circled vertex, all body coordinates, and therefore all
differences from the triangle centre, have order strictly greater than $-2$.
Every crossed
force term $\delta_{jk}^3(z_j-z_k)$ and $\delta_{jk}^3(w_j-w_k)$, $j\in\{1,2\}$,
$k\in T$, has order strictly greater than $-2$, since no stroke joins an outer
body to the triangle.  Finally $\ord r_{ij}\le1$, with equality exactly on
$zw$-edges.

The scale of the triangle is thus fixed.  Relative to this scale, each outer
body belongs to exactly one of the following three types:
\begin{itemize}
 \item \emph{far}: $\ord X_j<1$ and $\ord Y_j<1$;
 \item \emph{near in $z$}: $\ord X_j\ge1$ and $\ord Y_j=\beta_j<1$;
 \item \emph{near in $w$}: the preceding type with the colours exchanged.
\end{itemize}
If both centre orders were at least $1$, every crossed distance would have
order at least $1$, producing additional $zw$-edges.  This is excluded by the
diagram.  Infinite centre orders are nevertheless allowed in a near type.

For a body near in $z$, the limit $\theta_j=\lim t^{-1}X_j$ exists.
Equation~\eqref{d5-eq-offsets} shows that at least two of the differences
$z_j-z_k$ have order $1$.  The third can have a higher order $h_j>1$ only
when its leading coefficient vanishes at one triangle vertex.  We call the
body \emph{ordinary} if all three orders are $1$, and \emph{special}
otherwise; the vertex of higher-order contact is its special vertex.
We set $e_j=0$ for an ordinary body and $e_j=h_j-1$ for a special body.
The distance ties ensure that each edge difference is nonzero.  Thus $h_j$
is finite even if $X_j$ or a centre offset vanishes identically.  Summing
the three squared-distance orders gives
\begin{equation}\label{d5-eq-near-sum}
 \sum_{k\in T}\ord r_{jk}^2=3(1+\beta_j)+e_j
 \qquad\hbox{for a body near in }z.
\end{equation}

We now separate the forces exerted by the triangle on each outer body from
the mutual force of the two outer bodies.  These are, respectively,
\[
 \begin{gathered}
 F_j=\sum_{k\in T}m_k\delta_{jk}^3(z_j-z_k),\qquad
 G_j=\sum_{k\in T}m_k\delta_{jk}^3(w_j-w_k),\\
 P=\delta_{12}^3(X_1-X_2),\qquad H=\delta_{12}^3(Y_1-Y_2).
 \end{gathered}
\]
In these relative coordinates, the original equations and the total
mass-centre identity become
\begin{equation}\label{d5-eq-outer}
 \begin{aligned}
 X_1-X_2&=(m_1+m_2)P+F_1-F_2,\\
 X_1&=m_2P+(1+m_1/M_T)F_1+(m_2/M_T)F_2,\\
 X_2&=-m_1P+(m_1/M_T)F_1+(1+m_2/M_T)F_2.
 \end{aligned}
\end{equation}
The same identities hold with $(X,P,F)$ replaced by $(Y,H,G)$.  Their weighted
sums are
\begin{equation}\label{d5-eq-centre}
 m_1X_1+m_2X_2=\frac M{M_T}(m_1F_1+m_2F_2),\qquad
 m_1Y_1+m_2Y_2=\frac M{M_T}(m_1G_1+m_2G_2).
\end{equation}
For example, multiplying the two original outer $z$-equations by $m_1,m_2$
and adding cancels their mutual force.  The mass-centre identity gives
$m_1z_1+m_2z_2=(M_T/M)(m_1X_1+m_2X_2)$, which proves the first identity.
The coefficients of $(F_1,F_2)$ in the last two equations of
\eqref{d5-eq-outer} form the matrix
\begin{equation}\label{d5-eq-force-matrix}
 C=\begin{pmatrix}1+m_1/M_T&m_2/M_T\\m_1/M_T&1+m_2/M_T\end{pmatrix},
 \qquad
 C^{-1}=\frac1M\begin{pmatrix}M_T+m_2&-m_2\\-m_1&M_T+m_1\end{pmatrix}.
\end{equation}
Its determinant is $M/M_T>0$.  In particular, leading terms of the two
triangle forces cannot cancel in both outer-body equations simultaneously.

Throughout this subsection put
\[
 a=\ord(X_1-X_2),\qquad b=\ord(Y_1-Y_2),\qquad S=a+b,
 \qquad d=\ord\Delta.
\]
Because $12$ is an actual edge, $a,b$ are finite.  The distance ties give
\begin{equation}\label{d5-eq-pair-orders}
 \ord P=a-\frac32S=-\frac{a+3b}{2},\qquad
 \ord H=b-\frac32S=-\frac{3a+b}{2}.
\end{equation}
If the coordinate differences on a crossed edge have orders $h,v$, its
contributions to $F_j,G_j$ have orders $-(h+3v)/2$ and $-(3h+v)/2$,
respectively.  Cancellation in the sum can only increase the order of a
force.  For a special body near in $z$, however, the special edge gives
the unique term of lowest order in each force.  Hence
\begin{equation}\label{d5-eq-special-force}
 \ord F_j=-\frac{h_j+3\beta_j}{2},\qquad
 \ord G_j=-\frac{3h_j+\beta_j}{2}.
\end{equation}

When both bodies are far, we shall also need the inertia.
Lemma~\ref{lem:U-constant} shows that $I$ is constant on every irreducible
complex solution curve.  Its decomposition relative to the triangle centre is
\begin{equation}\label{d5-eq-inertia}
 \begin{gathered}
 I=I_T+J,\qquad I_T=\sum_{k\in T}m_k\xi_k\eta_k=O(t^2),\\
 J=\frac{m_1m_2(X_1-X_2)(Y_1-Y_2)+m_1M_TX_1Y_1+m_2M_TX_2Y_2}{M}.
 \end{gathered}
\end{equation}
Since the internal inertia tends to zero, $J$ is bounded.  This remains true
without assuming a real point on the curve or a nonzero value of $I$.

\subsubsection{The leading external force on the triangle}\label{d5-mb}
The outer-body equations alone do not control possible cancellation when
both outer bodies are near in the same colour.  We shall use the triangle
equations to show that, at the lowest order, the external force is the same
at all three vertices.  Such a force moves the centre but disappears from
the centred equations.

To study the internal forces, put $\hat x=\xi/t$, $\hat y=\eta/t$, and let
\[
 \widehat U(\hat x,\hat y)=\sum_{k<l,\ k,l\in T}m_km_l
 \bigl((\hat x_k-\hat x_l)(\hat y_k-\hat y_l)\bigr)^{-1/2},
\]
where the inverse-distance branches are those selected by the germ.
Let $\phi,\psi$ denote the reduced triangle forces in the $z,w$ directions.
Differentiation and homogeneity give
\begin{equation}\label{d5-eq-potential-identities}
 d\widehat U=-\frac12\sum_km_k(\psi_k\,d\hat x_k+\phi_k\,d\hat y_k),
 \qquad
 \widehat U=\sum_km_k\hat x_k\psi_k=\sum_km_k\hat y_k\phi_k.
\end{equation}
In particular, both forces tend to zero at $(\hat x,\hat y)=(s,s)$.

\begin{lemma}\label{d5-lem-mb}
For a germ of reduced triangles tending to $(s,s)$, put
$\kappa_\psi=\ord\psi$ and $\kappa_\phi=\ord\phi$, where the order of a vector is
the least order of its components.  If $\kappa_\psi<\kappa_\phi$, then
$\kappa_\phi\le2\kappa_\psi$.  The assertion also holds with the two forces
exchanged.
\end{lemma}
\begin{proof}
Translations and independent scalings in the two colours preserve the
zero-multiplier equations.  The resulting orbit is
\[
 O=\{(As+u\mathbf1,Bs+v\mathbf1):A,B\ne0\}.
\]
On this orbit $d\widehat U=0$ and $\widehat U=0$ by
\eqref{d5-eq-potential-identities}.  Let $d_{kl}$ be the inverse distances at
$(s,s)$ and define the symmetric matrix
\[
 L_{kl}=-m_km_ld_{kl}^3\quad(k\ne l),\qquad
 L_{kk}=\sum_{l\ne k}m_km_ld_{kl}^3.
\]
The relations $L\mathbf1=Ls=0$ express the translation and scaling
directions.  Since its off-diagonal entries are nonzero, the matrix has
rank one and kernel $\operatorname{span}\{\mathbf1,s\}$.  Explicitly,
$L=c\,nn^{\mathsf T}$, $c\ne0$, for
$n=(s_4-s_5,s_5-s_3,s_3-s_4)^{\mathsf T}$, so $\ell:=L_{33}\ne0$.
At $(As,Bs)$, choose $\varrho^2=AB$ consistently with the branches.  Differentiating
each potential term twice gives
\[
 \nabla^2\widehat U=\frac1{4\varrho}
 \begin{pmatrix}3A^{-2}L&(AB)^{-1}L\\(AB)^{-1}L&3B^{-2}L\end{pmatrix}.
\]
Its kernel is precisely the tangent space of $O$.  The remaining two
directions describe changes of shape.  We may choose the $\hat x_3,\hat y_3$
directions as a complement because $s_4\ne s_5$.  On this complement, the
Hessian $K$ satisfies
\[
 \det K=\frac{\ell^2}{2\varrho^6}\ne0,\qquad
 (K^{-1})_{11}=\frac{3A^2\varrho}{2\ell}\ne0,
 \qquad (K^{-1})_{22}=\frac{3B^2\varrho}{2\ell}\ne0.
\]
A nearby triangle can therefore be written as
$o(u)+v_1e_3^x+v_2e_3^y$, with $o(u)\in O$.
Expanding in the two transverse directions gives
$\widehat U=\frac12v^{\mathsf T}K(u)v+O(|v|^3)$ and
$\partial_v\widehat U=K(u)v+O(|v|^2)=-\frac{m_3}{2}(\psi_3,\phi_3)$.
Here $v=(v_1,v_2)$ measures the departure from the zero-multiplier orbit;
$u$ may vary along the germ.  Since $K(u(t))$ tends to an invertible matrix,
$v(t)$ and $\partial_v\widehat U$ have the same order.  The Taylor estimates
are uniform for nearby orbit coordinates.
Consequently
\[
 \ord\widehat U\ge2\ord v
 =2\ord\partial_v\widehat U\ge2\min(\kappa_\psi,\kappa_\phi).
\]
Suppose $\kappa_\psi<\kappa_\phi$.  Then $0<\kappa_\psi<\infty$.  The leading
coefficient of $\psi$ satisfies both $\sum_km_k\psi_{k,0}=0$ and
$\sum_km_ks_k\psi_{k,0}=0$, the latter by
\eqref{d5-eq-potential-identities} and $\ord\widehat U>\kappa_\psi$.
These two relations imply that its third component is nonzero: otherwise
$s_4\ne s_5$ would force the other two components to vanish as well.
Thus the transverse gradient has leading term
$t^{\kappa_\psi}(\chi,0)$, $\chi\ne0$.  Substituting it into
\[
 \widehat U=\frac12(\partial_v\widehat U)^{\mathsf T}K(u)^{-1}
 (\partial_v\widehat U)+O(|v|^3)
\]
and using $(K^{-1})_{11}\ne0$ gives $\ord\widehat U=2\kappa_\psi$.
The last identity in \eqref{d5-eq-potential-identities} also gives
$\ord\widehat U\ge\kappa_\phi$, which proves the assertion.
The same argument applies if one force vanishes identically.
\end{proof}

\begin{lemma}[Equality of the leading external forces]\label{d5-lem-uniform}
Suppose both outer bodies are near in $z$.  Let $\kappa_*$ be the least order of
the six individual terms $t^2m_j\delta_{jk}^3(w_k-w_j)$, $j\in\{1,2\}$, $k\in T$.
After summing over $j$, the coefficient of $t^{\kappa_*}$ is the same at all
three triangle vertices.  The same assertion holds with the colours exchanged.
\end{lemma}
\begin{proof}
Write $\ord(z_j-z_k)=1+e_{jk}$, $e_{jk}\ge0$, and
$\ord(w_j-w_k)=\beta_j\in(-2,1)$.  The orders of the reduced forcing terms in
the two colours are
\[
 \kappa_{jk}=\frac{1-\beta_j-3e_{jk}}2,
 \qquad
 \frac{3-3\beta_j-e_{jk}}2
 =2\kappa_{jk}+\frac{1-\beta_j+5e_{jk}}2>2\kappa_{jk}.
\]
The absence of crossed strokes gives $\kappa_*>0$.  There are ordinary edges
with $e_{jk}=0$, and $\beta_j>-2$, so also $\kappa_*<3/2$.
Let $g_k=t^2\sum_jm_j\delta_{jk}^3(w_k-w_j)$ and let $g^z$ be the analogous
$z$-forcing.  A hat on either forcing denotes subtraction of its mass-weighted
mean over $T$.  The exact centred triangle equations are
\[
 \psi=t^3\hat y-\widehat g,\qquad
 \phi=t^3\hat x-\widehat {g^z}.
\]
If the coefficient of $t^{\kappa_*}$ in $g$ were not constant across the
vertices, it would survive subtraction of the mean.  Then
$\kappa_\psi=\kappa_*$, while every term contributing to $\phi$ has order
strictly greater than $2\kappa_*$, including $t^3\hat x$.
This contradicts Lemma~\ref{d5-lem-mb}.
\end{proof}

\begin{lemma}[Support of the leading external force]\label{d5-lem-support}
For a single ordinary body near in $z$, its three lowest $w$-forcing coefficients
are not all equal.  For a special body the lowest forcing is supported only at
its special vertex.  If both outer bodies are special and near in $z$, their
special vertices coincide, their lowest forcings cancel there, and
\begin{equation}\label{d5-eq-special-balance}
 3h_1+\beta_1=3h_2+\beta_2.
\end{equation}
\end{lemma}
\begin{proof}
For an ordinary body,
$z_j-z_k=t(\theta_j-s_k+o(1))$ with all $\theta_j-s_k\ne0$.
The squares of the three forcing coefficients equal a common nonzero constant times
$(\theta_j-s_k)^{-3}$.  If the coefficients were equal, the three distinct
numbers $\theta_j-s_k$ would have equal nonzero cubes and hence would be all
three cube roots of one number.  Their sum would be zero, making
$\theta_j=(s_3+s_4+s_5)/3$ real.  Injectivity of the real cube would then force
the three numbers to coincide, a contradiction.

For a special body the special edge gives the unique lowest-order force
term, so the leading force acts only at its special vertex.
For two special bodies, unequal lowest orders would leave the lowest coefficient
supported at one vertex.  Equal orders with different special vertices would
leave it supported at two vertices.  Both are forbidden by
Lemma~\ref{d5-lem-uniform}: a coefficient supported at one or two vertices
cannot be constant at all three.  Therefore their vertices and orders agree,
and their leading coefficients cancel at the common vertex.  The unscaled order is
$-(3h_j+\beta_j)/2$, which proves \eqref{d5-eq-special-balance}.
\end{proof}

\subsubsection{Two far bodies}\label{d5-far-far}
We first consider the case in which both outer bodies are far.
Let $u_j=\ord X_j<1$, $v_j=\ord Y_j<1$, and
$A=u_1+v_1$, $B=u_2+v_2$.  All crossed squared edges from body $1$ have order
$A$, and those from body $2$ have order $B$, so
\begin{equation}\label{d5-eq-far-product}
 d=6+3A+3B+S.
\end{equation}
Individual force terms give
\[
 \ord F_1\ge u_1-\frac32A,\quad \ord G_1\ge v_1-\frac32A,
 \qquad
 \ord F_2\ge u_2-\frac32B,\quad \ord G_2\ge v_2-\frac32B.
\]
Thus $A<0$ makes the first two force orders strictly greater than $u_1,v_1$,
and $B<0$ gives the analogous conclusion for body $2$.

To show that the product has positive order, we exclude negative orders
among these three groups of distances.  The three terms in $J$ in
\eqref{d5-eq-inertia} have orders $S,A,B$ and nonzero mass coefficients.
Since $J$ is bounded, a negative minimum must occur at least twice.
Thus, if any order is negative, either
$A,B<0$, or, after exchanging the outer bodies, $A\ge0>B$ and $S=B<0$.
Indeed $S<0$ alone would be a unique negative minimum.

Suppose first that $A,B<0$.  In the first equation of \eqref{d5-eq-centre} the
right side has order greater than $\min(u_1,u_2)$, so
\[
 u_1=u_2=u,\qquad m_1\widehat X_1+m_2\widehat X_2=0.
\]
Here the hats denote nonzero leading coefficients.  Since $m_1+m_2>0$, they
are different and $a=u$.  The $w$ equation similarly gives $v_1=v_2=v$ and
$b=v$.  Hence $A=B=S=u+v<0$.  In the first equation of
\eqref{d5-eq-outer}, the left side has order $u$, whereas $P,F_1,F_2$ all have
orders greater than $u$, a contradiction.

Suppose instead that $A\ge0>B$ and $S<0$.  The original equation
$z_2=-m_1P+F_2$ and the total mass centre give
\[
 (m_1+M_T)X_2-m_1X_1=M(-m_1P+F_2).
\]
Its right side has order greater than
$\min(a,u_2)=\min(u_1,u_2)$.  The equality follows directly from
$X_1-X_2$: if $u_1\ne u_2$, then $a=\min(u_1,u_2)$; otherwise $a\ge u_2$.
Consequently $u_1=u_2=u$ and
$(m_1+M_T)\widehat X_2=m_1\widehat X_1$.  Since $M_T>0$, the leading
coefficients differ and $a=u$.  The same argument in $w$ gives
$v_1=v_2=v$ and $b=v$.  This forces $A=B=S$, contradicting $A\ge0>B$.
Exchanging $A,B$ treats the remaining possibility.
Thus $A,B,S\ge0$, and \eqref{d5-eq-far-product} gives $d\ge6>0$.

\subsubsection{One far body and one near body}\label{d5-far-near}
By exchanging the outer bodies and the colours if necessary, we may suppose
that body $1$ is far and body $2$ is near in
$z$.  Write
\[
 u=\ord X_1<1,\quad v=\ord Y_1<1,\quad
 \alpha=\ord X_2\ge1,\quad \beta=\ord Y_2<1.
\]
The orders $u,v,\beta$ are finite and greater than $-2$; $\alpha$ may be infinite.
Here $a=u$.  Put $c=\ord(Y_1-Y_2)\ge\min(v,\beta)>-2$,
$x=u+v$ and $y=1+\beta$, and let $e\ge0$ be the special-edge excess of body $2$.
Then $S=u+c$ and
\begin{equation}\label{d5-eq-mixed-product}
 d=6+S+3x+3y+e.
\end{equation}
If $x\ge0$, then $u>-1$, $S>-3$ and $y>-1$, so $d>0$.

It remains to consider $x<0$.  In this case the triangle forces on body $1$
have higher order than its coordinates:
\[
 \ord F_1\ge u-\frac32x>u,\qquad
 \ord G_1\ge v-\frac32x>v.
\]
The first equation of \eqref{d5-eq-centre} has left side of order $u$, hence
$\ord F_2=u$.  The equation for $X_2$ in \eqref{d5-eq-outer} then forces
$\ord P=u$, since its left side and its $F_1$ term have higher orders.
But $\ord P=-(u+3c)/2$, so $c=-u$ and $S=0$.
As $v<c$, cancellation in $Y_1-Y_2$ requires $\beta=v$ and equal nonzero
leading coefficients in $Y_1,Y_2$.  Positive masses therefore make
$m_1Y_1+m_2Y_2$ have order $v$.  The second equation of
\eqref{d5-eq-centre} now implies $\ord G_2=v$.

If body $2$ is special, \eqref{d5-eq-special-force} gives
\[
 u=-\frac{h+3v}{2},\qquad v=-\frac{3h+v}{2}.
\]
These identities give $v=-h$ and $u=h>1$, a contradiction.
Thus body $2$ is ordinary, so $e=0$ and
$u=\ord F_2\ge-(1+3v)/2$.  Since $u<1$, this also gives $v>-1$.
Using $S=0$, $\beta=v$ in \eqref{d5-eq-mixed-product}, we obtain
\[
 d=9+3u+6v\ge\frac{15+3v}{2}>6>0.
\]

\subsubsection{Two near bodies of opposite colours}\label{d5-opposite-near}
Let body $1$ be near in $z$ and body $2$ near in $w$, and write
\[
 \ord X_1\ge1,\quad v=\ord Y_1<1,\qquad
 u=\ord X_2<1,\quad \ord Y_2\ge1.
\]
Here $u,v>-2$, $a=u$, $b=v$ and $S=u+v$.  If $e_1,e_2\ge0$ are the two
special-edge excesses, summing the ten squared-edge orders gives
\begin{equation}\label{d5-eq-opposite-product}
 d=12+4S+e_1+e_2.
\end{equation}
Suppose $d\le0$.  Then $S\le-3$, and $u,v>-2$ imply $u,v<-1$.
Equation~\eqref{d5-eq-pair-orders} gives $\ord P>u$ and $\ord H>v$.
The mutual forces therefore do not contribute at the leading coordinate
orders.  In matrix form, the outer-body equations read
\[
 C\binom{F_1}{F_2}=\binom{X_1-m_2P}{X_2+m_1P},\qquad
 C\binom{G_1}{G_2}=\binom{Y_1-m_2H}{Y_2+m_1H}.
\]
The right sides have leading terms $(0,X_2)^{\mathsf T}$ at order $u$ and
$(Y_1,0)^{\mathsf T}$ at order $v$, respectively.  Using
\eqref{d5-eq-force-matrix} yields
\[
 F_1\sim-\frac{m_2}{M}X_2,\quad F_2\sim\frac{M_T+m_1}{M}X_2,
 \qquad
 G_1\sim\frac{M_T+m_2}{M}Y_1,\quad G_2\sim-\frac{m_1}{M}Y_1.
\]
All four coefficients are nonzero, so $\ord F_1=\ord F_2=u$ and
$\ord G_1=\ord G_2=v$.
If body $1$ is ordinary, then $\ord G_1\ge-(3+v)/2>v$, a contradiction.
If it is special, then
$v=-(3h_1+v)/2$ gives $v=-h_1$, whence
$\ord F_1=-(h_1+3v)/2=h_1>1$, also a contradiction.
Thus $d>0$.

\subsubsection{Two near bodies of the same colour}\label{d5-same-near}
It remains to consider two bodies near in $z$; the other colour is symmetric.
Put $\beta_j=\ord Y_j\in(-2,1)$ and retain $e_j\ge0$ for their special-edge
excesses.  Now $a\ge1$, $b\ge\min(\beta_1,\beta_2)$, so $S>-1$, and
\begin{equation}\label{d5-eq-same-product}
 d=12+3(\beta_1+\beta_2)+S+e_1+e_2.
\end{equation}
Suppose $d\le0$.  If $S\ge0$, the right side is positive since $\beta_j>-2$.
If, for example, $\beta_1\ge-1$, then $\beta_2>-2$ and $S>-1$ make it greater
than $2$.  Therefore
\begin{equation}\label{d5-eq-negative-orders}
 S<0,\qquad \beta_1<-1,\qquad \beta_2<-1.
\end{equation}

If both bodies are ordinary, then $\ord G_j\ge-(3+\beta_j)/2>\beta_j$.
The $w$ equation in \eqref{d5-eq-centre} rules out unequal $\beta_1,\beta_2$,
since its left side would have a unique lowest term.  Thus
$\beta_1=\beta_2=\beta$; writing $Y_j=v_jt^\beta(1+o(1))$, it gives
$m_1v_1+m_2v_2=0$.  Positivity makes $v_1\ne v_2$, hence $b=\beta$.
By $S<0$, $\ord H=b-3S/2>\beta$.  Both outer $w$-equations in
\eqref{d5-eq-outer} now have right sides of order greater than $\beta$, whereas
their left sides have order $\beta$, a contradiction.

Suppose next that exactly one body is special, say body $1$, with special
order $h>1$.  If the lowest individual $w$-forcing orders differed, only
one body would contribute at the lowest order.  By
Lemma~\ref{d5-lem-support}, its coefficient vector could not be constant
across the triangle vertices, whether that body were ordinary or special.
Lemma~\ref{d5-lem-uniform} therefore requires the two lowest orders to agree:
\[
 3h+\beta_1=3+\beta_2,\qquad \beta_1<\beta_2.
\]
Their common unscaled order is $q=-(3h+\beta_1)/2=-(3+\beta_2)/2$.
Equation~\eqref{d5-eq-negative-orders} implies
\[
 h+\beta_1=\frac{3+\beta_2+2\beta_1}{3}<0,\qquad
 q-\beta_1=-\frac32(h+\beta_1)>0.
\]
Both $G_j$ have order at least $q$.  But the $w$ equation of
\eqref{d5-eq-centre} has a unique lowest term of order $\beta_1$ on the left,
and right side of order at least $q>\beta_1$, a contradiction.

Finally suppose both bodies are special.  Lemma~\ref{d5-lem-support} gives a
common special vertex and $3h_1+\beta_1=3h_2+\beta_2$.
If $h_1=h_2=h$, then $\beta_1=\beta_2=\beta$, $a\ge h$ and $b\ge\beta$.
Thus $h+\beta\le S<0$ and
\[
 \ord G_j=-\frac{3h+\beta}{2}
 =\beta-\frac32(h+\beta)>\beta,\qquad
 \ord H=b-\frac32S>b\ge\beta.
\]
The outer $w$-equations again give a contradiction.
If $h_1\ne h_2$, interchange the bodies to arrange $h_1<h_2$, and put
$\eta=h_2-h_1>0$.  Then
\[
 \beta_1=\beta_2+3\eta,\qquad
 a=h_1=h_2-\eta,\qquad b=\beta_2.
\]
Here $a=h_1$ follows by subtracting the two actual differences to their common
special vertex.  Both $G_j$ have the same exact order
$q=-(3h_2+\beta_2)/2$, while $\ord H=\beta_2-3S/2>\beta_2$.
If $q>\beta_2$, the second outer $w$-equation has no term on the right of order
$\beta_2$.  If $q<\beta_2$, both left sides and $H$ have orders greater than
$q$, so the two nonzero leading coefficients of $G_1,G_2$ would be annihilated
by the invertible matrix $C$, also impossible.  Thus $q=\beta_2$, giving
$\beta_2=-h_2$ and $S=-\eta$.
The exact special-force orders and \eqref{d5-eq-pair-orders} now give
\[
 \ord F_1=h_2-4\eta,\qquad
 \ord F_2=h_2,\qquad \ord P=h_2+\frac\eta2.
\]
In the first equation of \eqref{d5-eq-outer}, $F_1$ is the unique lowest term, forcing
$a=h_2-4\eta$, contrary to $a=h_2-\eta$.
Hence $d>0$ also when both bodies are near in the same colour.

\subsubsection{Completion of the proof}
Up to exchanging the outer bodies and the colours, the four cases cover
all possible positions of the outer bodies.  They show that every Diagram~5
germ on a complex solution curve has $\ord\Delta>0$, proving
Theorem~\ref{thm:d5}.

The argument also allows infinite orders for differences from the triangle
centre; such a difference occurs only in a near type.  Some internal centre
offsets may vanish identically as well.  This does not affect the proof,
because the internal pairwise differences retain order $1$, and the
ordinary/special distinction is made using crossed edge differences.
Every special-edge order and every order entering a pair force is finite.
Thus the proof applies directly to the Puiseux germ, including these
possible degeneracies.
\qed
}

\section{Diagram 6}\label{sec:d6}

\subsection{The diagram and the butterfly chart}

{
Diagram~6 is the \emph{butterfly} of \cite[\S7.4]{AK}: a $z$-triangle and a
$w$-triangle, called the \emph{wings}, meet at one vertex.  We use the numbering
of \cite[\S8.6]{AK}.  The $z$-wing $\{1,2,5\}$ has strokes $12,15,25$ and
$z$-circles at $1,2$; the $w$-wing $\{3,4,5\}$ has strokes $34,35,45$ and
$w$-circles at $3,4$.  Vertex $5$ is uncircled, and there is no stroke between
$\{1,2\}$ and $\{3,4\}$.

The three possible circlings of the butterfly give Diagrams~6--8;
$5$ is uncircled in each of them \cite[\S7.4]{AK} (Fact~\ref{fact:fig11}).
Diagram~6 is the case in which both wings are circled.  Rules~1a and~1d place
the circles of the $z$-wing exactly at $1,2$, and give the corresponding
statement for $w$.
}

{
Permuting the bodies together with their masses, or exchanging
$z\leftrightarrow w$, preserves \eqref{eq:system}.  Thus the thirty coloured
labels of Diagram~6, specified by a common vertex and an ordered pair of
wings, are obtained from the numbering above.  The exchange
$z\leftrightarrow w$ composed with $(13)(24)$ preserves this numbering.
It suffices to treat this case at arbitrary positive masses.

The two circled wings have opposite coordinate scales.  We use a chart
adapted to these scales, in which a Diagram~6 germ has a finite limit.
The chart coordinates are rotation invariants, so the gauge $z_{12}=w_{12}$
of $A_m$ does not enter the construction.
}

\begin{definition}\label{d6-def-chart}
{
The \emph{butterfly chart} has $21$ coordinates: $t$, $Z_1,\dots,Z_5$,
$W_1,\dots,W_5$, and $D_{kl}$.  They are related to \eqref{eq:system} by
$(z_k,w_k)=(t^{-1}Z_k,tW_k)$ for $k=1,2$,
$(z_k,w_k)=(tZ_k,t^{-1}W_k)$ for $k=3,4$, and
$(z_5,w_5)=(tZ_5,tW_5)$.  On the \emph{crossed pairs} $13,14,23,24$ we set
$\delta_{kl}=tD_{kl}$; on the \emph{wing pairs} $12,15,25,34,35,45$ we set
$\delta_{kl}=D_{kl}$.

The defining equations, also called the \emph{rows}, are obtained as follows.
For each body, we substitute these expressions into the $z$-equation
$z_k-\sum_{l\ne k}m_l\delta_{kl}^3(z_k-z_l)=0$ and the corresponding equation
for $w$.  We multiply the two equations by $(t,t^{-1})$ for $k=1,2$,
by $(t^{-1},t)$ for $k=3,4$, and by $(t,t)$ for $k=5$; the resulting rows are
$R^z_k,R^w_k$.  We also retain the ten distance equations
$\delta_{kl}^2z_{kl}w_{kl}-1=0$ and impose the gauges
$Z_2-Z_1-1=0$, $W_4-W_3-1=0$.

We denote these $22$ rows by $F=0$ and call
$\mathcal B=F^{-1}(0)\cap\{t=0\}$ the \emph{boundary fibre}.
The notation $Z_k,W_k$ for chart coordinates is distinct from the force
terms $Z_{kl}=z_{kl}\delta_{kl}^3$ in \cite[\S3]{AK}.
}
\end{definition}

{\begin{lemma}\label{d6-lem-chart}
Let $m\in\Rp^5$.
\begin{romanlist}
  \item The rows are polynomials in the $21$ coordinates, of degree at most $5$ in $t$.
  \item Let $(z,w,\delta)\in A_m$ (so $z_{12}w_{34}\ne0$), choose $t$ with
  $t^2=(z_{12}w_{34})^{-1}$ and put $\nu=(tz_{12})^{-1}$.  The chart point of the rotated
  solution $(\nu z,\nu^{-1}w,\delta)$ has $Z_k=z_k/z_{12}$, $W_k=w_kz_{12}$ $(k=1,2)$,
  $Z_k=z_kw_{34}$, $W_k=w_k/w_{34}$ $(k=3,4)$, $Z_5=z_5w_{34}$, $W_5=w_5z_{12}$,
  $D_{kl}=\delta_{kl}/t$ on crossed pairs and $D_{kl}=\delta_{kl}$ on wing pairs, and satisfies
  $F=0$: each coordinate row is $t^{\pm1}\nu^{\pm1}$ times an equation of \eqref{eq:system}, the
  distance rows are those of \eqref{eq:system}, and both gauges hold.
  \item At $t=0$ the crossed distance rows read $D_{kl}^2Z_kW_l=-1$, and the wing distance rows
  read $D_{12}^2(Z_2-Z_1)(W_2-W_1)=1$, $-D_{15}^2Z_1(W_5-W_1)=1$, $-D_{25}^2Z_2(W_5-W_2)=1$,
  $D_{34}^2(Z_4-Z_3)(W_4-W_3)=1$, $-D_{35}^2(Z_5-Z_3)W_3=1$, $-D_{45}^2(Z_5-Z_4)W_4=1$.  So
  every $D_{kl}$ is nonzero on $\mathcal B$.
  \item Identically,
  $m_1R^z_1+m_2R^z_2+m_5R^z_5+t^2(m_3R^z_3+m_4R^z_4)=m_1Z_1+m_2Z_2+t^2(m_3Z_3+m_4Z_4+m_5Z_5)$,
  and likewise after $z\leftrightarrow w$, $1\leftrightarrow3$, $2\leftrightarrow4$.  Hence on
  $\mathcal B$: $Z_1=-m_2/(m_1+m_2)$, $Z_2=m_1/(m_1+m_2)$, $W_3=-m_4/(m_3+m_4)$,
  $W_4=m_3/(m_3+m_4)$.
  \item At $t=0$ the rows $R^z_1,R^z_2,R^z_5,R^w_1,R^w_2$, the distance rows of $12,15,25$ and
  the first gauge involve only the \emph{$z$-wing coordinates}
  $Z_1,Z_2,W_1,W_2,W_5,D_{12},D_{15},D_{25}$; the remaining rows except the crossed distance rows
  involve only the \emph{$w$-wing coordinates} $Z_3,Z_4,Z_5,W_3,W_4,D_{34},D_{35},D_{45}$; and
  $D_{13},D_{14},D_{23},D_{24}$ occur only in the crossed distance rows.
\end{romanlist}
\end{lemma}}
{\begin{proof}
Put $I=\{1,2\}$ and $J=\{3,4\}$, and let $i'$ be the other element of $I$.
Substitution in the defining equations gives
\begin{align*}
R_i^z={}&Z_i-m_{i'}D_{ii'}^3(Z_i-Z_{i'})
 -m_5D_{i5}^3(Z_i-t^2Z_5)
 -t^3\sum_{j\in J}m_jD_{ij}^3(Z_i-t^2Z_j),\\
R_i^w={}&W_i-m_{i'}D_{ii'}^3(W_i-W_{i'})
 -m_5D_{i5}^3(W_i-W_5)
 -t\sum_{j\in J}m_jD_{ij}^3(t^2W_i-W_j),\\
R_5^z={}&t^2Z_5-\sum_{i\in I}m_iD_{i5}^3(t^2Z_5-Z_i)
 -t^2\sum_{j\in J}m_jD_{j5}^3(Z_5-Z_j).
\end{align*}
The other coordinate rows follow by exchanging the wings and $z,w$.
The distance rows for $ii'$, $i5$, $j5$, and $ij$ are, respectively,
\begin{gather*}
D_{ii'}^2(Z_i-Z_{i'})(W_i-W_{i'})-1,\qquad
D_{i5}^2(Z_i-t^2Z_5)(W_i-W_5)-1,\\
D_{j5}^2(Z_j-Z_5)(W_j-t^2W_5)-1,\qquad
D_{ij}^2(Z_i-t^2Z_j)(t^2W_i-W_j)-1,
\end{gather*}
and the $34$ row is the wing exchange of the $12$ row.
These formulas prove (i), (iii), and (v). For (ii), the displayed change of
variables is precisely the substitution defining these rows, with
$t^2z_{12}w_{34}=1$ and $\nu tz_{12}=1$; the two gauges therefore hold.
Finally, in the weighted sum of the original $z$-equations every pair force
cancels with the force of the reversed pair. Multiplying this identity by
$t$ gives (iv). At $t=0$, its first identity and the first gauge give
$Z_1=-m_2/(m_1+m_2)$ and $Z_2=m_1/(m_1+m_2)$; the other wing is identical.
\end{proof}

}\subsection{Limits in the butterfly chart}

\begin{proposition}[Limits in the chart]\label{prop:d6-cover}
Let $m\in\Rp^5$ and let $\gamma(s)$, $s\to0$, be a Diagram-6 germ, in the numbering above, on a
curve in $A_m$.  Along $\gamma$ define $t$ and $\nu$ as in Lemma~\ref{d6-lem-chart}(ii), and
let $\tilde\gamma(s)$ be the chart point of the rotated solution.  Then $t\to0$, every chart
coordinate is bounded along $\tilde\gamma$, and $\tilde\gamma(s)$ converges to a point
$p\in\mathcal B$ at which $Z_1,Z_2,W_3,W_4$ and all ten $D_{kl}$ are nonzero.
\end{proposition}

\begin{proof}
{
We normalize as in \cite[p.~547]{AK}, with the largest components of
$\mathcal Z$ and $\mathcal W$ both having modulus $\epsilon^{-2}$.
Lemma~\ref{lem:germ} realizes this normalization along $\gamma$, with
$\epsilon$ comparable to a positive power of $s$.  We write $\ord$ for order
in $\epsilon$; thus $\ord f>\alpha$ means $f\epsilon^{-\alpha}\to0$.
Since the chart coordinates are invariant under $z\mapsto az$, $w\mapsto a^{-1}w$,
we may compute their orders in this frame.
}

{
The circles give $\ord z_1=\ord z_2=\ord w_3=\ord w_4=-2$.
The uncircled coordinates $z_3,z_4,z_5,w_1,w_2,w_5$ have order $>-2$,
so $\ord z_{15}=\ord z_{25}=-2$.  A $z$-stroke $kl$ means
$\ord(z_{kl}\delta_{kl}^3)=\ord(z_{kl}^{-1/2}w_{kl}^{-3/2})=-2$, or equivalently
$\ord z_{kl}+3\ord w_{kl}=4$.  Thus $\ord w_{15}=\ord w_{25}=2$, and hence
$\ord w_{12}\ge2$.  The stroke $12$ then gives $\ord z_{12}\le-2$.
Combining this with $\ord z_{12}\ge-2$ (Estimate~1, Fact~\ref{fact:rules}),
we obtain $\ord z_{12}=-2$, $\ord w_{12}=2$.
Consequently $r_{kl}^2=z_{kl}w_{kl}$ and $\delta_{kl}$ have order $0$ on
all three $z$-wing pairs; these are the maximal edges with $r\approx1$ in
Estimate~2.  The same argument applies to the $w$-wing.

On a crossed pair, both $z_{kl}$ and $w_{kl}$ are differences of a circled
and an uncircled coordinate.  Both therefore have order $-2$, and
$\ord\delta_{kl}=2$.

We also need bounds on the small coordinate of each outer vertex.
In the $w$-equation $w_1=\sum_{l\ne1}m_l\delta_{1l}^3(w_1-w_l)$,
the terms have orders $2$ for $l=2,5$ and $4$ for $l=3,4$.
It follows that $\ord w_1\ge2$.  Similarly,
$\ord w_2,\ord z_3,\ord z_4\ge2$, while
$w_5=w_1+w_{15}$ and $z_5=z_3+z_{35}$ have order $\ge2$.
These are the estimates $A_w,B_z,m_1w_1+m_2w_2,w_5,B_w,A_z\preccurlyeq\epsilon^2$
of \cite[\S8.6]{AK}.

Finally, $\ord t=-\frac12(\ord z_{12}+\ord w_{34})=2$ and
$\ord\nu=-\ord t-\ord z_{12}=0$.
Lemma~\ref{d6-lem-chart}(ii) now shows that $Z_1,Z_2,W_3,W_4$ and all ten
$D_{kl}$ have order exactly $0$; for the crossed pairs the order is
$2+\frac12(-2-2)=0$.  The remaining coordinates
$W_1,W_2,W_5,Z_3,Z_4,Z_5$ have order $\ge0$.
Thus all chart coordinates are bounded.  Their Puiseux series in $s$
converge, giving $\tilde\gamma(s)\to p\in\mathcal B$.
Coordinates of order exactly $0$ have nonzero limits.
}
\end{proof}
In particular $D_{12}=\delta_{12}$ and $D_{34}=\delta_{34}$ have finite nonzero limits along every
Diagram-6 germ, at every mass.

\subsection{A necessary mass condition}

{
The torque identity imposes a condition on the masses before any local
calculation at the boundary points.  Its leading crossed contribution is
$\Sigma=\sum_{k\in\{1,2\},\,l\in\{3,4\}}m_km_lD_{kl}^3Z_kW_l$.
}

{\begin{lemma}\label{d6-lem-torque}
There is a polynomial $Q$ in the chart coordinates, with coefficients polynomial in $m$, such
that identically
\begin{equation}\label{d6-eq-torque}
  \sum_{k=1,2}m_k\bigl(W_kR^z_k-Z_kR^w_k\bigr)+m_5W_5R^z_5=-t\,\Sigma+t^2Q.
\end{equation}
\end{lemma}}
{\begin{proof}
Let $I=\{1,2\}$ and $J=\{3,4\}$. The polynomial in the statement can be chosen as
\begin{align*}
Q={}&t^3\sum_{i\in I,\,j\in J}m_im_jD_{ij}^3W_iZ_j
 +m_5W_5Z_5\\
&+m_5Z_5\sum_{i\in I}m_iD_{i5}^3(W_i-W_5)
 -m_5W_5\sum_{j\in J}m_jD_{j5}^3(Z_5-Z_j).
\end{align*}
Indeed, the terms for the pair $12$ cancel in
$\sum_{i\in I}m_i(W_iR_i^z-Z_iR_i^w)$.
Each crossed pair contributes
$m_im_jD_{ij}^3(-tZ_iW_j+t^5W_iZ_j)$,
and each pair $i5$ contributes
$m_im_5D_{i5}^3(t^2W_iZ_5-Z_iW_5)$.
The terms $-m_im_5D_{i5}^3Z_iW_5$ cancel with the constant part of
$m_5W_5R_5^z$. The remaining terms are exactly $-t\Sigma+t^2Q$.
\end{proof}

}{
To relate this identity to \cite[\S8.6]{AK}, let $\rho^z_k,\rho^w_k$ be
the left-hand sides of \eqref{eq:system}.  The pair torque
$\sum_{k=1,2}m_k(w_k\rho^z_k-z_k\rho^w_k)$ is
$\sum_{k=1,2}\sum_{l=3,4,5}m_km_l\delta_{kl}^3(w_kz_l-z_kw_l)$,
since the terms with $l\in\{1,2\}$ cancel.  It is $2i$ times (31) of that paper.
In the chart, this torque becomes the first sum in \eqref{d6-eq-torque}.
The terms with $l=3,4$ contribute $-t\Sigma+t^5\sum m_km_lD_{kl}^3W_kZ_l$,
and those with $l=5$ contribute
$\sum_{k=1,2}m_km_5D_{k5}^3(t^2W_kZ_5-Z_kW_5)$.
The additional row in the lemma removes the latter contribution at order
zero, because $R^z_5-(m_1D_{15}^3Z_1+m_2D_{25}^3Z_2)$ is divisible by $t^2$.
}

\begin{proposition}[Necessary mass condition]\label{prop:d6-first}
Let $m\in\Rp^5$, let $\gamma$ be a Diagram-6 germ in the numbering above, and let $p$ be its
limit (Proposition~\ref{prop:d6-cover}).  Then:
\begin{romanlist}
  \item $\Sigma(p)=0$, and $\Sigma(p)=-(m_1m_3D_{13}+m_1m_4D_{14}+m_2m_3D_{23}+m_2m_4D_{24})(p)$;
  \item $m_1=m_2$ and $m_3=m_4$; this is relation (32) of \cite[\S8.6]{AK};
  \item $(D_{13},D_{14},D_{23},D_{24})(p)=(2i\eta_{13},\,2\eta_{14},\,2\eta_{23},\,2i\eta_{24})$
  with signs $\eta_{kl}\in\{\pm1\}$ satisfying $\eta_{14}=-\eta_{23}$ and $\eta_{13}=-\eta_{24}$.
\end{romanlist}
In another numbering, wings $\{a,b,c\}$ and $\{d,e,c\}$ force $m_a=m_b$ and $m_d=m_e$.
\end{proposition}

\begin{proof}
All rows vanish along $\tilde\gamma$, so \eqref{d6-eq-torque} gives $t\Sigma=t^2Q$, and $t\ne0$
gives $\Sigma=tQ$.  By Proposition~\ref{prop:d6-cover}, $Q$ is bounded and $t\to0$, so
$\Sigma(p)=0$.  By Lemma~\ref{d6-lem-chart}(iii), $D_{kl}^2Z_kW_l=-1$ at $p$ for the crossed
pairs, so $\Sigma(p)=\sum m_km_lD_{kl}(D_{kl}^2Z_kW_l)=-\sum m_km_lD_{kl}$; this is (i).

By Lemma~\ref{d6-lem-chart}(iv), $Z_1W_3=m_2m_4/s_0^2$, $Z_1W_4=-m_2m_3/s_0^2$,
$Z_2W_3=-m_1m_4/s_0^2$ and $Z_2W_4=m_1m_3/s_0^2$ at $p$, where $s_0>0$,
$s_0^2=(m_1+m_2)(m_3+m_4)$.  Write $m_k=\mu_k^2$ with $\mu_k>0$.  Then $D_{kl}^2=-1/(Z_kW_l)$ gives
\[
  D_{13}=\frac{i\eta_{13}s_0}{\mu_2\mu_4},\quad D_{14}=\frac{\eta_{14}s_0}{\mu_2\mu_3},\quad
  D_{23}=\frac{\eta_{23}s_0}{\mu_1\mu_4},\quad D_{24}=\frac{i\eta_{24}s_0}{\mu_1\mu_3}
\]
with signs $\eta_{kl}\in\{\pm1\}$, and
\[
  \Sigma(p)=-\frac{s_0}{\mu_1\mu_2\mu_3\mu_4}\bigl(P_{\mathrm{re}}+iP_{\mathrm{im}}\bigr),\qquad
  P_{\mathrm{re}}=\eta_{14}\mu_1^3\mu_4^3+\eta_{23}\mu_2^3\mu_3^3,\quad
  P_{\mathrm{im}}=\eta_{13}\mu_1^3\mu_3^3+\eta_{24}\mu_2^3\mu_4^3.
\]
Up to the factor $\pm i$, $P_{\mathrm{re}}+iP_{\mathrm{im}}$ is the left-hand side of (12) of
\cite[\S5.1]{AK} written in the $\mu_k$, and each of its eight sign choices arises from exactly two
of the sixteen sign classes $\eta$.
Since the $\mu_k$ and $s_0$ are real, $\Sigma(p)=0$ forces $P_{\mathrm{re}}=P_{\mathrm{im}}=0$.  If
$\eta_{14}=\eta_{23}$, then $P_{\mathrm{re}}\ne0$, and if $\eta_{13}=\eta_{24}$, then
$P_{\mathrm{im}}\ne0$; this excludes twelve classes.  In the other four,
$P_{\mathrm{re}}=\eta_{14}(x-y)(x^2+xy+y^2)$ with $x=\mu_1\mu_4$, $y=\mu_2\mu_3$ and
$x^2+xy+y^2>0$, so $\mu_1\mu_4=\mu_2\mu_3$; likewise $P_{\mathrm{im}}=0$ gives
$\mu_1\mu_3=\mu_2\mu_4$.  Then
$\mu_3(\mu_1\mu_4-\mu_2\mu_3)-\mu_4(\mu_1\mu_3-\mu_2\mu_4)=\mu_2(\mu_4-\mu_3)(\mu_4+\mu_3)$
vanishes, so $\mu_3=\mu_4$ and then $\mu_1=\mu_2$: this is (ii).  With $m_1=m_2$ and $m_3=m_4$
the displayed values of $D_{kl}$ become those of (iii).
\end{proof}
{
When $m_1=m_2$, $m_3=m_4$, the expression $\Sigma(p)$ vanishes for exactly
the four sign classes in (iii).  This recovers relation~(32) of
Albouy and Kaloshin \cite[\S8.6]{AK}, where it is deduced ``as in 5.1''
from the estimate $S\preccurlyeq\epsilon^4$ in (31).  Here the polynomial
torque identity gives the relation directly at the boundary point.
}

\subsection{The boundary points}

{
We henceforth write $m=(a,a,b,b,c)$ with $a,b,c>0$.
At the boundary, the equations separate into the two wings and the crossed
distance equations.  We first solve the wing equations explicitly.  This
will give a finite list of boundary points at each positive mass.
}

{\begin{lemma}\label{d6-lem-atlas}
The points of $\mathcal B$ are exactly the points with $Z_1=W_3=-\frac12$, $Z_2=W_4=\frac12$,
\begin{gather*}
  (D_{12},D_{15},D_{25})=(d,\lambda d,\omega\lambda d),\\
  W_1=-d(a+2c\lambda),\quad
  W_2=d(a+2c\omega\lambda),\quad W_5=ad\bigl(4\lambda^{-2}-1\bigr),\\
  (D_{34},D_{35},D_{45})=(e,\kappa e,\omega'\kappa e),\\
  Z_3=-e(b+2c\kappa),\quad
  Z_4=e(b+2c\omega'\kappa),\quad Z_5=be\bigl(4\kappa^{-2}-1\bigr),\\
  (D_{13},D_{14},D_{23},D_{24})=(2i\eta_{13},\,2\eta_{14},\,2\eta_{23},\,2i\eta_{24}),\quad
  \eta_{kl}\in\{\pm1\},
\end{gather*}
where
\begin{equation}\label{d6-eq-wing}
  \omega^3=1,\quad\lambda^2=2(1+\omega),\quad(2a+c\lambda^3)d^3=1,\qquad
  \omega'^3=1,\quad\kappa^2=2(1+\omega'),\quad(2b+c\kappa^3)e^3=1.
\end{equation}
For each $\omega$ both roots $\lambda$ are nonzero.  The scale denominator $2a+c\lambda^3$
vanishes only for $(\omega,\lambda)=(1,-2)$ and $a=4c$: for $\omega=1$, $\lambda=\pm2$; for
$\omega\ne1$, $\lambda^6=-8$ and the two denominators have product $4a^2+8c^2$.  Hence a wing
has $18$ classes $(\omega,\lambda,d)$ if $a\ne4c$ and $15$ if $a=4c$, and $\mathcal B$ has
$18\cdot18\cdot16=5184$ points if $a\ne4c\ne b$, $4320$ if exactly one of $a,b$ equals $4c$, and
$3600$ if $a=b=4c$.
\end{lemma}}
{\begin{proof}
By Lemma~\ref{d6-lem-chart}(iv),(v), the boundary equations separate
into the two wing systems and an independent square-root choice on each
crossed pair. On the $z$-wing,
$Z_1=-1/2$, $Z_2=1/2$. Write
$d=D_{12}$, $s=D_{15}$, $q=D_{25}$; all three are nonzero.
The equation $R_5^z=0$ gives $q^3=s^3$, whereas $R_1^z=0$ gives
$2ad^3+cs^3=1$. The distance equations give
\[
W_2-W_1=d^{-2},\qquad W_5-W_1=2s^{-2},\qquad
W_2-W_5=2q^{-2}.
\]
Thus, with $\lambda=s/d$ and $\omega=q/s$,
\[
\omega^3=1,\qquad \lambda^2=2(1+\omega),\qquad
(2a+c\lambda^3)d^3=1.
\]
The equation $R_1^w=0$ now reads $W_1+ad+2cs=0$.
Consequently
\begin{align*}
W_1&=-d(a+2c\lambda),\\
W_2&=W_1+d^{-2}=d(a+2c\omega\lambda),\\
W_5&=W_1+2s^{-2}=ad(4\lambda^{-2}-1),
\end{align*}
where the second equality uses $\lambda^3-2\lambda=2\omega\lambda$.
Conversely, these formulas satisfy the three distance equations,
$R_1^w=0$, and
$R_2^w=W_2-ad-2cq=0$. They also satisfy $R_5^z=0$ and
$R_1^z=R_2^z=0$ by the two cubic equations above, as well as the gauge.
This accounts for all nine wing equations.  Exchanging the wings gives
the second system.  Each crossed equation has the two square roots
displayed in the statement, and the four choices are independent.

For $\omega=1$ the two ratios are $\lambda=\pm2$; only $\lambda=-2$
can have zero scale denominator, precisely at $a=4c$.
For $\omega\ne1$, $\lambda^6=-8$, so the product of the two scale
denominators is $4a^2+8c^2>0$. Every nonzero denominator gives three distinct
cube-root scales.  Counting these choices gives the stated number of
boundary points, including those on the exceptional mass faces.
\end{proof}

}{
The germs at infinity are complex (Remark~\ref{rem:levels}).  The boundary
analysis must therefore include all eighteen wing classes, even though the
positive real class has $(\omega,\lambda)=(1,2)$ and $D_{15}=D_{25}=2D_{12}$.
}

\begin{definition}\label{d6-def-labels}
{
A \emph{label} specifies a point of $\mathcal B$ by the choices
$(\omega,\lambda,d)$, $(\omega',\kappa,e)$ and $\eta$ of
Lemma~\ref{d6-lem-atlas}.

We distinguish two kinds of wing \emph{shape}.
We classify the shapes $(\omega,\lambda)$ and $(\omega',\kappa)$ separately for each wing.
The first is called \emph{Gaussian} when $\omega=1$ (so $\lambda=\pm2$), and \emph{primitive}
otherwise (so $\lambda^4-2\lambda^2+4=0$).  The same definition, with the primed
parameters, applies to the second wing.
Thus each wing has two Gaussian and four primitive shapes.
A label is called \emph{corank-one} when $\eta_{14}=-\eta_{23}$ and
$\eta_{13}=-\eta_{24}$, and \emph{regular} otherwise.
When $a\ne4c\ne b$, there are $12\cdot324=3888$ regular labels and
$4\cdot324=1296$ corank-one labels.
}
\end{definition}

{
The symmetries reduce the number of symbolic calculations required.
At masses $(a,a,b,b,c)$, $F^{-1}(0)$ is preserved by exchanging $1$ and $2$
together with $t\mapsto it$,
$(Z_1,Z_2,W_1,W_2,W_5)\mapsto(-Z_2,-Z_1,-W_2,-W_1,-W_5)$ and
$D_{kl}\mapsto-iD_{kl}$ on crossed pairs.  It is also preserved by
$(t,D_{13},D_{14},D_{23},D_{24})\mapsto-(t,D_{13},D_{14},D_{23},D_{24})$
and by complex conjugation.  The \emph{wing exchange}
($z\leftrightarrow w$, $1\leftrightarrow3$, $2\leftrightarrow4$) maps the chart
at $(a,a,b,b,c)$ to the chart at $(b,b,a,a,c)$.

A \emph{state} records a $z$-wing shape, a $w$-wing shape and $\eta$,
but leaves the two cube-root scales unspecified.  The $576$ states form
$29$ symmetry orbits: $21$ regular and $8$ corank-one
(Appendix~\ref{app:verification}).
In the computations, we express the masses through the scales as
$a=\frac12(d^{-3}-c\lambda^3)$ and $b=\frac12(e^{-3}-c\kappa^3)$.
One symbolic calculation for a state therefore covers its nine choices
of $(d,e)$ at every positive $(a,b,c)$.
}

\begin{remark}[Regular labels]\label{d6-rem-regular}
{
Proposition~\ref{prop:d6-first}(iii) already excludes regular labels as
limits of Diagram~6 germs.  There is also an independent local exclusion:
$DF$ has rank $21$ at every regular label, so each is an isolated point
of $F^{-1}(0)$.

For one representative state of each regular orbit, the minor obtained
by omitting $R^z_1$ is a unit in $\{\pm1,\pm i,\pm1\pm i\}$ times
$-\frac{81}{128}a^2b\,d^2e^4H(\lambda)H(\kappa)(2a+4c\lambda)/(2a+c\lambda^3)$,
where $H(x)=x^3(x^2-2)^2((x^2-2)^2+4)$.  This expression is nonzero
at every finite label.  The minors are computed in exact arithmetic over
$\mathbb Q(i,\sqrt2,\sqrt3)(c,d,e)$, and nonvanishing is checked by
resultants and norms (Appendix~\ref{app:verification}).
}
\end{remark}

{
It remains to exclude the corank-one labels allowed by the torque
condition.  The next lemma does this by solving twenty equations locally
and showing that a remaining equation has a nonzero leading Taylor
coefficient.  The determinant and coefficient calculations are
computer-assisted and use exact arithmetic.
}

\begin{lemma}[Computer-assisted isolation]\label{d6-lem-corank}
Let $p$ be a corank-one label.
\begin{romanlist}
  \item The $20$ rows of $F$ other than $R^w_2$ and $R^w_4$, together with the coordinate $t$,
  have invertible derivative at $p$.  Hence near $p$ the common zeros of these $20$ rows form a
  smooth curve $C_p$, the graph of an analytic map of $t$.
  \item Along $C_p$, $R^w_2=c_kt^k+O(t^{k+1})$ with $c_k\ne0$, where $k$ and $c_k$ are as in
  Table~\ref{d6-tab-corank}.
  \item Consequently $p$ is an isolated point of $F^{-1}(0)$.
\end{romanlist}
\end{lemma}

\begin{table}[ht]
\centering
\small
\begin{tabular}{p{3.3cm}p{0.9cm}p{0.3cm}p{6.6cm}}
\hline
wing shapes ($z$, $w$) & labels & $k$ & $c_k$ is a nonzero multiple of \\
\hline
Gaussian, Gaussian & $144$ & $4$ &
$Q_{2,2}=-\frac{2a+2b+c}{60(a+4c)(b+4c)}$, $Q_{2,-2}=\frac{2a+2b+c}{4(a+4c)(b-4c)}$,
$Q_{-2,-2}=\frac{2a+2b+c}{68(a-4c)(b-4c)}$ for $(\lambda,\kappa)=(2,2),(2,-2),(-2,-2)$;
$(-2,2)$ by wing exchange \\
one Gaussian, one primitive & $576$ & $3$ &
$P=A\kappa^3-4A\kappa+B\kappa^3-B\kappa^2-4B\kappa+B$ for $\lambda=\pm2$, $\kappa$ primitive;
product over the primitive $\kappa$: $9(8A^2+16AB+9B^2)^2$ \\
primitive, primitive: orbit of $(\lambda_0,\lambda_0)$ & $288$ & $2$ &
$1-2\kappa^3$; product over the primitive $\kappa$: $1089$ \\
primitive, primitive: orbits of $(\lambda_0,-\lambda_0)$ & $288$ & $2$ &
$2B/(2B+\kappa^3)$; product over the primitive $\kappa$: $B^4/(B^2+2)^2$ \\
\hline
\end{tabular}
\caption{Order $k$ and leading coefficient of $R^w_2$ along $C_p$ at the corank-one labels
(Lemma~\ref{d6-lem-corank}), at one representative state of each orbit.  Counts are for
$a\ne4c\ne b$; $A=a/c$, $B=b/c$; $\lambda_0=\frac12(\sqrt6+i\sqrt2)$ is the primitive ratio for
$\omega=\frac12(-1+i\sqrt3)$.  In the first row $R^w_4$ has the same order.}\label{d6-tab-corank}
\end{table}

{
\begin{proof}
The assertions in (i) and (ii) are verified by exact symbolic computation
at one representative state of each of the eight corank-one orbits.
The coefficient field is $\mathbb Q(i)(c,d,e)$ when both wings are Gaussian;
otherwise it is extended by a root $\kappa$ of $\kappa^4-2\kappa^2+4$.
The computations are given in Appendix~\ref{app:verification}.

For (i), eliminating the diagonal derivative block of the distance rows
reduces the determinant to a Schur complement.
When both wings are Gaussian, its determinant, after multiplying the
rows by $d^3e^3$, is
$82944\,d^{61}e^{59}(cd^3-\frac{\lambda}{16})(ce^3-\frac{\kappa}{16})$.
The first factor is nonzero because
$cd^3-\frac{\lambda}{16}=-\frac{\lambda a}{8(2a+c\lambda^3)}\ne0$;
the second is treated by exchanging the wings.

In the other cases, after removing the unit row for $t$, the complement
is block diagonal with one $5\times5$ block for each wing
(Lemma~\ref{d6-lem-chart}(v)).  The computation checks the vanishing of an
off-diagonal block and evaluates the two diagonal blocks.
For one Gaussian and one primitive wing, the determinants are nonzero
multiples of $8cd^3-\frac{\lambda}{2}=-\lambda A/(2A+\lambda^3)$ and of
a polynomial in $ce^3$ and $\kappa$.  The product of the latter polynomial
over the primitive $\kappa$ is
$(\frac9{16})^4\cdot64B^4(B^4-8B^2+64)/(B^2+2)^4$.
Its nonvanishing follows from $B^4-8B^2+64=(B^2-4)^2+48$.
This product, and those in Table~\ref{d6-tab-corank}, are evaluated using
resultants with $\kappa^4-2\kappa^2+4$.

For (ii), the Taylor coefficients of $C_p$ are found by solving the
$20$ rows in (i) successively in powers of $t$.  Substitution into all
$22$ rows checks the resulting expansions and gives the first nonzero
coefficient listed in Table~\ref{d6-tab-corank}.

Finally, (iii) is the analytic consequence of these two computations.
Near $p$, (i) gives $F^{-1}(0)\subset C_p$.
On this curve the expansion in (ii) shows that
$R^w_2$ vanishes only at $t=0$, which is the point $p$.
Thus $p$ is isolated in the full zero set.
\end{proof}

Together with Remark~\ref{d6-rem-regular}, the lemma shows that every
label is isolated in $F^{-1}(0)$.  When $a\ne4c\ne b$, there are
$3888+1296=5184$ such labels.
}

\subsection{Exclusion}

\begin{theorem}\label{thm:d6}
Let $m\in\Rp^5$.  No curve in $A_m$ has a Diagram-6 germ, for any numbering of the bodies and
either colouring.  In particular the case $r=6$ of Theorem~\ref{thm:exclusion}(a) holds, at the
complex level.
\end{theorem}

\begin{proof}
{
The symmetries at the beginning of the section allow us to use the
numbering above.  Suppose that $\gamma(s)$ is a Diagram-6 germ on a curve
in $A_m$.  Proposition~\ref{prop:d6-cover} gives a chart germ
$\tilde\gamma(s)$ with $t(s)\to0$, $t(s)\ne0$ for $s\ne0$, bounded
coordinates, and a limit $p\in\mathcal B$.

Proposition~\ref{prop:d6-first} forces $m=(a,a,b,b,c)$ and
$\eta_{14}=-\eta_{23}$, $\eta_{13}=-\eta_{24}$ at $p$.
Lemma~\ref{d6-lem-atlas} therefore identifies $p$ as a label of corank one.
The computer-assisted Lemma~\ref{d6-lem-corank} shows that this point is
isolated in $F^{-1}(0)$.

On the other hand, the coordinates of $\tilde\gamma$ are bounded Puiseux
series in $s$.  After a substitution $s=\sigma^N$, they are convergent
power series in $\sigma$.  Thus $\sigma\mapsto\tilde\gamma(\sigma^N)$
is a holomorphic map from a disc into $F^{-1}(0)$, taking the value $p$
at $\sigma=0$.  Since $p$ is isolated, the map is constant near
$\sigma=0$.  This contradicts $t\not\equiv0$.
}
\end{proof}

{
Lemma~\ref{d6-lem-atlas} describes the whole boundary fibre $\mathcal B$,
including the faces $a=4c$, $b=4c$.  The argument therefore applies to
every crossed sign class and to every ramification of the germ.
}

\begin{remark}\label{d6-rem-projective}
{
On the face $a=4c$, the wing class $(\omega,\lambda)=(1,-2)$ has no finite
scale: it lies at $D_{12}=\infty$.  It cannot occur as the limit used in
Theorem~\ref{thm:d6}, since Proposition~\ref{prop:d6-cover} gives finite
nonzero limits for $D_{12}$ and $D_{34}$ along every Diagram-6 germ.

For completeness, Appendix~\ref{app:verification} also considers a
projective chart for one wing with $h=1/D_{12}$ and, on $a=b=4c$, a
projective ratio chart for both wings.  There, no analytic branch through
a point of $\{t=h=0\}$ leaves the divisor $h=0$.  This conclusion does not establish that these boundary points are
isolated.  The projective calculation is not required for
the proof above.
}
\end{remark}



\section{Diagram 7}\label{sec:d7}

Diagram~7 is the butterfly without circles in \cite[\S7.4]{AK}.  Albouy and Kaloshin
derive the mass relations (36), (39) and (40) for this diagram in
\cite[\S\S8.7.1--8.7.3]{AK}.  The argument below applies independently of these relations.

\begin{theorem}\label{thm:d7}
Let $m\in\Rp^5$.  No Diagram~7 germ with an unbounded four-product lies on a curve through a
real point: if an irreducible curve $\Gamma\subset A_m$ carries such a germ, then $\Gamma$
contains no real point.
\end{theorem}

This proves part~(b) of Theorem~\ref{thm:exclusion} for $r=7$.  The proof separates the
germs according to the sum of the two wing orders.  When $a+b\ge0$ in
Proposition~\ref{d7-cover}, every four-product is bounded (Lemma~\ref{d7-bounded}).  The
remaining germs lie in the open cone of the nested chart introduced below.  There an
identity on the solution set of \eqref{eq:system} implies $I\to0$
(Lemma~\ref{d7-vanish}).  Constancy of the inertia then gives $I=U\equiv0$ on $\Gamma$,
contradicting positivity at a real point.

\subsection{The diagram and the nested chart}
In the numbering of \cite[\S8.7]{AK}, the \emph{upper wing} $\{1,2,5\}$ is an isolated triangle
of $z$-edges without $z$-circles, the \emph{lower wing} $\{3,4,5\}$ is an isolated triangle of
$w$-edges without $w$-circles, and the four \emph{crossed} pairs $kj$ ($k\in\{1,2\}$,
$j\in\{3,4\}$) are not edges.
Since Diagram~7 has no circles, Rule~1e shows that neither wing has a maximal stroke (see the
proof of Proposition~\ref{d7-cover}(i)).  By Definition~\ref{def:germ} and the convention after
Theorem~\ref{thm:exclusion}, we may use this numbering and colouring.

We first normalize the orders as in Lemma~\ref{lem:germ}.  After renormalization,
$Z$ and $W$ have a common order.  Choosing a power $t$ of the Puiseux parameter makes this
order $-2$.  Thus $t$ plays the role of $\epsilon$ in \cite[\S3]{AK}, and $\ord$ denotes
order in $t$.  The renormalizations $(z,w)\mapsto(\lambda z,\lambda^{-1}w)$ preserve the
diagram, the ties, $I$, $U$ and every $r_{kl}^2=z_{kl}w_{kl}$.

To describe the two clusters at their different scales, we use the nested chart
with exponents $P,Q,R$ and coordinates $A_1,A_2,C,u_3,u_4$,
$B_3,B_4,D,v_1,v_2$ and $D_{kl}$ ($k<l$):
\begin{gather*}
  z_1=A_1,\quad z_2=A_2,\qquad w_3=t^{-Q}B_3,\quad w_4=t^{-Q}B_4,\\
  z_j=C+t^Pm_5u_j\ \ (j=3,4),\qquad z_5=C-t^P(m_3u_3+m_4u_4),\\
  w_k=t^{-Q}D+t^Rm_5v_k\ \ (k=1,2),\qquad w_5=t^{-Q}D-t^R(m_1v_1+m_2v_2),
\end{gather*}
and $\delta_{kl}=t^{e_{kl}}D_{kl}$, where $e_{kl}=-R/2$ on the upper wing, $Q/2$ on the crossed
pairs and $(Q-P)/2$ on the lower wing ($D$ without indices is not an inverse distance).
Here $C$ and $t^{-Q}D$ are the centres of mass of the $z$-cluster $\{3,4,5\}$ and
the $w$-cluster $\{1,2,5\}$.  The exponents $P$ and $R$ describe the internal separations
of these clusters, while $-Q$ describes the common $w$-scale.  The choices of $e_{kl}$
make the two sides of all ten ties $\delta_{kl}^2z_{kl}w_{kl}=1$ have the same order.
For the pairs $12$ and $34$, the ties become
\[
  m_5D_{12}^2(A_2-A_1)(v_2-v_1)=1,\qquad m_5D_{34}^2(u_4-u_3)(B_4-B_3)=1.
\]
The \emph{open cone} is $P>Q>0$, $R>0$.

We leave the rotational freedom in this chart unfixed.  In particular, the gauge
$z_{12}=w_{12}$ is not imposed.  The displayed ties show that $A_2-A_1$ and $B_4-B_3$
never vanish, so either could be used to normalize the rotations and the rescalings of
$t$.  Such a normalization is unnecessary here, since $I$, $U$, the ties and the solution
set of \eqref{eq:system} are invariant under the rotations.

\subsection{Orders of the germs and the chart}

\begin{proposition}[Representation in the nested chart]\label{d7-cover}
Let $\gamma$ be a Diagram~7 germ.  There are $a,b$ with $-2<a<1$ and $-2<b<1$ such that the
following hold.
\begin{romanlist}
  \item On the upper wing $\ord z_{kl}=a$ and $\ord w_{kl}=(4-a)/3$; on the lower wing
  $\ord w_{kl}=b$ and $\ord z_{kl}=(4-b)/3$; on the crossed pairs $\ord z_{kl}=a$ and
  $\ord w_{kl}=b$.  Moreover $\ord z_k\ge a$ and $\ord w_k\ge b$ for every $k$.
  \item After the rotation $(z,w)\mapsto(t^{-a}z,t^{a}w)$, $\gamma$ is a germ of the nested
  chart with exponents
  \[
    P=\frac{4-b}{3}-a,\qquad Q=-(a+b),\qquad R=\frac{4+2a}{3};
  \]
  all chart coordinates have nonnegative order, so they converge, and every $D_{kl}$ has order
  zero.  Moreover $R=\frac23(a+2)>0$ and $P-Q=\frac23(b+2)>0$, so $(P,Q,R)$ lies in the open
  cone if and only if $a+b<0$.
\end{romanlist}
\end{proposition}

\begin{proof}
(i) We use Fact~\ref{fact:rules}.  By Rule~2e \cite[p.~552]{AK} the three $z$-separations of a
wing have one order, and so do its three $w$-separations; let $a$ be the $z$-order of the upper
wing and $b$ the $w$-order of the lower wing.  As $\delta_{kl}^2z_{kl}w_{kl}=1$,
$\ord Z_{kl}=-(\ord z_{kl}+3\ord w_{kl})/2$ and $\ord W_{kl}=-(3\ord z_{kl}+\ord w_{kl})/2$;
since the strokes have order $-2$, this gives the other two wing orders.  The upper wing is an
isolated component of the $z$-diagram without $z$-circles, so it has no maximal $z$-stroke
(Rule~1e, p.~546), and along a $z$-stroke Estimate~2 (p.~548) gives $-2\le\ord z_{kl}\le1$, with
$-2$ exactly for a maximal $z$-edge and $1$ exactly for a $zw$-edge.  Hence $-2<a<1$, and dually
$-2<b<1$.  As $z_{13}=z_{15}+z_{53}$
with $\ord z_{53}=(4-b)/3>1>a$, and $w_{13}=w_{15}+w_{53}$ with $\ord w_{15}=(4-a)/3>1>b$, the
crossed orders are $a$ and $b$; likewise for $14$, $23$, $24$.  Finally
$\sum_km_kz_k=\sum_km_kw_k=0$ on $A_m$ (proof of Lemma~\ref{lem:distances}), so each $z_k$ is a
combination of $z$-differences, all of order at least $a$; likewise $\ord w_k\ge b$.

(ii) After the rotation the $z$-differences have order $0$ on the upper wing and the crossed
pairs and $P$ on the lower wing, the $w$-differences have order $R$ on the upper wing and $-Q$
elsewhere, and $\ord z_k\ge0$, $\ord w_k\ge-Q$.  Inverting the chart map, $A_k=z_k$,
$B_j=t^Qw_j$, $C$ and $t^{-Q}D$ are the centres of the two clusters, $m_5t^Pu_j=z_j-C$ and
$m_5t^Rv_k=w_k-t^{-Q}D$ are combinations of differences inside them, and
$D_{kl}=t^{-e_{kl}}\delta_{kl}$, where $\ord\delta_{kl}=-(\ord z_{kl}+\ord w_{kl})/2=e_{kl}$.  So
every coordinate has nonnegative order and $\ord D_{kl}=0$.
\end{proof}

The condition $a+b<0$ therefore characterizes germs in the open cone.
As shown in Lemma~\ref{d7-bounded} below, germs with $a+b\ge0$ have no unbounded
four-product.  The case $a+b<0$ includes what \cite[\S8.7]{AK} call a wing in the
general case, but the argument here uses only the inequality itself.

\subsection{The inertia identity}
Write
\[
  F^z_k=z_k-\sum_{l\ne k}m_l\delta_{kl}^3(z_k-z_l),\qquad
  F^w_k=w_k-\sum_{l\ne k}m_l\delta_{kl}^3(w_k-w_l)
\]
for the coordinate rows of \eqref{eq:system} and $T_{kl}=\delta_{kl}^2z_{kl}w_{kl}-1$ for the
ties.  Then $I-U=\sum_km_kw_kF^z_k+\sum_{k<l}m_km_l\delta_{kl}T_{kl}$, so $U=I$ on $A_m$
\cite[p.~541]{AK}.

\begin{lemma}[Inertia identity]\label{d7-identity}
Let
\[
  S_5=\sum_{j=3,4}\sum_{l=1,2}m_jm_l\delta_{jl}^3
  \bigl((z_j-z_5)(w_l-w_5)-(w_j-w_5)(z_l-z_5)\bigr).
\]
Then, as polynomials in the masses and in $(z,w,\delta)$,
\begin{multline*}
  I=\sum_{k=1,2}m_kz_k(w_k-w_5)+\sum_{j=3,4}m_jw_j(z_j-z_5)+S_5\\
  -\sum_{j=3,4}m_j\bigl((z_j-z_5)F^w_j-(w_j-w_5)F^z_j\bigr)+w_5\sum_{k=1}^5m_kF^z_k.
\end{multline*}
In particular $I=\sum_{k=1,2}m_kz_k(w_k-w_5)+\sum_{j=3,4}m_jw_j(z_j-z_5)+S_5$ on $A_m$.
\end{lemma}

\begin{proof}
Put $\hat z_k=z_k-z_5$ and $\hat w_k=w_k-w_5$.  By antisymmetry of the force terms,
$\sum_km_kF^z_k=\sum_km_kz_k$.  Since $z_j\hat w_j-w_j\hat z_j=z_5w_j-w_5z_j$,
\begin{align*}
  I&=\sum_km_kz_k\hat w_k+w_5\sum_km_kz_k\\
   &=\sum_{k=1,2}m_kz_k\hat w_k+\sum_{j=3,4}m_j\bigl(w_j\hat z_j+z_5w_j-w_5z_j\bigr)
    +w_5\sum_km_kF^z_k.
\end{align*}
For $j\in\{3,4\}$, using $z_j-z_l=\hat z_j-\hat z_l$ and $w_j-w_l=\hat w_j-\hat w_l$,
\[
  \hat z_jF^w_j-\hat w_jF^z_j=(w_5z_j-z_5w_j)
  +\sum_{l\ne j}m_l\delta_{jl}^3(\hat z_j\hat w_l-\hat w_j\hat z_l).
\]
In $\sum_{j=3,4}m_j(\hat z_jF^w_j-\hat w_jF^z_j)$ the terms with $l=5$ vanish, those with
$l\in\{3,4\}$ cancel in pairs, and those with $l\in\{1,2\}$ add up to $S_5$.  Hence
$\sum_{j=3,4}m_j(z_5w_j-w_5z_j)=S_5-\sum_{j=3,4}m_j(\hat z_jF^w_j-\hat w_jF^z_j)$, and
substituting this into the display gives the identity.
\end{proof}

The term $S_5$ is the torque about body~$5$ of the crossed forces on the lower wing; compare
\cite[(31)]{AK}.

\begin{lemma}\label{d7-vanish}
Let $\gamma$ be a Diagram~7 germ on $A_m$ with $a+b<0$ (Proposition~\ref{d7-cover}).  Then
$I(\gamma(t))\to0$ as $t\to0$.
More precisely, in the nested chart of Proposition~\ref{d7-cover}(ii) the three terms
$\sum_{k=1,2}m_kz_k(w_k-w_5)$, $\sum_{j=3,4}m_jw_j(z_j-z_5)$ and $S_5$ are sums of powers $t^e$
times polynomials in the chart coordinates, with exponents
\[
  e=R;\qquad e=P-Q;\qquad e\in\bigl\{\tfrac Q2,\ P+\tfrac Q2,\ \tfrac{3Q}2+R,\ P+\tfrac{3Q}2+R\bigr\},
\]
respectively.
\end{lemma}

\begin{proof}
On $A_m$ the last two sums in Lemma~\ref{d7-identity} vanish, and the three remaining terms are
invariant under the rotations, so we may evaluate them on the rotated germ in the chart.  There
\begin{gather*}
  w_k-w_5=t^R\bigl((m_k+m_5)v_k+m_{k'}v_{k'}\bigr),\qquad
  z_j-z_5=t^P\bigl((m_j+m_5)u_j+m_{j'}u_{j'}\bigr),\\
  w_j-w_5=t^{-Q}(B_j-D)+t^R(m_1v_1+m_2v_2),\\
  z_l-z_5=A_l-C+t^P(m_3u_3+m_4u_4),
\end{gather*}
where $\{k,k'\}=\{1,2\}$, $\{j,j'\}=\{3,4\}$ and $l\in\{1,2\}$, and
$\delta_{jl}^3=t^{3Q/2}D_{jl}^3$.  Substituting, together with $z_k=A_k$ and $w_j=t^{-Q}B_j$,
gives the stated exponents.  Since $a+b<0$, Proposition~\ref{d7-cover}(ii) puts $(P,Q,R)$ in the
open cone, so every exponent is positive, and all chart coordinates are bounded.  Hence the three
terms, and with them $I$, tend to zero.
\end{proof}

\subsection{Germs outside the open cone}

\begin{lemma}\label{d7-bounded}
Let $\gamma$ be a Diagram~7 germ with $a+b\ge0$ (Proposition~\ref{d7-cover}).  Then every
four-product tends to zero along $\gamma$.  If $a+b=0$, then $-1<a<1$, and $p_{12}$, $p_{34}$
and each of the other eight four-products have the orders $(16-4a)/3$, $(16+4a)/3$ and $8/3$.
If $a+b>0$, every $r_{kl}$ tends to zero, and so does $I$.
\end{lemma}

\begin{proof}
By Proposition~\ref{d7-cover}(i), $r_{kl}^2=z_{kl}w_{kl}$ has order $(4+2a)/3>0$ on the upper
wing, $(4+2b)/3>0$ on the lower wing and $a+b\ge0$ on the crossed pairs.  The four-product
$p_{ij}$ contains the three squared sides of the triangle $\{k,l,m\}$ complementary to
$\{i,j\}$, and this triangle contains a wing pair: if $5\in\{k,l,m\}$, both pairs through~$5$
are wing pairs; otherwise $\{k,l,m\}\subset\{1,2,3,4\}$ contains $\{1,2\}$ or $\{3,4\}$ (the
crossed pairs form the bipartite graph $K_{2,2}$).  Hence
$\ord p_{ij}\ge\min\{(4+2a)/3,(4+2b)/3\}>0$.  If $a+b=0$, the bounds $-2<a,b<1$ give $-1<a<1$,
and adding orders gives the stated values.  If $a+b>0$, every $r_{kl}^2$ has positive order, and
the Lagrange identity $MI=(\sum_km_kz_k)(\sum_km_kw_k)+\sum_{k<l}m_km_lz_{kl}w_{kl}$, with
$\sum_km_kz_k=\sum_km_kw_k=0$, gives $I\to0$.
\end{proof}

\begin{remark}\label{d7-balanced}
Germs with $a+b=0$ carry no unbounded four-product.  They therefore fall outside
the conclusion of Proposition~\ref{prop:reduction}, and their exclusion is unnecessary
for the present proof.
\end{remark}

\subsection{Conclusion of the Diagram 7 argument}

\begin{proof}[Proof of Theorem~\ref{thm:d7}]
Let $\gamma$ be a Diagram~7 germ with an unbounded four-product on an irreducible curve
$\Gamma\subset A_m$, numbered as above, and let $a,b$ be as in Proposition~\ref{d7-cover}.  Some
four-product has negative order along $\gamma$ (Definition~\ref{def:germ}), so $a+b<0$ by
Lemma~\ref{d7-bounded}, and $I(\gamma(t))\to0$ by Lemma~\ref{d7-vanish}.  By
Lemma~\ref{lem:U-constant}, $I=U$ is constant on $\Gamma$, with value
$\lim_{t\to0}I(\gamma(t))=0$.  At a real point $w_k=\bar z_k$, so $I=\sum_km_k|z_k|^2$,
which is positive because $m_k>0$ and the tie of the pair $12$ forces $z_1\ne z_2$.  Hence
$\Gamma$ contains no real point.
\end{proof}

\begin{remark}\label{d7-masses}
Theorem~\ref{thm:d7} holds for all positive masses.  The inertia identity in
Lemma~\ref{d7-identity} is polynomial; the order estimates in
Proposition~\ref{d7-cover} and Lemma~\ref{d7-bounded} follow from Rules~1e, 2e and
Estimate~2 of \cite[\S3]{AK}.  Lemma~\ref{d7-vanish} then uses the positivity of the
exponents, and the conclusion follows from Lemma~\ref{lem:U-constant} and $m_k>0$.
Thus the mass relations (36), (39) and (40) play no role in this argument, which derives
no additional mass condition for Diagram~7.
\end{remark}



\section{Diagram 8}\label{sec:d8}

Diagram~8 is the butterfly of \cite[\S7.4]{AK} with exactly one circled wing
(Fact~\ref{fact:fig11}).  We use the numbering of \cite[\S8.8]{AK}: the upper wing $\{1,2,5\}$
is an isolated triangle of $z$-edges without $z$-circles, the lower wing $\{3,4,5\}$ is a
triangle of $w$-edges with $w$-circles at $3$ and $4$ (vertex $5$ is never circled
\cite[\S7.4]{AK}, and Rules~1a, 1d place the circles of a circled wing at its other two
vertices), and the crossed pairs $13$, $14$, $23$, $24$ are not edges.
The other circling becomes this one after exchanging $z$ and $w$ and renumbering
$1\leftrightarrow3$, $2\leftrightarrow4$; both operations are symmetries of \eqref{eq:system}
(the second permutes the masses), so it suffices to treat this numbering for every
$m\in\Rp^5$.
The mass condition of the diagram is relation (40) of \cite[\S8.8]{AK},
\begin{equation}\label{eq:d8-40}
  m_5^2(m_1+m_2)^4=m_1^3m_2^3,
\end{equation}
which \cite[p.~578]{AK} obtain ``by the same arguments'' as for Diagram~7.  For the other
circling it reads $m_5^2(m_3+m_4)^4=m_3^3m_4^3$, which is their (39).
The proof begins with a chart in which the lower wing has a finite limit.
Every Diagram~8 germ is represented by a ray in this chart
(Proposition~\ref{prop:d8-cover}), and its limiting equations imply
\eqref{eq:d8-40} (Proposition~\ref{prop:d8-first}).  We then use the
{exact distance equations on the upper wing and the coefficients of order nine
to restrict the possible germs to one curve of}
masses.  A second-order expansion excludes the remaining germs, including complex
ones (Theorem~\ref{thm:d8}).

We shall also use two symmetries to normalize the masses and coordinates.  Scaling the
masses by $\lambda>0$ and $(z,w,\delta)$ by
$(\lambda^{1/3},\lambda^{1/3},\lambda^{-1/3})$ maps $A_m$ onto $A_{\lambda m}$.
Every coordinate and every component of $\mathcal Z$, $\mathcal W$ is multiplied by a
nonzero constant, so the orders and diagrams are preserved (Lemma~\ref{lem:germ}).
The equations of \eqref{eq:system}, apart from $z_{12}=w_{12}$, are also invariant under
$z\mapsto cz$, $w\mapsto c^{-1}w$.  This changes the $z$-orders by $\ord c$ and the
$w$-orders by $-\ord c$, preserving the diagram \cite[p.~547]{AK}.

\subsection{The ratio chart}
Let $\gamma$ be a Diagram~8 germ.  We rotate it by $z_{12}^{-1}$, so that $z_{12}=1$ along
$\gamma$.  The \emph{ratio chart} has two parameters $t,h$ and nineteen coordinates, the
positions and shapes $A_1,A_2,C,B_3,B_4,D,u,v,R_{12},R_{25},P,Q$ and the inverse distances
$D_{13},D_{14},D_{23},D_{24},D_{34},D_{35},D_{45}$; we write $x$ for all of them.  It is given by
\begin{gather*}
  z_1=A_1,\quad z_2=A_2,\quad z_3=C+m_5t^6u,\quad z_4=C+m_5t^6v,\\
  z_5=C-t^6(m_3u+m_4v),\quad w_3=t^{-6}B_3,\quad w_4=t^{-6}B_4,\\
  w_1=t^{-6}D+m_5t^4h^{-2}P,\quad w_2=t^{-6}D+m_5t^4h^{-2}Q,\\
  w_5=t^{-6}D-t^4h^{-2}(m_1P+m_2Q),\quad \delta_{15}=t^{-2}h,\\
  \delta_{12}=t^{-2}hR_{12},\quad \delta_{25}=t^{-2}hR_{25},\\
  \delta_{bc}=t^3D_{bc}\ \ (b\in\{1,2\},\ c\in\{3,4\}),\qquad
  \delta_{kl}=D_{kl}\ \ (k,l\in\{3,4,5\}).
\end{gather*}
{On the gauge slice defined below, we take a local sixth root $t$ of
$(w_4-w_3)^{-1}$.  Once this root is fixed and $th\ne0$, the inverse chart map is uniquely
determined: $C$ and $t^{-6}D$ are the centres of mass of $z_3,z_4,z_5$ and of}
$w_1,w_2,w_5$, $h=t^2\delta_{15}$, $R_{12}=\delta_{12}/\delta_{15}$,
$R_{25}=\delta_{25}/\delta_{15}$, $u=(z_3-C)/(m_5t^6)$, $P=h^2(w_1-t^{-6}D)/(m_5t^4)$, and
{the remaining coordinates are recovered in the same way.  Along a germ,
a finite ramification allows all these roots to be chosen consistently.}
The rows of the chart are the gauges $E_{20}:=A_2-A_1-1$ and $E_{21}:=B_4-B_3-1$ (that is,
$z_{12}=1$ and $w_4-w_3=t^{-6}$) and $E_0,\dots,E_{19}$, obtained by substituting the chart into
the twenty equations of \eqref{eq:system} other than $z_{12}=w_{12}$ and multiplying each by a
monomial in $t$ and $h^{1/2}$: $E_{2k-2}$ and $E_{2k-1}$ come from the $z$- and the $w$-equation
of body $k$, and $E_{10},\dots,E_{19}$ from the ties of the pairs
$12,13,14,15,23,24,25,34,35,45$.
In the corner coordinates
\[
  a:=th^{-1/2},\qquad \eta:=t^2h^{1/2}\qquad(\text{so } t^3=a\eta,\ h^{3/2}=\eta\,a^{-2})
\]
we use the monomials
\[
  \mu:=a^6=t^6h^{-3},\qquad \nu:=a^2\eta^2=t^6,\qquad T:=a\eta=t^3,\qquad \eta^2=t^4h,
\]
and we write $X:=A_1-C$, $Y:=A_2-C$, $M_{12}:=m_1+m_2$, $M:=m_3+m_4+m_5$, and
$C_5:=C-\nu(m_3u+m_4v)$, which is $z_5$.

{The chart parameters may be complex.  In asymptotic estimates we therefore
interpret sums of scales as sums of absolute values: $O(\mu+\rho^2)$ means
$O(|\mu|+|\rho|^2)$, and $o(\mu+\rho^2)$ means $o(|\mu|+|\rho|^2)$.
Products inside $O$ and $o$ are interpreted in the same way.  Algebraic identities and
equations between leading coefficients retain their usual complex meaning.}

{\begin{lemma}\label{lem:d8-rows}
The chart equations have the following polynomial identities, valid for all masses.
\begin{romanlist}
\item Every $E_j$ is a polynomial in $a$, $\eta$ and $x$, and its \emph{face}
$f_j:=E_j|_{a=\eta=0}$ is
\begin{align*}
  f_0&=-\bigl(m_2(X-Y)R_{12}^3+m_5X\bigr), & f_8&=m_1X+m_2YR_{25}^3,\\
  f_2&=-\bigl(m_1(Y-X)R_{12}^3+m_5YR_{25}^3\bigr), & f_{20}&=A_2-A_1-1,\\
  f_{10}&=m_5(X-Y)(P-Q)R_{12}^2-1, & f_{13}&=X\bigl((m_1+m_5)P+m_2Q\bigr)-1,\\
  f_{16}&=YR_{25}^2\bigl(m_1P+(m_2+m_5)Q\bigr)-1, & f_4&=f_6=C,\quad f_1=f_3=D,\\
  f_{11}&=X(D-B_3)D_{13}^2-1, & f_{12}&=X(D-B_4)D_{14}^2-1,\\
  f_{14}&=Y(D-B_3)D_{23}^2-1, & f_{15}&=Y(D-B_4)D_{24}^2-1,\\
  f_5&=B_3+m_4(B_4-B_3)D_{34}^3-m_5(B_3-D)D_{35}^3, & f_{21}&=B_4-B_3-1,\\
  f_7&=B_4-m_3(B_4-B_3)D_{34}^3-m_5(B_4-D)D_{45}^3,\\
  f_9&=D+m_3(B_3-D)D_{35}^3+m_4(B_4-D)D_{45}^3,\\
  f_{17}&=m_5(B_3-B_4)(u-v)D_{34}^2-1,\\
  f_{18}&=(B_3-D)\bigl((m_3+m_5)u+m_4v\bigr)D_{35}^2-1,\\
  f_{19}&=(B_4-D)\bigl(m_3u+(m_4+m_5)v\bigr)D_{45}^2-1.
\end{align*}
\item $m_1E_0+m_2E_2+m_5E_8+a^6(m_3E_4+m_4E_6)=a^6(m_1A_1+m_2A_2+MC)$.
\item $E_1=D+\eta^2\beta_1+T^3\tau_1+\sum_{i\ge6,\,j\ge2}a^i\eta^j(\cdots)$, and $E_3$ is the
same with $\beta_3,\tau_3$, where
\begin{align*}
  \beta_1&=-m_5\bigl(m_2(P-Q)R_{12}^3+(m_1+m_5)P+m_2Q\bigr),\\
  \beta_3&=-m_5\bigl(m_1(Q-P)R_{12}^3+(m_1P+(m_2+m_5)Q)R_{25}^3\bigr),\\
  \tau_1&=m_3(B_3-D)D_{13}^3+m_4(B_4-D)D_{14}^3,\\
  \tau_3&=m_3(B_3-D)D_{23}^3+m_4(B_4-D)D_{24}^3;
\end{align*}
the monomials of $E_1-E_3$ are $\eta^2$, $T^3$, $a^6\eta^2$ and $a^9\eta^5$.
\item For $k\in\{0,2,8,10,13,16,20\}$, $E_k=f_k(\mathcal I)+\sum_{i\ge6}a^i\eta^j(\cdots)$, where
$f_k$ is read as a function of $(X,Y,R_{12},R_{25},P,Q)$ and
$\mathcal I:=(A_1-C_5,\,A_2-C_5,\,R_{12},\,R_{25},\,P,\,Q)$.
\item $E_4=C+\nu L_3+T^3F_3+T^5G_3$ exactly, where
$L_3=-m_5\bigl(m_4(u-v)D_{34}^3+((m_3+m_5)u+m_4v)D_{35}^3-u\bigr)$,
$F_3=-\bigl(m_1(C-A_1)D_{13}^3+m_2(C-A_2)D_{23}^3\bigr)$ and
$G_3=-m_5u\bigl(m_1D_{13}^3+m_2D_{23}^3\bigr)$.
\end{romanlist}
\end{lemma}}
{\begin{proof}
We write out the chart equations to identify both their limiting terms and the
monomials that enter the later order comparisons.
Put $I=\{1,2\}$, $J=\{3,4\}$, $P_1=P$, $P_2=Q$, $u_3=u$, $u_4=v$,
and set
\[
q=R_{12},\quad s=R_{25},\quad V=m_3u+m_4v,\quad S=m_1P+m_2Q,
\quad \gamma=\mu\eta^2=\nu a^4.
\]
The unscaled positions are
\[
z_i=A_i,\quad z_j=C+\nu m_5u_j,\quad z_5=C_5=C-\nu V,
\]
whereas their scaled $w$-coordinates are
\[
\nu w_i=D+\gamma m_5P_i,\qquad
\nu w_j=B_j,\qquad \nu w_5=D-\gamma S.
\]
The inverse-distance cubes are $\mu^{-1}q^3,\mu^{-1},\mu^{-1}s^3$
on $12,15,25$, respectively, $T^3D_{ij}^3$ on crossed pairs, and
$D_{kl}^3$ on the lower wing.
The chart equations are obtained by multiplying the $z$-equations of $1,2,5$
by $\mu$, leaving those of $3,4$ unscaled, and multiplying all $w$-equations by $\nu$.
Writing $X'=A_1-C_5$ and $Y'=A_2-C_5$, the resulting upper $z$-rows are
\begin{align*}
E_0={}&\mu A_1-m_2(A_1-A_2)q^3-m_5X'
 -\mu T^3\sum_{j\in J}m_j(A_1-C-\nu m_5u_j)D_{1j}^3,\\
E_2={}&\mu A_2-m_1(A_2-A_1)q^3-m_5Y's^3
 -\mu T^3\sum_{j\in J}m_j(A_2-C-\nu m_5u_j)D_{2j}^3,\\
E_8={}&\mu C_5+m_1X'+m_2Y's^3
 +\mu\nu\sum_{j\in J}m_j(V+m_5u_j)D_{j5}^3.
\end{align*}
If $j'$ denotes the other element of $J$, the lower $z$-rows are
\begin{align*}
E_{2j-2}={}&C+\nu m_5u_j
 -\nu m_{j'}m_5(u_j-u_{j'})D_{jj'}^3
 -\nu m_5(V+m_5u_j)D_{j5}^3\\
&-T^3\sum_{i\in I}m_i(C-A_i+\nu m_5u_j)D_{ij}^3.
\end{align*}
This formula at $j=3$ proves (v).

The upper $w$-rows are
\begin{align*}
E_1&=D+\eta^2\beta_1+T^3\tau_1
 +\gamma m_5P-\gamma T^3m_5P\sum_{j\in J}m_jD_{1j}^3,\\
E_3&=D+\eta^2\beta_3+T^3\tau_3
 +\gamma m_5Q-\gamma T^3m_5Q\sum_{j\in J}m_jD_{2j}^3.
\end{align*}
Thus (iii) follows, including its complete list of monomials.
The lower and central $w$-rows are
\begin{align*}
E_{2j-1}={}&B_j-m_{j'}(B_j-B_{j'})D_{jj'}^3-m_5(B_j-D)D_{j5}^3
 -\gamma m_5 S D_{j5}^3\\
&-T^3\sum_{i\in I}m_i(B_j-D)D_{ij}^3
 +\gamma T^3m_5\sum_{i\in I}m_iP_iD_{ij}^3,\\
E_9={}&D+\sum_{j\in J}m_j(B_j-D)D_{j5}^3
 +\eta^2\beta_5+\gamma S\Bigl(-1+\sum_{j\in J}m_jD_{j5}^3\Bigr),
\end{align*}
where
\[
\beta_5=m_1((m_1+m_5)P+m_2Q)
 +m_2s^3(m_1P+(m_2+m_5)Q).
\]

Finally the distance rows are obtained by subtracting $1$ from the
following expressions:
\begin{center}
\begin{tabular}{c l}
$12$ & $m_5(A_1-A_2)(P-Q)q^2$\\
$15$ & $X'((m_1+m_5)P+m_2Q)$\\
$25$ & $Y's^2(m_1P+(m_2+m_5)Q)$\\
$ij$ & $(A_i-C-\nu m_5u_j)(D-B_j+\gamma m_5P_i)D_{ij}^2$\\
$34$ & $m_5(u-v)(B_3-B_4)D_{34}^2$\\
$j5$ & $(V+m_5u_j)(B_j-D+\gamma S)D_{j5}^2$.
\end{tabular}
\end{center}
These formulas account for all ten distance rows; the last two rows are
$E_{20}=A_2-A_1-1$ and $E_{21}=B_4-B_3-1$.
Every expression is polynomial.  Setting $a=\eta=0$ gives the boundary
equations listed in (i).
The three upper distance rows and the gauge are already exact functions
of $\mathcal I$; in the upper $z$-rows all other terms have a factor $\mu$.
This proves (iv).  Finally, the mass-weighted sum of the original $z$-equations
has no force terms, since the two terms contributed by each pair cancel.
Multiplying this sum by $\mu$ gives (ii), because the weighted position of
the lower cluster is $MC$.
\end{proof}

}\begin{definition}\label{def:d8-ray}
A germ in the ratio chart is a \emph{ray germ} if along it $t\to0$, $x$ converges, $\ord a>0$
and $\ord\eta>0$.  Its \emph{ray} is $r:=\ord\eta/\ord a\in(0,\infty)$ and its \emph{exponent}
is $p:=\ord h/\ord t$.
\end{definition}
Since $t^3=a\eta$ and $h^{3/2}=\eta a^{-2}$, $p=(2r-4)/(r+1)$, an increasing bijection of
$(0,\infty)$ onto $(-4,2)$; the rays $r=\frac12,2,3$ are $p=-2,0,\frac12$, and the faces
$\{a=0\}$, $\{\eta=0\}$ of the corner chart are $p=-4$ and $p=2$.

\subsection{Representation of all Diagram 8 germs}

\begin{proposition}[Representation in the ratio chart]\label{prop:d8-cover}
Every Diagram~8 germ is a ray germ of the ratio chart with exponent $p\in(-4,2)$ and a finite
lower wing.  Precisely: let $\gamma$ be a Diagram~8 germ, rotated so that $z_{12}=1$, and put
$t:=(w_4-w_3)^{-1/6}$ and $h:=t^2\delta_{15}$ along $\gamma$.  Let $\theta$ be the order of
$z_{12}$ before the rotation, in the normalization $Z=W=\epsilon^{-2}$ with orders in $\epsilon$.
Then $-2<\theta<1$; every coordinate in $x$ converges along $\gamma$, and $R_{12}$, $R_{25}$,
$D_{13},D_{14},D_{23},D_{24}$, $D_{34},D_{35},D_{45}$ have nonzero limits; and
\begin{gather*}
  \ord t=\tfrac{2-\theta}6,\qquad \ord h=-\tfrac{2\theta}3,\qquad
  p=-\tfrac{4\theta}{2-\theta},\\
  \ord a=\tfrac{2+\theta}6,\qquad \ord\eta=\tfrac{2-2\theta}3,\qquad
  r=\tfrac{4(1-\theta)}{2+\theta}.
\end{gather*}
\end{proposition}

\begin{proof}
By Fact~\ref{fact:rules} and Lemma~\ref{lem:germ} the rules and estimates of \cite[\S3]{AK}
hold for the orders along $\gamma$.
We begin with the lower wing.  The triangle $\{3,4,5\}$ is an isolated component of
the $w$-diagram with $w$-circles, so it has a maximal $w$-stroke
\cite[Rule~1e, p.~546]{AK}.  By Rule~2e its three
$w$-separations have one order, so all three $w$-strokes are maximal, and the right equality
case of Estimate~2 \cite[p.~548]{AK}, in the colour $w$, gives $\ord w_{kl}=-2$,
$\ord z_{kl}=2$, $\ord r_{kl}=0$ for $kl\in\{34,35,45\}$.
The upper wing $\{1,2,5\}$ is an isolated component of the $z$-diagram without
circles.  It has no maximal $z$-stroke by Rule~1e and no $w$-stroke, hence no $zw$-edge.
Thus
neither equality case of Estimate~2 holds, and by Rule~2e the three $z$-separations have one
order $\theta$ with $-2<\theta<1$.  The $z$-strokes have $\ord Z_{kl}=-2$, so
$\ord r_{kl}=(\theta+2)/3$ and $\ord w_{kl}=(4-\theta)/3$.
These wing orders determine the crossed separations.  For $b\in\{1,2\}$, $c\in\{3,4\}$, $z_{bc}=z_{b5}+z_{5c}$ has order
$\theta$ (as $\theta<2$) and $w_{bc}=w_{b5}+w_{5c}$ has order $-2$ (as $(4-\theta)/3>-2$), so
$\ord r_{bc}=(\theta-2)/2$.
We next recover the orders of the positions.  Every $w_k$ has order $\ge-2$, with equality for $k=3,4$ ($w$-circles).  Let
$C$ be the centre of mass of $z_3,z_4,z_5$.  Then $(m_1+m_2+M)C=-m_1(z_1-C)-m_2(z_2-C)$, and
$z_b-C=z_{5b}+(z_5-C)$ with $\ord z_{5b}=\theta$ and $\ord(z_5-C)\ge2$; so $C$, $z_1$, $z_2$
have order $\ge\theta$, and $z_k-C$ has order $\ge2$ for $k=3,4,5$.
After the rotation, all $z$-orders decrease by $\theta$ and all $w$-orders increase
by $\theta$.  Consequently $\ord(w_4-w_3)=\theta-2$ and
$\ord t=(2-\theta)/6>0$.  The inverse chart formulas now give the orders of the new
coordinates: $A_1,A_2,C$ have order $\ge0$; $u,v$ have order $\ge(2-\theta)-6\ord t=0$; $B_3,B_4$
have order $0$ and $D$ order $\ge0$; $\ord h=2\ord t-\ord r_{15}=-2\theta/3$; $R_{12},R_{25}$
have order $0$; $P,Q$ have order $\ge2\ord h+\frac{4+2\theta}3-4\ord t=0$, because
$w_1-t^{-6}D=(m_2w_{21}+m_5w_{51})/(m_1+m_2+m_5)$ and similarly for $w_2$; $D_{bc}=t^{-3}\delta_{bc}$
has order $\frac{2-\theta}2-3\ord t=0$; and $D_{kl}=\delta_{kl}$ has order $0$.  Hence every
coordinate converges (Lemma~\ref{lem:germ}), the inverse-distance coordinates to nonzero limits,
and $p=\ord h/\ord t=-4\theta/(2-\theta)\in(-4,2)$.  Finally
$\ord a=\ord t-\frac12\ord h=(2+\theta)/6$ and $\ord\eta=2\ord t+\frac12\ord h=(2-2\theta)/3$,
both positive.
\end{proof}
In particular, the lower inverse distances have finite nonzero limits.
A Diagram~8 germ therefore cannot approach the projective end at which these
coordinates degenerate.

\subsection{The mass relation and the limiting configuration}

\begin{proposition}[Mass relation and limiting coordinates]\label{prop:d8-first}
At positive masses every Diagram~8 germ satisfies \eqref{eq:d8-40}:
$m_5^2(m_1+m_2)^4=m_1^3m_2^3$.  Moreover, in the chart of Proposition~\ref{prop:d8-cover} its
limit $x_0$ has $A_1=-m_2/M_{12}$, $A_2=m_1/M_{12}$, $C=0$, $R_{25}=1$,
$R_{12}=-m_5M_{12}/(m_1m_2)$, $(P,Q)=(P_0,Q_0)$ the unique solution of $f_{10}=f_{13}=0$,
$D=0$, $B_3=-m_4/(m_3+m_4)$, $B_4=m_3/(m_3+m_4)$, and, with $L:=(M_{12}(m_3+m_4))^{1/2}>0$ and
signs $s_{bc}\in\{\pm1\}$,
\begin{gather*}
  D_{13}=\frac{is_{13}L}{\sqrt{m_2m_4}},\qquad D_{14}=\frac{s_{14}L}{\sqrt{m_2m_3}},\\
  D_{23}=\frac{s_{23}L}{\sqrt{m_1m_4}},\qquad D_{24}=\frac{is_{24}L}{\sqrt{m_1m_3}}.
\end{gather*}
\end{proposition}
We write $U:=(A_1,A_2,C,R_{12},R_{25},P,Q)$, $U_0$ for its limit, and
$\mathcal I_0:=(X_0,Y_0,R_{12}^0,1,P_0,Q_0)$, with $X_0=-m_2/M_{12}$, $Y_0=m_1/M_{12}$ and
$R_{12}^0=-m_5M_{12}/(m_1m_2)$, for the limit of $\mathcal I$ (note $C_5\to C\to0$).

\begin{proof}
Let $m\in\Rp^5$ and let $\gamma$ be a Diagram~8 germ, in the chart of
Proposition~\ref{prop:d8-cover}, with limit $x_0$.
By Lemma~\ref{lem:d8-rows}\,(i) and Proposition~\ref{prop:d8-cover},
$E_j=f_j(x)+O(|a|+|\eta|)$ along $\gamma$, so $f_j(x_0)=0$ for all $j$.
The centre-of-mass identity in Lemma~\ref{lem:d8-rows}\,(ii), together with
$a\ne0$, gives $m_1A_1+m_2A_2+MC=0$ along $\gamma$.  At the limit this identity,
$f_4=C$ and $f_{20}$ form a linear system with determinant $M_{12}$.
It follows that $A_1=-m_2/M_{12}$, $A_2=m_1/M_{12}$ and $C=0$ at $x_0$.
Thus $X=A_1$, $Y=A_2$ and $X-Y=-1$ there.

At these positions the upper force equations reduce to
$f_8=\frac{m_1m_2}{M_{12}}(R_{25}^3-1)$ and $f_0=\frac{m_2}{M_{12}}F$,
where $F:=M_{12}R_{12}^3+m_5$.  Hence $R_{25}^3=1$ and $F=0$.

The distance equations impose a further relation, expressing the closure of
the $w$-triangle $125$.  The identity
\begin{align*}
  &(X-Y)R_{12}^2Xf_{16}-(X-Y)R_{12}^2YR_{25}^2f_{13}+XYR_{25}^2f_{10}\\
  &\qquad=(X-Y)R_{12}^2(YR_{25}^2-X)-XYR_{25}^2,
\end{align*}
eliminates $P,Q$.  At $x_0$ its right-hand side is $-G/M_{12}^2$, where
$G:=M_{12}R_{12}^2(m_1R_{25}^2+m_2)-m_1m_2R_{25}^2$, so $G=0$.

To combine this relation with the force equations, set
$H:=M_{12}(m_1+m_2R_{25})R_{12}^2-m_1m_2$ and
$\Phi(s):=m_5^2M_{12}(m_1+m_2s)^3-(m_1m_2)^3$.  The identities
\begin{align*}
  H&=R_{25}G+(1-R_{25}^3)\,m_1(M_{12}R_{12}^2-m_2),\\
  \Phi(R_{25})&=H\bigl(H^2+3m_1m_2H+3m_1^2m_2^2\bigr)
  -M_{12}(m_1+m_2R_{25})^3F(F-2m_5),
\end{align*}
give $H=0$ and $\Phi(R_{25})=0$ at $x_0$.

Positivity of the masses excludes the nonreal cube roots.  Indeed, if
$R_{25}\ne1$, then $R_{25}^2+R_{25}+1=0$ and
$\Phi(R_{25})=\Phi_0+\Phi_1R_{25}$ with real $\Phi_0,\Phi_1$.  Since
$R_{25}\notin\R$, both coefficients must vanish: $\Phi_0=\Phi_1=0$.
However, $\Phi_1=3m_1m_2m_5^2M_{12}(m_1-m_2)$ forces $m_1=m_2$, after which
$\Phi_0=-m_1^4(m_1^2+2m_5^2)<0$.  Thus $R_{25}=1$, and $\Phi(1)=0$ is precisely
\eqref{eq:d8-40}.

The remaining limit coordinates are now determined successively.  With $R_{25}=1$,
the equations $F=0$ and $H=M_{12}^2R_{12}^2-m_1m_2=0$ give
$R_{12}=R_{12}^3/R_{12}^2=-m_5M_{12}/(m_1m_2)$.  The equations $f_{10},f_{13}$
are linear in $(P,Q)$ with determinant
$m_2m_5(m_1+m_2+m_5)R_{12}^2/M_{12}\ne0$, so their solution is unique.
Next, $f_1=D$ gives $D=0$.  The identity
$m_1f_1+m_2f_3+m_3f_5+m_4f_7+m_5f_9=(m_1+m_2+m_5)D+m_3B_3+m_4B_4$
and $f_{21}$ determine $B_3$, $B_4$.  Finally, the crossed distance equations
$f_{11},f_{12},f_{14},f_{15}$ yield
$D_{13}^2=-M_{12}(m_3+m_4)/(m_2m_4)$, $D_{14}^2=M_{12}(m_3+m_4)/(m_2m_3)$,
$D_{23}^2=M_{12}(m_3+m_4)/(m_1m_4)$ and $D_{24}^2=-M_{12}(m_3+m_4)/(m_1m_3)$,
which give the stated square roots.
\end{proof}
\begin{remark}\label{rem:d8-ak}
\cite[\S8.8]{AK} place the cluster $z_3,z_4,z_5$ at the centre of mass of $z_1,z_2$, invoke
their relation (34), $(m_1+m_2)W_{12}\sim m_5W_{51}+m_5W_{25}$, and refer to their \S8.7.3.  In
the chart $(m_1+m_2)W_{12}-m_5(W_{51}+W_{25})=t^{-2}h(\beta_1-\beta_3)=t^{-6}\eta^2
(\beta_1-\beta_3)$, and $\eta^2(\beta_1-\beta_3)$ is the $\eta^2$-part of $E_1-E_3$; so (34)
holds at the limit of every germ (Lemma~\ref{lem:d8-rows}\,(iii)), and given $R_{25}^3=1$, the
ties and the centre of mass it is equivalent to $F=0$.  The centre of mass cannot be dropped: at
$(m_1,m_2,m_5)=(4,1,\frac49)$, which violates \eqref{eq:d8-40}, the point
$A_1=R_{12}=-\frac23$, $A_2=\frac13$, $C=0$, $R_{25}=2$, with $P,Q$ from $f_{10}=f_{13}=0$,
satisfies every upper face and (34), but $m_1A_1+m_2A_2+MC=-\frac73$.  Relation (39) is not
used: at $D=0$, $B_3=-m_4/(m_3+m_4)$, $B_4=m_3/(m_3+m_4)$, with $D_{35}=g$, $D_{45}=\omega g$,
$D_{34}=\lambda g$, the face $f_5$ reads $g^3\bigl((m_3+m_4)\lambda^3+m_5\bigr)=1$, and the
escape $(m_3+m_4)\lambda^3+m_5=0$ of the lower scale forces $\omega=1$ and (39); by
Proposition~\ref{prop:d8-cover} no Diagram~8 germ escapes.
\end{remark}

\subsection{Exclusion}
{We first use the distance equations on the upper wing to exclude the range
$0<p<2$.  For the remaining rays, we compare the}
coefficients of $t^9=T^3$ in $E_1$ and $E_3$.  When $-4<p<\frac12$, this
comparison forces the coefficient $\Xi$ defined below to vanish.  A second-order
expansion then excludes the masses and signs on the locus $\Xi=0$.

{\begin{lemma}\label{lem:d8-order-six}
No Diagram~8 germ at positive masses has $0<p<2$.
\end{lemma}
\begin{proof}
Suppose that such a germ exists.  Its ray then satisfies $r>2$, so
$\nu=o(\mu)$, $T^3=o(\mu)$ and $T^5=o(\mu)$.
The exact lower row $E_4=C+\nu L_3+T^3F_3+T^5G_3=0$ gives
$C=o(\mu)$ and hence $C_5=o(\mu)$.
The centre of mass and the gauge then give
\[
A_1=X_0+o(\mu),\quad A_2=Y_0+o(\mu),\quad
X:=A_1-C_5=X_0+o(\mu),\quad Y:=A_2-C_5=Y_0+o(\mu),
\]
where $X_0=-m_2/M_{12}$ and $Y_0=m_1/M_{12}$.
Write $q=R_{12}$, $s=R_{25}$ and $q_0=R_{12}^0\ne0$.
In the upper force equations, the crossed forces contribute only
$O(\mu T^3)$.  The remaining terms are
\begin{align*}
0=E_0&=m_2q^3-m_5X+\mu A_1+O(\mu T^3),\\
0=E_2&=-m_1q^3-m_5Ys^3+\mu A_2+O(\mu T^3).
\end{align*}
The first row gives
\[
q^3-q_0^3=\frac{\mu}{M_{12}}+o(\mu),\qquad
q-q_0=\frac{\mu}{3M_{12}q_0^2}+o(\mu).
\]
In the second row the terms
$-m_1\mu/M_{12}$ and $\mu Y_0$ cancel, so
$s^3-1=o(\mu)$ and therefore $s-1=o(\mu)$.

The three exact upper distance equations give
\[
P-Q=-\frac1{m_5q^2},\qquad
(m_1+m_5)P+m_2Q=X^{-1},\qquad
m_1P+(m_2+m_5)Q=(Ys^2)^{-1}.
\]
Subtracting the last two equations yields the closure identity
\[
q^{-2}=(Ys^2)^{-1}-X^{-1}.
\]
The right-hand side differs from its limit by $o(\mu)$, while the left-hand side
differs by $-2\mu/(3M_{12}q_0^5)+o(\mu)$.  Since $M_{12}>0$ and $q_0\ne0$,
the coefficient of the latter difference is nonzero, giving a contradiction.
\end{proof}

\begin{lemma}\label{lem:d8-kappa}
Let $m\in\Rp^5$ satisfy \eqref{eq:d8-40}; let $x_0$, $\mathcal I_0$, $U_0$ and the signs
$s_{bc}$ be as in Proposition~\ref{prop:d8-first}, $\beta:=\beta_1-\beta_3$,
$b:=(f_0,f_2,f_{10},f_{13},f_{16},f_{20})$ as a function of
$(X,Y,R_{12},R_{25},P,Q)$, and $B:=\partial b/\partial\mathcal I$ at $\mathcal I_0$.
\begin{alphalist}
\item $\beta(\mathcal I_0)=0$.
\item If $m_1\ne m_2$, then $B$ is invertible.  If $m_1=m_2$ $(=4m_5)$, then $\ker B$ is spanned
by a vector $k_0$, equal to $(\frac{27}{272},\frac{27}{272},\frac9{272},-\frac9{68},-\frac{19}{17},1)$
at $(m_1,m_2,m_5)=(1,1,\frac14)$; $d\beta(\mathcal I_0)\,k_0=0$; and the Jacobian of the eight
functions $f_0,f_2,f_{10},f_{13},f_{16},f_{20}$ (evaluated at $X=A_1-C$, $Y=A_2-C$), $C$ and
$m_1A_1+m_2A_2+MC$ with respect to $U$ at $U_0$ has rank $7$.  (These three facts are proved
at $(m_1,m_2,m_5)=(1,1,\frac14)$ with $m_3,m_4$ free, which suffices by the mass scaling.)
\item With $\Xi:=s_{14}(m_1m_4)^{3/2}+s_{23}(m_2m_3)^{3/2}
+i\bigl(s_{13}(m_1m_3)^{3/2}+s_{24}(m_2m_4)^{3/2}\bigr)$,
\[
  (\tau_1-\tau_3)(x_0)=\frac{m_3m_4}{m_3+m_4}\,L^3\,(m_1m_2m_3m_4)^{-3/2}\,\Xi .
\]
\end{alphalist}
Consequently, if a Diagram~8 germ at the masses $m$ has $-4<p<\frac12$, then $\Xi=0$, that is,
$m_1=m_2=4m_5$, $m_3=m_4$, $s_{13}=-s_{24}$ and $s_{14}=-s_{23}$.
\end{lemma}}
{\begin{proof}[Proof of (a)--(c)]
On the zero set of the three upper distance faces, writing $q=R_{12}$
and $s=R_{25}$, the same subtraction as above gives
\[
q^{-2}=(Ys^2)^{-1}-X^{-1},\qquad
\beta=M_{12}q-m_5(X^{-1}-s/Y).
\]
The second identity follows by replacing $P-Q$ in the definition of
$\beta$ by $-1/(m_5q^2)$ and the other two linear combinations by their
distance values. At the face point it gives $\beta=0$, proving (a).

To prove (b), the gauge eliminates $Y$ by $Y=X+1$.
The equations $f_{10}=f_{13}=0$ determine $P,Q$ analytically, because
their linear coefficient matrix has nonzero determinant proportional to
$m_5q^2X(M_{12}+m_5)$.
After this elimination the remaining three equations are equivalent to
\[
m_2q^3-m_5X=0,\qquad
-m_1q^3-m_5Ys^3=0,\qquad
q^{-2}-(Ys^2)^{-1}+X^{-1}=0.
\]
After scaling, we may take $m_2=1$, $m_1=\alpha^2$.  The mass relation then gives
$m_5=\alpha^3/(\alpha^2+1)^2$; we set $N=\alpha^2+1$.
At the face point,
$(X,Y,q,s)=(-1/N,\alpha^2/N,-\alpha/N,1)$.
The linearizations of the first two equations give
\[
\delta q=\frac\alpha3\,\delta X,\qquad
\delta s=-\frac{N^2}{3\alpha^2}\,\delta X.
\]
Substitution in the linearization of the third gives
\[
\frac{N^2(1-\alpha^4)}{3\alpha^4}\,\delta X=0.
\]
Thus $B$ is invertible for $\alpha\ne1$.
For $\alpha=1$, its kernel is one-dimensional; choosing $\delta s=1$
gives
\[
(\delta X,\delta Y,\delta q,\delta s,\delta P,\delta Q)
=(-\tfrac34,-\tfrac34,-\tfrac14,1,\tfrac{76}{9},-\tfrac{68}{9}).
\]
Multiplication by $-9/68$ gives precisely the vector $k_0$ in (b).
Differentiating the displayed formula for $\beta$ on this tangent gives
$d\beta=2\delta q+\tfrac12\delta s=0$.
Finally $C$ fixes the translation variable and the centre-of-mass
condition is nonzero on the remaining tangent, since
$\delta(X+Y)=-3/2$. The eight functions in (b) therefore have rank $7$.

For (c), $D=0$, $B_3=-m_4/(m_3+m_4)$ and $B_4=m_3/(m_3+m_4)$ give
\[
(\tau_1-\tau_3)(x_0)
=\frac{m_3m_4}{m_3+m_4}
(-D_{13}^3+D_{23}^3+D_{14}^3-D_{24}^3).
\]
Substituting the four square roots from Proposition~\ref{prop:d8-first}
gives the following imaginary and real terms, respectively:
\[
\begin{gathered}
 iL^3s_{13}(m_2m_4)^{-3/2},\quad iL^3s_{24}(m_1m_3)^{-3/2},\\
 L^3s_{23}(m_1m_4)^{-3/2},\quad L^3s_{14}(m_2m_3)^{-3/2}.
\end{gathered}
\]
Factoring the common denominator gives
exactly the stated formula involving $\Xi$.
\end{proof}

}\begin{proof}[Proof of the last assertion]
Let $\gamma$ be such a germ.  Along $\gamma$, $\mathcal I\to\mathcal I_0$ and
$b(\mathcal I)=O(a^6)$ by Lemma~\ref{lem:d8-rows}\,(iv).  We claim
\begin{equation}\label{eq:d8-bracket}
  \beta(\mathcal I)=O(a^6+a^4\eta^4).
\end{equation}
When $m_1\ne m_2$, the invertibility in (b) and the inverse function theorem
give $\mathcal I-\mathcal I_0=O(a^6)$.  Since (a) gives the vanishing of the constant
term, this proves \eqref{eq:d8-bracket}.

For $m_1=m_2$, we use the eight functions in (b) to control the direction left
undetermined by the upper equations.  These functions are $O(a^6+\nu+T^3)$ along
$\gamma$.  For the first six, this follows from
$f_k(A_1-C,A_2-C,\dots)-f_k(\mathcal I)=O(\nu)$; the seventh is controlled by
Lemma~\ref{lem:d8-rows}\,(v).  The last is identically zero by
Lemma~\ref{lem:d8-rows}\,(ii), since $a\ne0$.  Their rank therefore gives the bounds
$U-U_0$ and $\mathcal I-\mathcal I_0$ of order $O(a^6+\nu+T^3)$.

Taylor expansion also gives
$B(\mathcal I-\mathcal I_0)=O(a^6+|\mathcal I-\mathcal I_0|^2)$, so the component of
$\mathcal I-\mathcal I_0$ transverse to $k_0$ is
$O(a^6+|\mathcal I-\mathcal I_0|^2)$.  Because $\beta(\mathcal I_0)=0$ and
$d\beta(\mathcal I_0)k_0=0$, the uncontrolled direction has no linear contribution
to the difference of the upper equations.  Consequently
$\beta(\mathcal I)=O(a^6+|\mathcal I-\mathcal I_0|^2)=O(a^6+\nu^2+T^6)$.
This is \eqref{eq:d8-bracket}, since $\nu^2=a^4\eta^4$ and $T^6=a^6\eta^6$.

Now $E_1=E_3=0$ along $\gamma$, so Lemma~\ref{lem:d8-rows}\,(iii) and \eqref{eq:d8-bracket}
give $T^3(\tau_1-\tau_3)(x)=O(a^6\eta^2+a^4\eta^6)$, that is,
$(\tau_1-\tau_3)(x)=O(a^3\eta^{-1}+a\eta^3)$.  As $-4<p<\frac12$ means $r<3$,
$\ord(a^3\eta^{-1})=(3-r)\ord a>0$; so $(\tau_1-\tau_3)(x_0)=0$, and $\Xi=0$ by (c).  Since all
$m_k>0$, the real part $s_{14}(m_1m_4)^{3/2}+s_{23}(m_2m_3)^{3/2}$ of $\Xi$ vanishes only if
$s_{14}=-s_{23}$ and $m_1m_4=m_2m_3$, and its imaginary part only if $s_{13}=-s_{24}$ and
$m_1m_3=m_2m_4$.  These give $m_3=m_4$ and $m_1=m_2$, and then \eqref{eq:d8-40} gives
$m_1=4m_5$.
\end{proof}

The order-nine coefficient vanishes precisely on the mass curve
$m=(4,4,y,y,1)$, $y>0$, up to scaling, with the four sign classes
$s_{13}=-s_{24}$, $s_{14}=-s_{23}$.  On this locus $\Xi=0$, so the preceding
comparison gives no contradiction.  We therefore examine the second-order terms.

{\begin{lemma}\label{lem:d8-cancel}
Let $m_1=m_2=4m_5$ and $m_3=m_4$, scaled to
$(1,1,m,m,\frac14)$ with $m>0$.
No Diagram~8 germ at these masses has a limit with
$s_{13}=-s_{24}$ and $s_{14}=-s_{23}$.
\end{lemma}

\begin{proof}
Suppose, for contradiction, that such a germ exists, and let $x_0$ be its limit in
the chart of Proposition~\ref{prop:d8-cover}.
We shall use the distance equations to reduce the upper-wing equations to
two scalar relations, then compare their second-order terms.

We first show that the limiting equations on the lower wing have an invertible
Jacobian.  Put $d=D_{34}$, $g=D_{35}$ and $q=D_{45}$.  At a common zero of the lower faces,
$f_1=0$ gives $D=0$. The identity
\[
m f_5+m f_7+\tfrac14 f_9=m(B_3+B_4)+\tfrac14D
\]
and $f_{21}=0$ give $B_3=-1/2$, $B_4=1/2$.
Then $f_9=0$ is $q^3=g^3$, and $f_5=0$ is $8md^3+g^3=4$.
The equations $f_{18}=f_{19}=0$ read
\[
(m+\tfrac14)u+mv=-2g^{-2},\qquad
mu+(m+\tfrac14)v=2q^{-2}.
\]
Their coefficient determinant is $\tfrac14(2m+\tfrac14)>0$.
Subtracting them and using $f_{17}=0$ gives
$2d^2(g^{-2}+q^{-2})=1$.
Writing $q=\omega g$ and $d=\lambda g$, we have thus obtained
\[
\omega^3=1,\qquad 2\lambda^2(1+\omega^{-2})=1,
\qquad g^3(8m\lambda^3+1)=4.
\]
The representation in the chart ensures that all inverse distances are nonzero.
The preceding elimination gives an invertible Jacobian for the eight lower faces:
$f_1$, the displayed weighted sum, and $f_{21}$ determine $D,B_3,B_4$;
the next three successively determine
$\omega,\lambda,g$, with derivatives
\[
3\omega^2,\qquad4\lambda(1+\omega^{-2}),\qquad
3g^2(8m\lambda^3+1)=12/g,
\]
all nonzero; and the last linear system determines $u,v$.

We now turn to the upper wing and write its deviations from the limit as
\[
X=A_1-C_5=-\tfrac12+x,\qquad
Y=A_2-C_5=\tfrac12+x,\qquad
R_{12}=-\tfrac12+v_*,\qquad R_{25}=1+h_*.
\]
The equality $Y-X=1$ follows from the gauge.  The symbols $v_*,h_*$
denote upper-wing deviations and are distinct from the chart coordinates $v,h$.
Write $r_*=R_{12}$ and $s_*=R_{25}$.
The three upper distance equations hold along the whole germ:
\[
P-Q=-4r_*^{-2},\qquad
\tfrac54 P+Q=X^{-1},\qquad
P+\tfrac54Q=(Ys_*^2)^{-1}.
\]
Consequently
\begin{equation}\label{eq:d8-upper-closure-analytic}
r_*^{-2}=(Ys_*^2)^{-1}-X^{-1}.
\end{equation}
They determine $P,Q$ analytically once this closure equation holds.
Also, direct grouping of the upper forces gives
\begin{align*}
E_0&=r_*^3-\tfrac14X+\mu A_1+O(\mu T^3),\\
E_2&=-r_*^3-\tfrac14Ys_*^3+\mu A_2+O(\mu T^3),
\end{align*}
because every crossed force in an upper $z$-row has factor $\mu T^3$.
Here and below all estimates are along the germ, whose chart coordinates
are bounded.

Expanding \eqref{eq:d8-upper-closure-analytic} and $E_2=0$ to second order gives
\begin{align*}
0={}&16v_*+4h_*+48v_*^2-16x^2-8xh_*-6h_*^2
 +O\bigl((|x|+|v_*|+|h_*|)^3\bigr),\\
\tfrac12\mu={}&\tfrac14x+\tfrac34v_*+\tfrac38h_*
 -\tfrac32v_*^2+\tfrac34xh_*+\tfrac38h_*^2
 +O\bigl((|x|+|v_*|+|h_*|)^3\bigr)+o(\mu).
\end{align*}
The derivative of these two equations with respect to $(x,v_*)$ is invertible
at zero. Solving the linear and quadratic terms therefore gives
\begin{equation}\label{eq:d8-upper-expansion-analytic}
x=-\tfrac34h_*+2\mu+o(\mu+h_*^2),\qquad
v_*=-\tfrac14h_*+\tfrac38h_*^2+o(\mu+h_*^2).
\end{equation}
More precisely, the analytic implicit solution before substituting the
remaining bounded coordinates has remainders
$O(h_*^3+\mu h_*+\mu^2)+o(\mu)$; thus the displayed estimates hold also
when one of the deviations vanishes identically.
Substituting in $E_0=0$ now gives
\begin{equation}\label{eq:d8-row0-analytic}
-\mu+\tfrac3{16}h_*^2=o(\mu+h_*^2).
\end{equation}

The same exact distance equations simplify the difference
$\beta=\beta_1-\beta_3$ to
\[
\beta=2r_*-\tfrac14(X^{-1}-s_*/Y).
\]
Its quadratic expansion is
$2v_*+\tfrac12h_*+4x^2-xh_*+O((|x|+|v_*|+|h_*|)^3)$.
By \eqref{eq:d8-upper-expansion-analytic},
\begin{equation}\label{eq:d8-beta-analytic}
\beta=\tfrac{15}{4}h_*^2+o(\mu+h_*^2).
\end{equation}
Finally, direct substitution in the two upper $w$-equations gives
\begin{equation}\label{eq:d8-w-difference-analytic}
E_1-E_3=\eta^2\beta+T^3(\tau_1-\tau_3)
 +\tfrac14\mu\eta^2(P-Q)+O(\mu\eta^2T^3).
\end{equation}
In particular, $P_0-Q_0=-16$.

To compare the remaining terms, we need to control the lower and crossed
coordinates.  Their deviations from the limits are
\begin{equation}\label{eq:d8-lower-bound-analytic}
O(\nu+\mu+|h_*|).
\end{equation}
Indeed, in the eight lower rows, the terms beyond the limiting equations carry
factors $T^3$, $\mu\eta^2$, $\mu\eta^2T^3$, or $\eta^2\beta_j$;
the last kind occurs in the upper or central $w$-row. Its coefficient
vanishes at the upper face point, where $(P,Q,r_*,s_*)=(-8,8,-1/2,1)$.
For the central row this coefficient is $-4(\beta_1+\beta_3)$.
Thus these errors are $O(\nu+\mu+|h_*|)$, since $T^3=o(\nu)$ and
\eqref{eq:d8-upper-expansion-analytic} bounds the upper deviations by
$O(\mu+|h_*|)$. The crossed distance rows have only the additional factors
$\nu$, $\mu\eta^2$ and $\nu\mu\eta^2$; each crossed face has derivative
$2/D_{bc}\ne0$ in its own inverse distance. The implicit function theorem
and the lower nondegeneracy proved above give
\eqref{eq:d8-lower-bound-analytic}.

Put $M=2m+\tfrac14$ and
\[
\Psi=L_3+\frac{2m(u+v)}{2+M}.
\]
The centre-of-mass identity and the exact row
$E_4=C+\nu L_3+T^3F_3+T^5G_3=0$ give
\[
X+Y=(2+M)(\nu\Psi+T^3F_3+T^5G_3).
\]
Since $X+Y=2x=-\tfrac32h_*+O(\mu+h_*^2)$, first
$h_*=O(\nu+T^3+\mu)$ and then, by
\eqref{eq:d8-lower-bound-analytic},
\[
h_*=h_\nu\nu+O(T^3+\nu^2+\mu),\qquad
h_\nu=-\tfrac23(2+M)\Psi(x_0).
\]
If $h_\nu=0$, then
$h_*^2=O(T^6+\nu^4+\mu^2)=o(\mu)$, contradicting
\eqref{eq:d8-row0-analytic}.
If $h_\nu\ne0$, then
$h_*^2=h_\nu^2\nu^2+o(\mu+\nu^2)$.
Equation~\eqref{eq:d8-row0-analytic} forces
$\operatorname{ord}\eta=\tfrac12\operatorname{ord}a$ and gives
\begin{equation}\label{eq:d8-first-balance-analytic}
-a_0^2+\tfrac3{16}h_\nu^2\eta_0^4=0,
\end{equation}
where $a_0,\eta_0\ne0$ are the leading coefficients.

The crossed signs in the statement give $(\tau_1-\tau_3)(x_0)=0$.
By \eqref{eq:d8-lower-bound-analytic},
$T^3(\tau_1-\tau_3)=O(T^3(\nu+\mu+|h_*|))$.
On the ray just obtained, this is higher order than $a^6\eta^2$.
Equations~\eqref{eq:d8-beta-analytic} and
\eqref{eq:d8-w-difference-analytic} therefore give the second balance
\[
-4a_0^2+\tfrac{15}{4}h_\nu^2\eta_0^4=0.
\]
Together with \eqref{eq:d8-first-balance-analytic}, this is a homogeneous
linear system with determinant $-15/4+3/4=-3\ne0$.  It forces $a_0=0$,
contrary to the choice of a nonzero leading coefficient.  The germ is therefore excluded.
\end{proof}

}\begin{theorem}\label{thm:d8}
For $m\in\Rp^5$, no curve in $A_m$ carries a Diagram~8 germ, real or complex.
In particular, Theorem~\ref{thm:exclusion}\,(a) holds for $r=8$.
\end{theorem}

\begin{proof}
By the symmetries at the beginning of this section we may use the numbering of
\cite[\S8.8]{AK}.  Let $\gamma$ be a Diagram~8 germ at $m\in\Rp^5$.  By
Proposition~\ref{prop:d8-cover} it is a ray germ with ray $r\in(0,\infty)$, and by
Proposition~\ref{prop:d8-first} the masses satisfy \eqref{eq:d8-40} and the limit has signs
$s_{bc}$.  The case $r>2$, equivalently $0<p<2$, is excluded by
Lemma~\ref{lem:d8-order-six}.  If $r\le2$, equivalently $-4<p\le0$,
Lemma~\ref{lem:d8-kappa} gives $m_1=m_2=4m_5$, $m_3=m_4$, $s_{13}=-s_{24}$
and $s_{14}=-s_{23}$.  These masses and signs are excluded by
Lemma~\ref{lem:d8-cancel}.  Both ranges therefore lead to a contradiction, proving that
$\gamma$ does not exist and, in particular, excluding germs with an unbounded four-product.
\end{proof}

\appendix

\section{Computational statements used in the proof}\label{app:verification}

The computer-assisted statements used in the proofs of Diagrams~2, 3 and~6 are listed in Table~\ref{tab:verification}. Each entry identifies the corresponding assertion in the main text and its associated program modules.

\begin{table}[ht]
\centering\footnotesize
\begin{tabular}{p{4.8cm}p{8.2cm}}
\hline
Computational statement & Associated program modules (prefix \code{ak\_} omitted)\\
\hline
Lemma~\ref{lem:d2-fibre}, Definition~\ref{def:d2-mass}
 & \code{blowup}, \code{resultants}, \code{p\_source\_completeness},
 \code{p\_general}, \code{diagram2\_symmetry}\\
Lemma~\ref{lem:d2-empty}
 & \code{l4\_same\_sign}, \code{resultants}, \code{exceptional}, \code{l5\_obstruction}\\
Lemma~\ref{lem:d2-L5}: jets, minors and label orbits
 & \code{l5\_obstruction}\\
Lemmas~\ref{lem:d2-rays}, \ref{lem:d2-Lray}
 & \code{exceptional\_branch\_positive\_feasibility}, \code{l1\_complex\_fibre},
 \code{l\_block\_symmetric\_germs}\\
Lemma~\ref{lem:d2-K}
 & \code{k1\_curve}, \code{k1\_complex\_fibre}, \code{p4\_curve},
 \code{p5\_completion}, \code{resultants}, \code{diagram2\_symmetry}\\
Lemmas~\ref{lem:d2-P}, \ref{lem:d2-fold}
 & \code{p\_sheet\_second\_order}, \code{p\_sheet\_special\_ends},
 \code{p\_sheet\_e5\_corner}, \code{p\_sheet\_euler\_closed\_form}, \code{p\_sheet\_p5\_exact}\\
\hline
Lemmas~\ref{lem:d3-order3}, \ref{lem:d3-order4}
 & \code{diagram3\_blowup}\\
Lemma~\ref{lem:d3-order5}: jet and product identities
 & \code{diagram3\_fifth\_order}\\
Lemma~\ref{lem:d3-order6}: $144$ excluded labels
 & \code{diagram3\_equal\_shape}\\
\hline
Diagram~6 orbit coverage and Lemma~\ref{d6-lem-corank}
 & \code{diagram6\_blowup}, \code{diagram6\_primitive\_cubic},
 \code{diagram6\_primitive\_quadratic}, \code{diagram6\_primitive\_opposite}\\
\hline
\end{tabular}
\caption{Computational statements used in the proof.}
\label{tab:verification}
\end{table}


\section*{Acknowledgments}
%\addcontentsline{toc}{section}{Acknowledgments}
The initial version of this paper was submitted to arXiv on October 5, 2026 (submission ID: 8179056) by the current second author on behalf of the first author; however, this submission did not comply with arXiv proxy submission policies. Subsequently, the manuscript was resubmitted by the first author on October 6, 2026 (submission ID: 8186741), but has remained on hold. Shortly thereafter, \cite{FiveT} (arXiv:2610.13100) appeared reporting identical results. The present version was submitted by the second author.

This manuscript was prepared with AI assistance and preliminary automated verification. Rigorous manual verification by the second author is currently underway and will be incorporated into a forthcoming version, alongside a complete formal proof in Lean. The authors thank T. Tomar for pointing out a non-essential gap.


\begin{thebibliography}{99}

\bibitem{AK}
A.~Albouy and V.~Kaloshin, \emph{Finiteness of central configurations of five bodies in the
plane}, Ann. of Math. (2) \textbf{176} (2012), 535--588.

\bibitem{BCR}
J.~Bochnak, M.~Coste and M.-F.~Roy, \emph{Real Algebraic Geometry}, Ergeb. Math. Grenzgeb. (3)
\textbf{36}, Springer, 1998.

\bibitem{CC1}
K.-M.~Chang and K.-C.~Chen, \emph{Toward finiteness of central configurations for the planar
six-body problem by symbolic computations. (I) Determine diagrams and orders}, J. Symbolic
Comput. \textbf{123} (2024), 102277; arXiv:2303.02853.

\bibitem{CC2}
K.-M.~Chang and K.-C.~Chen, \emph{Toward finiteness of central configurations for the planar
six-body problem by symbolic computations. (II) Determine mass relations}, SIAM J. Appl. Dyn.
Syst. \textbf{24} (2025), 2369--2404.

\bibitem{DH}
Y.~Deng and M.~Hampton, \emph{Equilateral chains and cyclic central configurations of the
planar $5$-body problem}, arXiv:2112.06755, 2021.

\bibitem{HJ}
M.~Hampton and A.~N.~Jensen, \emph{Finiteness of spatial central configurations in the
five-body problem}, Celestial Mech. Dynam. Astronom. \textbf{109} (2011), 321--332.

\bibitem{HM}
M.~Hampton and R.~Moeckel, \emph{Finiteness of relative equilibria of the four-body problem},
Invent. Math. \textbf{163} (2006), 289--312.

\bibitem{JL}
A.~N.~Jensen and A.~Leykin, \emph{Smale's 6th problem for generic masses}, J. Exp. Math.
\textbf{1} (2025), 306--314; arXiv:2301.02305.

\bibitem{LS}
T.-L.~Lee and M.~Santoprete, \emph{Central configurations of the five-body problem with equal
masses}, Celestial Mech. Dynam. Astronom. \textbf{104} (2009), 369--381,
doi:10.1007/s10569-009-9219-0.

\bibitem{MZ19}
M.~Moczurad and P.~Zgliczy\'nski, \emph{Central configurations in planar $n$-body problem with
equal masses for $n=5,6,7$}, Celestial Mech. Dynam. Astronom. \textbf{131} (2019),
doi:10.1007/s10569-019-9920-6; arXiv:1812.07279.

\bibitem{MZ26}
M.~Moczurad and P.~Zgliczy\'nski, \emph{Central configurations with unequal masses: finiteness
in several exceptional cases of five bodies}, arXiv:2601.01165, 2026.
\bibitem{Moulton}
F.~R.~Moulton, \emph{The straight line solutions of the problem of $n$ bodies}, Ann. of Math. (2)
\textbf{12} (1910), 1--17.

\bibitem{Roberts}
G.~E.~Roberts, \emph{A continuum of relative equilibria in the five-body problem}, Phys. D
\textbf{127} (1999), 141--145.

\bibitem{Smale}
S.~Smale, \emph{Mathematical problems for the next century}, Math. Intelligencer \textbf{20}
(1998), no.~2, 7--15.

\bibitem{FiveT}
T. Tomar, \emph{Finiteness of planar central configurations of five bodies for all positive masses}, arXiv:2610.13100v1, 2026.


\bibitem{Wintner}
A.~Wintner, \emph{The Analytical Foundations of Celestial Mechanics}, Princeton University
Press, 1941.

\bibitem{YZ}
X.~Yu and S.~Zhu, \emph{On the finiteness of four-body central configurations},
arXiv:2103.08906, 2021.

\end{thebibliography}

\end{document}
